% Leçon 16 — Grammaire : Les compl.de durée ; redoublement des adjectifs — Chinois 1ère A — Cours signé OnBuch+
% Programme officiel MINESEC. Compiler avec : tectonic ZHA16-grammaire-les-compl-de-duree-redoublement-des-adje.tex
\newcommand{\DOCMATIERE}{Chinois}
\newcommand{\DOCNIVEAU}{1ère A}
\newcommand{\DOCMODULE}{Module 2 : Environnement et santé}
\newcommand{\DOCLECON}{ZHA16}
\newcommand{\DOCTITRE}{Leçon 16 — Grammaire : Les compl.de durée ; redoublement des adjectifs}
% OnBuch+ — préambule « cours signés » ENRICHI (compilé avec tectonic / XeLaTeX)
% Système visuel « Pop » d'OnBuch+ : fond crème (jamais blanc), bordures encre
% épaisses, ombres « dures » décalées, titres Archivo Black en capitales.
% Version MATHÉMATIQUES : graphes (pgfplots/tikz), boîtes pédagogiques
% (définition, propriété, méthode, attention, le savais-tu, exercice résolu,
% exercice, TP, bilan), page de garde et signature OnBuch+ ancrée en bas.
%
% Macros attendues dans main.tex AVANT \input{preamble} :
%   \DOCMATIERE  \DOCNIVEAU  \DOCMODULE  \DOCLECON (numéro)  \DOCTITRE
\documentclass[11pt]{article}

\usepackage{fontspec}
\usepackage[french]{babel} % français = langue principale ; le chinois via \zh{..} / textebox (xeCJK)
\usepackage{xeCJK}
\usepackage{pifont} % \ding{51} \ding{55} utilisés par les modèles
\usepackage[a4paper,margin=2cm,top=3.4cm,bottom=2.8cm,headheight=1.4cm,footskip=1.3cm]{geometry}
\usepackage{amsmath,amssymb,amsfonts}
\usepackage{mathtools} % \xmapsto, \coloneqq... souvent utilisés par les modèles
\usepackage{cancel} % simplifications barrées (bilans de forces)
% \abs{z} (module, valeur absolue) et \norm{u} (norme), utilisés par les modèles
\providecommand{\abs}{}\let\abs\relax\DeclarePairedDelimiter{\abs}{\lvert}{\rvert}
\providecommand{\norm}{}\let\norm\relax\DeclarePairedDelimiter{\norm}{\lVert}{\rVert}
\usepackage[table]{xcolor} % [table] : fournit aussi \rowcolors (alternance de lignes)
\usepackage{pagecolor}
\usepackage{tikz}
\usetikzlibrary{arrows.meta,positioning,calc,shapes.geometric,shapes.misc,shapes.arrows,decorations.pathmorphing,decorations.pathreplacing,decorations.markings,patterns,shadows,mindmap,trees,fit,backgrounds,matrix,intersections,angles,quotes}
\usepackage{pgfplots}
\usepackage[version=4]{mhchem} % équations nucléaires \ce{^{235}_{92}U}
\usepackage[european,siunitx]{circuitikz} % schémas électriques
\pgfplotsset{compat=1.18}
\usepgfplotslibrary{fillbetween,groupplots} % aires entre courbes (calcul intégral)
\usepackage{enumitem}
\usepackage{fancyhdr}
% [most] charge toutes les bibliothèques tcolorbox (listings, minted, raster,
% documentation...) et gonfle inutilement la mémoire TeX sur un document long
% (~300 boîtes) : on ne charge que ce qui est réellement utilisé.
\usepackage[skins,breakable]{tcolorbox}
\usepackage{graphicx}
\usepackage{colortbl}
\usepackage{booktabs}
\usepackage{tabularx}
\usepackage{diagbox} % case d'en-tête diagonale des tableaux à double entrée
% Colonnes extensibles centrée / gauche / droite (C, L, R), souvent utilisées par les modèles
\newcolumntype{C}{>{\centering\arraybackslash}X}
\newcolumntype{L}{>{\raggedright\arraybackslash}X}
\newcolumntype{R}{>{\raggedleft\arraybackslash}X}
\usepackage{array}
\usepackage{multicol}
\usepackage{multirow}
\usepackage{siunitx}
% parse-numbers=false : accepte aussi \SI{2,5 \times 10^{-3}}{...}
\sisetup{locale=FR,output-decimal-marker={,},group-minimum-digits=4,parse-numbers=false}
\DeclareSIUnit\ohm{\ensuremath{\symup{\Omega}}} % U+2126 absent de la police
\DeclareSIUnit\division{div} % réglages d'oscilloscope (ms/div, V/div)
\DeclareSIUnit\tr{tr} % tours (vitesse de rotation, ex. \SI{1500}{\tr\per\minute})
\DeclareSIUnit\molar{M} % concentration molaire (ex. \SI{3}{\molar}, \SI{1}{\micro\molar})
\DeclareSIUnit\an{an} % année (le modèle confond parfois avec \year, un compteur TeX) — ex. \SI{5}{\kilo\meter\per\an}
\DeclareSIUnit\FCFA{FCFA} % franc CFA (ex. \SI{500}{\FCFA\per\kilo\gram})
\DeclareSIUnit\degreeBrix{\textdegree Brix} % teneur en sucre d'un sirop/jus (réfractométrie)
\DeclareSIUnit\month{mois} % le modèle utilise parfois \month, un compteur TeX, comme unité de durée
\usepackage{qrcode}
% \degree : commande fréquemment utilisée par les modèles pour un angle ou une
% température en dehors de \SI{}{} (non fournie par les paquets chargés).
\providecommand{\degree}{^{\circ}}
% \qed (fin de démonstration), utilisé par les modèles sans amsthm
\providecommand{\qed}{\hfill$\square$}
% \barre : double barre des valeurs interdites dans un tableau de variations (tkz-tab)
\providecommand{\barre}{\ensuremath{\|}}
% environnement proof (amsthm non chargé) utilisé par les modèles
\ifdefined\proof\else\newenvironment{proof}[1][Démonstration]{\par\noindent\textit{#1.}\ }{\qed\par}\fi
\usepackage{lastpage}
\usepackage{hyperref}
\hypersetup{hidelinks}

% ============================================================
% Palette « Pop » OnBuch+ — mêmes hex que la classe P (plus_ui.dart)
% ============================================================
\definecolor{poporange}{HTML}{F59321}
\definecolor{popdark}{HTML}{E07A0C}
\definecolor{popcream}{HTML}{FFF6E8}
\definecolor{popink}{HTML}{1C1714}
\definecolor{poppurple}{HTML}{7A5AE0}
\definecolor{popgreen}{HTML}{2E9E6A}
\definecolor{poppink}{HTML}{E0457B}
\definecolor{popblue}{HTML}{2F7DE1}
\definecolor{popgold}{HTML}{C98A12}
\definecolor{popmuted}{HTML}{8A7F76}
\definecolor{popline}{HTML}{EADFD2}
\definecolor{popgray}{HTML}{8A7F76}
\colorlet{popgrey}{popgray}
% Alias tolérants (le modèle nomme parfois les couleurs différemment)
\colorlet{popwhite}{white}
\colorlet{popblack}{popink}
\colorlet{popred}{poppink}
\colorlet{popyellow}{popgold}
% Teintes claires (fonds de tableaux/encadrés) — le modèle nomme parfois
% les couleurs avec un suffixe L (ex. popblueL) sans les avoir définies.
\colorlet{poporangeL}{poporange!20}
\colorlet{popdarkL}{popdark!20}
\colorlet{poppurpleL}{poppurple!20}
\colorlet{popgreenL}{popgreen!20}
\colorlet{poppinkL}{poppink!20}
\colorlet{popblueL}{popblue!20}
\colorlet{popgoldL}{popgold!20}
\colorlet{popredL}{popred!20}
\colorlet{popgrayL}{popgray!20}
% Teintes claires pour fonds de tableaux / zones
\colorlet{popblueL}{popblue!10!white}
\colorlet{poporangeL}{poporange!14!white}
\colorlet{popgreenL}{popgreen!12!white}
\colorlet{poppurpleL}{poppurple!12!white}
\colorlet{poppinkL}{poppink!12!white}
\colorlet{popgoldL}{popgold!14!white}

\pagecolor{popcream}

% ============================================================
% Typographie
% ============================================================
\defaultfontfeatures{Ligatures=TeX}
\setmainfont{PlusJakartaSans-Medium.ttf}[Path=../../../fonts/,BoldFont=PlusJakartaSans-Bold.ttf]
% Symboles, API et emojis absents de la police principale : repli sur DejaVu Sans (caractères actifs)
\newfontface\symfont{DejaVuSans.ttf}[Path=../../../fonts/]
\makeatletter
\newcommand\symactive[1]{\catcode#1=13 \begingroup\lccode`\~=#1 \lowercase{\endgroup\def~}{{\symfont\char#1}}}
\newcommand\symblank[1]{\catcode#1=13 \begingroup\lccode`\~=#1 \lowercase{\endgroup\def~}{}}
\newcommand\symhyph[1]{\catcode#1=13 \begingroup\lccode`\~=#1 \lowercase{\endgroup\def~}{\mbox{-}}}
\newcount\symc
\newcommand\symrange[2]{\symc=#1 \@whilenum\symc<#2 \do{\expandafter\symactive\expandafter{\the\symc}\advance\symc by 1 }}
\newcommand\symblankrange[2]{\symc=#1 \@whilenum\symc<#2 \do{\expandafter\symblank\expandafter{\the\symc}\advance\symc by 1 }}
\makeatother
\symrange{"0250}{"02FF}\symactive{"2713}\symactive{"2714}\symactive{"2717}\symactive{"2718}\symactive{"2610}\symactive{"2611}\symactive{"2612}
\symrange{"2600}{"27BF}
\catcode"2705=13 \begingroup\lccode`\~="2705 \lowercase{\endgroup\def~}{{\symfont\char"2713}}\symblank{"26BD}\symblank{"26A1}
\symrange{"2460}{"257F}\symactive{"223C}\symactive{"2248}\symactive{"2260}
\symactive{"01F8}\symactive{"01F9}\symactive{"1E3E}\symactive{"1E3F}
\symhyph{"2011}\symhyph{"2E1A}\symblankrange{"1F300}{"1FAFF}\symblank{"FE0F}
% Caractères chinois : WenQuanYi Zen Hei (sous-ensemble GB2312 + pinyin), gras/italique simulés
\setCJKmainfont{WenQuanYiZenHei-subset.ttf}[Path=../../../fonts/,AutoFakeBold=3,AutoFakeSlant=0.2]
\xeCJKsetup{CJKspace=false}
% Pinyin : voyelles à caron (ǎ ǐ ǒ ǔ ǖ ǘ ǚ ǜ) absentes de la police principale -> caractères actifs
\newfontface\pyfont{WenQuanYiZenHei-subset.ttf}[Path=../../../fonts/]
\newcommand\pyactive[1]{\catcode#1=13 \begingroup\lccode`\~=#1 \lowercase{\endgroup\def~}{{\pyfont\char#1}}}
\pyactive{"01CD}\pyactive{"01CE}\pyactive{"01CF}\pyactive{"01D0}\pyactive{"01D1}\pyactive{"01D2}\pyactive{"01D3}\pyactive{"01D4}\pyactive{"01D5}\pyactive{"01D6}\pyactive{"01D7}\pyactive{"01D8}\pyactive{"01D9}\pyactive{"01DA}\pyactive{"01DB}\pyactive{"01DC}
% Maths en sans-serif (Fira Math) avec chiffres et lettres droites de Plus
% Jakarta, pour que les formules se fondent dans le texte.
\usepackage[math-style=ISO,bold-style=ISO]{unicode-math}
\setmathfont{FiraMath-Regular.otf}[Path=../../../fonts/]
\setmathfont{PlusJakartaSans-Medium.ttf}[Path=../../../fonts/,range={up/num,up/{Latin,latin},\mathup}]
\usepackage{icomma} % virgule décimale sans espace en mode math
% Caractères mathématiques Unicode absents des polices de texte (Plus Jakarta,
% Archivo Black) : rendus par leur équivalent mathématique, en texte comme en
% formule — sinon ils s'affichent en carré vide (ex. « Arithmétique dans ℤ »).
\usepackage{newunicodechar}
\newunicodechar{ℝ}{\ensuremath{\mathbb{R}}}
\newunicodechar{ℤ}{\ensuremath{\mathbb{Z}}}
\newunicodechar{ℂ}{\ensuremath{\mathbb{C}}}
\newunicodechar{ℕ}{\ensuremath{\mathbb{N}}}
\newunicodechar{ℚ}{\ensuremath{\mathbb{Q}}}
\newunicodechar{𝔻}{\ensuremath{\mathbb{D}}}
\newunicodechar{α}{\ensuremath{\alpha}}
\newunicodechar{β}{\ensuremath{\beta}}
\newunicodechar{γ}{\ensuremath{\gamma}}
\newunicodechar{δ}{\ensuremath{\delta}}
\newunicodechar{ε}{\ensuremath{\varepsilon}}
\newunicodechar{θ}{\ensuremath{\theta}}
\newunicodechar{λ}{\ensuremath{\lambda}}
\newunicodechar{μ}{\ensuremath{\mu}}
\newunicodechar{σ}{\ensuremath{\sigma}}
\newunicodechar{τ}{\ensuremath{\tau}}
\newunicodechar{φ}{\ensuremath{\varphi}}
\newunicodechar{ψ}{\ensuremath{\psi}}
\newunicodechar{ω}{\ensuremath{\omega}}
\newunicodechar{Σ}{\ensuremath{\Sigma}}
% majuscules produites par \MakeUppercase dans les titres (α-aminés -> Α)
\newunicodechar{Α}{\ensuremath{\alpha}}
\newunicodechar{Β}{\ensuremath{\beta}}
\newunicodechar{Γ}{\ensuremath{\Gamma}}
\newunicodechar{Δ}{\ensuremath{\Delta}}
\newunicodechar{Ε}{\ensuremath{\varepsilon}}
\newunicodechar{Θ}{\ensuremath{\Theta}}
\newunicodechar{Λ}{\ensuremath{\Lambda}}
\newunicodechar{Μ}{\ensuremath{\mu}}
\newunicodechar{Τ}{\ensuremath{\tau}}
\newunicodechar{Φ}{\ensuremath{\Phi}}
\newunicodechar{Ψ}{\ensuremath{\Psi}}
\newunicodechar{Ω}{\ensuremath{\Omega}}
\newunicodechar{↦}{\ensuremath{\mapsto}}
\newunicodechar{∈}{\ensuremath{\in}}
\newunicodechar{∉}{\ensuremath{\notin}}
\newunicodechar{∧}{\ensuremath{\wedge}}
\newunicodechar{⇔}{\ensuremath{\Leftrightarrow}}
\newunicodechar{⟺}{\ensuremath{\Longleftrightarrow}}
\newunicodechar{⇒}{\ensuremath{\Rightarrow}}
\newunicodechar{⟹}{\ensuremath{\Longrightarrow}}
\newunicodechar{∘}{\ensuremath{\circ}}
\newunicodechar{∪}{\ensuremath{\cup}}
\newunicodechar{∩}{\ensuremath{\cap}}
\newunicodechar{≡}{\ensuremath{\equiv}}
\newunicodechar{⊂}{\ensuremath{\subset}}
\newunicodechar{⊥}{\ensuremath{\perp}}
\newunicodechar{∃}{\ensuremath{\exists}}
\newunicodechar{∀}{\ensuremath{\forall}}
\newunicodechar{∛}{\ensuremath{\sqrt[3]{\,}}}
\newunicodechar{∜}{\ensuremath{\sqrt[4]{\,}}}
\newunicodechar{⃗}{\ensuremath{{}^{\to}}}
\newunicodechar{ⁿ}{\ifmmode^{n}\else\textsuperscript{\ensuremath{n}}\fi}
\newunicodechar{ˣ}{\ifmmode^{x}\else\textsuperscript{\ensuremath{x}}\fi}
\newunicodechar{ᵇ}{\ifmmode^{b}\else\textsuperscript{\ensuremath{b}}\fi}
\newunicodechar{ʳ}{\ifmmode^{r}\else\textsuperscript{\ensuremath{r}}\fi}
\newunicodechar{ᵘ}{\ifmmode^{u}\else\textsuperscript{\ensuremath{u}}\fi}
\newunicodechar{ʸ}{\ifmmode^{y}\else\textsuperscript{\ensuremath{y}}\fi}
\newunicodechar{ᵃ}{\ifmmode^{a}\else\textsuperscript{\ensuremath{a}}\fi}
\newunicodechar{ᵉ}{\ifmmode^{e}\else\textsuperscript{\ensuremath{e}}\fi}
\newunicodechar{ᵏ}{\ifmmode^{k}\else\textsuperscript{\ensuremath{k}}\fi}
\newunicodechar{ᵖ}{\ifmmode^{p}\else\textsuperscript{\ensuremath{p}}\fi}
\newunicodechar{ᴺ}{\ifmmode^{N}\else\textsuperscript{\ensuremath{N}}\fi}
\newunicodechar{ᶜ}{\ifmmode^{c}\else\textsuperscript{\ensuremath{c}}\fi}
\newunicodechar{ʰ}{\ifmmode^{h}\else\textsuperscript{\ensuremath{h}}\fi}
\newunicodechar{ᵠ}{\ifmmode^{q}\else\textsuperscript{\ensuremath{q}}\fi}
\newunicodechar{ₙ}{\ifmmode_{n}\else\textsubscript{\ensuremath{n}}\fi}
\newunicodechar{ₐ}{\ifmmode_{a}\else\textsubscript{\ensuremath{a}}\fi}
\newunicodechar{ᵢ}{\ifmmode_{i}\else\textsubscript{\ensuremath{i}}\fi}
\newunicodechar{ᵣ}{\ifmmode_{r}\else\textsubscript{\ensuremath{r}}\fi}
\newunicodechar{ₖ}{\ifmmode_{k}\else\textsubscript{\ensuremath{k}}\fi}
\newunicodechar{ᵦ}{\ifmmode_{\beta}\else\textsubscript{\ensuremath{\beta}}\fi}
\newunicodechar{ₓ}{\ifmmode_{x}\else\textsubscript{\ensuremath{x}}\fi}
\newfontfamily\popdisplay{ArchivoBlack-Regular.ttf}[Path=../../../fonts/]
\newfontfamily\poplogo{SpaceGrotesk-Bold.ttf}[Path=../../../fonts/]
\newfontfamily\popsemi{PlusJakartaSans-SemiBold.ttf}[Path=../../../fonts/]
\newfontfamily\popxbold{PlusJakartaSans-ExtraBold.ttf}[Path=../../../fonts/]
\newfontfamily\popxboldls{PlusJakartaSans-ExtraBold.ttf}[Path=../../../fonts/,LetterSpace=3.0]

\setlist[enumerate]{leftmargin=1.7em,itemsep=5pt,topsep=4pt}
\setlist[itemize]{leftmargin=1.7em,itemsep=4pt,topsep=4pt,label={\color{poporange}\textbullet}}
\setlength{\parindent}{0pt}
\setlength{\parskip}{5pt}
\renewcommand{\arraystretch}{1.25}

% pgfplots : style Pop par défaut
\pgfplotsset{
  every axis/.append style={axis line style={popink,line width=1pt},tick style={popink},
    grid style={popline,line width=0.5pt},label style={font=\small},tick label style={font=\footnotesize}},
  popcurve/.style={poporange,line width=1.8pt,no markers,smooth},
}
% Même style disponible dans un \draw TikZ simple (hors pgfplots)
\tikzset{popcurve/.style={poporange,line width=1.8pt}}
\tikzset{popgrid/.style={popline,very thin}}
\colorlet{popcurve}{poporange} % le nom du style est parfois utilisé comme couleur

% ============================================================
% Pill / squiggle
% ============================================================
\newcommand{\poppill}[3]{%
  \tikz[baseline=-0.55ex]\node[draw=popink,line width=1.2pt,rounded corners=2.6pt,%
    fill=#2,inner xsep=6.5pt,inner ysep=3pt]{%
    \popxboldls\fontsize{7}{7}\selectfont\color{#3}\MakeUppercase{#1}};%
}
\newcommand{\popsquiggle}[1]{%
  \tikz[baseline=-0.4ex]{
    \draw[#1,line width=2.6pt,line cap=round]
      plot [smooth] coordinates {(0,0.09) (0.34,0.01) (0.68,0.21) (1.02,0.01) (1.36,0.10)};
  }%
}
% Mot-clé surligné (marqueur orange sous le texte)
\newcommand{\cle}[1]{{\popxbold\color{popdark}#1}}

% ============================================================
% En-tête / pied
% ============================================================
\pagestyle{fancy}
\fancyhf{}
\renewcommand{\headrulewidth}{0pt}
\renewcommand{\footrulewidth}{0pt}
\lhead{\raisebox{-0.15ex}{\poplogo\fontsize{13}{13}\selectfont\color{popink}OnBuch\color{poporange}+}}
\rhead{\poppill{\DOCMATIERE\ \textbullet\ \DOCNIVEAU\ \textbullet\ Leçon \DOCLECON}{white}{popink}}
\lfoot{\popxbold\fontsize{7.6}{7.6}\selectfont\color{popmuted}\MakeUppercase{Cours signé OnBuch+}\ \textbullet\ {\popsemi\DOCTITRE}}
\rfoot{\tikz[baseline=-0.6ex]\node[rounded corners=8.5pt,fill=popink,minimum height=17pt,minimum width=17pt,inner xsep=5pt,inner sep=0]{\popxbold\fontsize{8}{8}\selectfont\color{popcream}\thepage\,/\,\pageref*{LastPage}};}

% ============================================================
% Titres
% ============================================================
\newcommand{\chaptitre}[2]{%
  \begin{tcolorbox}[enhanced,colback=poporange,colframe=popink,boxrule=2.6pt,arc=11pt,
      drop shadow=popink,left=15pt,right=15pt,top=12pt,bottom=11pt,before skip=3pt,after skip=18pt]
    \poppill{#2}{popink}{popcream}\par\vspace{9pt}
    {\raggedright\hyphenpenalty=10000\exhyphenpenalty=10000\popdisplay\fontsize{22}{24}\selectfont\color{popcream}\MakeUppercase{#1}\par}
  \end{tcolorbox}
}
% Section numérotée : pastille orange avec le numéro + titre Archivo
\newcounter{popsec}
\newcommand{\coursec}[1]{%
  \stepcounter{popsec}\setcounter{popsub}{0}%
  \par\vspace{16pt}\noindent
  \begin{minipage}{\linewidth}
  \tikz[baseline=-0.7ex]\node[draw=popink,line width=1.6pt,fill=poporange,rounded corners=4pt,minimum size=22pt,inner sep=0]{\popdisplay\fontsize{12}{12}\selectfont\color{popcream}\thepopsec};%
  \hspace{8pt}{\popdisplay\fontsize{15}{17}\selectfont\color{popink}\MakeUppercase{#1}}\par
  \vspace{4pt}\noindent{\color{poporange}\rule{36pt}{3.6pt}}
  \end{minipage}\par\nopagebreak\vspace{8pt}
}
\newcounter{popsub}
\newcommand{\courssub}[1]{%
  \stepcounter{popsub}%
  \par\vspace{9pt}\noindent{\popxbold\fontsize{11.5}{13}\selectfont\color{popink}{\color{poporange}\thepopsec.\thepopsub}\hspace{6pt}#1}\par\nopagebreak\vspace{4pt}
}

% ============================================================
% Boîtes pédagogiques — même recette : fond blanc, bordure épaisse colorée,
% ombre dure encre, badge plein. Titre optionnel [#1].
% ============================================================
\tcbset{popbase/.style={breakable,enhanced,colback=white,boxrule=2.3pt,arc=9pt,
  shadow={3pt}{-3pt}{0pt}{popink},left=14pt,right=14pt,top=11pt,bottom=11pt,before skip=11pt,after skip=15pt,
  fontupper=\popsemi\fontsize{10.6}{14.5}\selectfont\color{popink},
  fontlower=\popsemi\fontsize{10.6}{14.5}\selectfont\color{popink}}}
% Coche dessinée (le glyphe ✓ n'existe pas dans les polices du document)
\newcommand{\popcheck}[1]{\tikz[baseline=-0.5ex]\draw[#1,line width=1.6pt,line cap=round,line join=round](-0.13,0.0)--(-0.04,-0.09)--(0.13,0.1);}
\providecommand{\coche}{\popcheck{popgreen}}
\providecommand{\resultat}[1]{\par\smallskip\noindent\fcolorbox{popgreen}{popgreenL}{\parbox{\dimexpr\linewidth-2\fboxsep-2\fboxrule\relax}{\textbf{Résultat.} #1}}\par\smallskip}
\newcommand{\popboxhead}[3]{\poppill{#1}{#2}{white}\ifx\relax#3\relax\else\hspace{6pt}{\popxbold\fontsize{10.6}{12}\selectfont\color{popink}#3}\fi\par\vspace{7pt}}

\newtcolorbox{aretenir}[1][]{popbase,colframe=popgreen,before upper={\popboxhead{À retenir}{popgreen}{#1}}}
% Texte en chinois (document de lecture, dialogue, extrait) : cadre
% d'étude. Un titre optionnel (source, auteur) s'affiche dans l'en-tête.
\newtcolorbox{textebox}[1][]{popbase,colframe=popblue,before upper={\popboxhead{文本 · Texte}{popblue}{#1}},after upper={}}
% Marques ✓ / ✗ (forme correcte / fautive) souvent utilisées par les modèles
\providecommand{\cross}{\textcolor{poppink}{\ensuremath{\times}}}
\providecommand{\xmark}{\cross}\providecommand{\crossmark}{\cross}\providecommand{\cmark}{\ensuremath{\checkmark}}
% Ligne de réponse / justification à compléter (inventée par les modèles)
\providecommand{\justification}[1]{\par\noindent\textit{Justification~:} \dotfill\par}
% \textipa (package tipa non chargé) : la police ne couvre pas l'API -> simple passage
\providecommand{\textipa}[1]{#1}
% Soulignement (ulem non chargé) : \uline, \ul
\providecommand{\uline}[1]{\underline{#1}}\providecommand{\ul}[1]{\underline{#1}}
% \glossaire{\item ...} inventée par les modèles : liste de glossaire
\providecommand{\glossaire}[1]{\begin{itemize}#1\end{itemize}}
% \alert (beamer) inventée par les modèles : texte mis en évidence
\providecommand{\alert}[1]{\textbf{\textcolor{poppink}{#1}}}
% variantes ASCII d'ordinaux écrites en macro
\providecommand{\iere}{\textsuperscript{re}}\providecommand{\ieme}{\textsuperscript{e}}\providecommand{\ere}{\textsuperscript{re}}\providecommand{\eme}{\textsuperscript{e}}\providecommand{\er}{\textsuperscript{er}}\providecommand{\ie}{\textsuperscript{e}}
% Ordinaux français écrits en macro par les modèles : 1\ière, 3\ième
\newcommand{\ière}{\textsuperscript{re}}\newcommand{\ième}{\textsuperscript{e}}
% \exemple (inventée par les modèles) : étiquette « Exemple : »
\providecommand{\exemple}{\textbf{Exemple~:}\ }
% \square utilisable aussi hors mode math (cases à cocher)
\AtBeginDocument{\let\squareMath\square\renewcommand{\square}{\ensuremath{\squareMath}}}
% Mot ou phrase en chinois à l'intérieur d'un paragraphe français
\newcommand{\zh}[1]{\ifmmode\text{#1}\else{#1}\fi}
\let\eng\zh \let\esp\zh \let\all\zh % alias tolérés
% flèches d'intonation inventées par les modèles
\providecommand{\textsearrow}{}\renewcommand{\textsearrow}{\ensuremath{\searrow}}\providecommand{\textnearrow}{}\renewcommand{\textnearrow}{\ensuremath{\nearrow}}
% \fr{..} : mot français (inventé par les modèles)
\providecommand{\fr}[1]{\textit{#1}}
% zhbox : encadré de texte chinois inventé par les modèles
\newenvironment{zhbox}{\begin{quote}}{\end{quote}}
% \vs : « versus » inventé par les modèles
\providecommand{\vs}{\textit{vs}\ }
% \blank : case à compléter inventée par les modèles
\providecommand{\blank}[1][]{\underline{\hspace{1.6em}}}
\newtcolorbox{exemplebox}[1][]{popbase,colframe=poporange,before upper={\popboxhead{Exemple}{poporange}{#1}}}
\newtcolorbox{definition}[1][]{popbase,colframe=popblue,before upper={\popboxhead{Définition}{popblue}{#1}}}
\newtcolorbox{propriete}[1][]{popbase,colframe=poppurple,before upper={\popboxhead{Propriété}{poppurple}{#1}}}
\newtcolorbox{methode}[1][]{popbase,colframe=popgold,before upper={\popboxhead{Méthode}{popgold}{#1}}}
\newtcolorbox{attention}[1][]{popbase,colframe=poppink,before upper={\popboxhead{Attention}{poppink}{#1}}}
\newtcolorbox{experience}[1][]{popbase,colframe=popblue,colback=popblueL,before upper={\popboxhead{Expérience}{popblue}{#1}}}
% « Le savais-tu ? » : carte violette pleine (contexte, culture, Cameroun)
\newtcolorbox{savaistu}[1][]{popbase,colback=poppurple,colframe=popink,
  fontupper=\popsemi\fontsize{10.4}{14.2}\selectfont\color{white},
  before upper={\poppill{Le savais-tu\,?}{white}{poppurple}\ifx\relax#1\relax\else\hspace{6pt}{\popxbold\color{white}#1}\fi\par\vspace{7pt}}}
% Exercice résolu : énoncé en haut, \tcblower puis la solution sur fond crème
\newcounter{popexo}\newcounter{popexr}
\newtcolorbox{exoresolu}[1][]{popbase,colframe=poporange,colbacklower=popcream,
  segmentation style={popink,line width=1.2pt,dashed},
  before upper={\stepcounter{popexr}\popboxhead{Exercice résolu \thepopexr}{poporange}{#1}},
  before lower={\poppill{Solution}{popink}{popcream}\par\vspace{6pt}}}
% Exercice (non résolu en place) — la correction est dans la partie Corrigés
\newtcolorbox{exercice}[1][]{popbase,colframe=popink,
  before upper={\stepcounter{popexo}\popboxhead{Exercice \thepopexo}{popink}{#1}}}
% Corrigé
\newtcolorbox{corrige}[1][]{popbase,colframe=popgreen,colback=popgreenL,
  before upper={\popboxhead{Corrigé}{popgreen}{#1}}}
% Objectifs / prérequis / bilan
\newtcolorbox{objectifs}{popbase,colframe=popink,colback=poporangeL,
  before upper={\popboxhead{Objectifs de la leçon}{popink}{}}}
\newtcolorbox{prerequis}{popbase,colframe=popink,colback=popblueL,
  before upper={\popboxhead{Prérequis}{popblue}{}}}
\newtcolorbox{bilan}{popbase,colframe=popink,colback=white,boxrule=2.8pt,
  before upper={\popboxhead{Fiche bilan}{poporange}{L'essentiel à connaître pour l'examen}}}
% Figure encadrée avec légende
\newtcolorbox{popfigure}[1][]{enhanced,breakable=false,colback=white,colframe=popline,boxrule=1.4pt,arc=8pt,
  left=10pt,right=10pt,top=10pt,bottom=8pt,before skip=10pt,after skip=12pt,halign=center,
  lower separated=false,halign lower=center,
  fontlower=\popsemi\fontsize{9}{11.5}\selectfont\color{popmuted},
  #1}
\newcommand{\legende}[1]{\par\vspace{6pt}{\popsemi\fontsize{9}{11.5}\selectfont\color{popmuted}#1}}

% ============================================================
% Signature OnBuch+
% ============================================================
\newcommand{\poplogoinline}[1]{{\poplogo\fontsize{#1}{#1}\selectfont\color{popink}OnBuch\color{poporange}+}}
\newcommand{\signatureonbuch}{%
  \begin{tcolorbox}[enhanced,colback=popink,colframe=popink,boxrule=2.2pt,arc=10pt,
      drop shadow=poporange,left=14pt,right=14pt,top=10pt,bottom=10pt,before skip=0pt,after skip=0pt]
    \tikz[baseline=-0.7ex]\node[circle,fill=popgreen,draw=popcream,line width=1.3pt,minimum size=22pt,inner sep=0]{\popcheck{white}};%
    \hspace{9pt}\begin{minipage}[c]{0.62\linewidth}
      {\popxbold\fontsize{12}{13}\selectfont\color{popcream}Signé\ \textbullet\ {\poplogo OnBuch\color{poporange}+}}\par\vspace{2pt}
      {\raggedright\popsemi\fontsize{8}{10.5}\selectfont\color{popline}\MakeUppercase{Équipe pédagogique OnBuch+}\\\MakeUppercase{Conforme au programme officiel MINESEC}\par}
    \end{minipage}\hfill
    {\poplogo\fontsize{22}{22}\selectfont\color{popcream}OnBuch\color{poporange}+}
  \end{tcolorbox}%
}

% ============================================================
% Page de garde
% #1 titre · #2 matière · #3 niveau · #4 module · #5 n° leçon · #6 durée ·
% #7 description / objectifs (texte ou itemize)
% ============================================================
\newcommand{\pagegarde}[7]{%
  \thispagestyle{empty}
  \newgeometry{a4paper,margin=1.7cm,top=1.6cm,bottom=1.4cm}%
  \begin{tikzpicture}[remember picture,overlay]
    \fill[poporange!18] ([xshift=-3.2cm,yshift=-3.6cm]current page.north east) circle (4.6cm);
    \fill[poppurple!14] ([xshift=2.4cm,yshift=7.6cm]current page.south west) circle (3.8cm);
  \end{tikzpicture}%
  {\poplogo\fontsize{30}{30}\selectfont\color{popink}OnBuch\color{poporange}+}\hfill
  \raisebox{0.6ex}{\poppill{Cours signé \textbullet\ Premium}{popink}{popcream}}\par
  \vspace{1pt}\popsquiggle{poppurple}\par
  \vspace{8pt}
  \begin{tcolorbox}[enhanced,colback=poporange,colframe=popink,boxrule=2.8pt,arc=13pt,
      drop shadow=popink,left=18pt,right=18pt,top=12pt,bottom=13pt,after skip=10pt]
    \poppill{#2 \textbullet\ #3}{popink}{popcream}\hspace{5pt}\poppill{#4}{white}{popink}\par\vspace{8pt}
    {\popdisplay\fontsize{14}{14}\selectfont\color{popink}LEÇON #5}\par\vspace{4pt}
    {\raggedright\hyphenpenalty=10000\exhyphenpenalty=10000\popdisplay\fontsize{25}{27}\selectfont\color{popcream}\MakeUppercase{#1}\par}
    \vspace{8pt}
    {\popxbold\fontsize{9.5}{9.5}\selectfont\color{popink}\MakeUppercase{Durée conseillée : #6}}
  \end{tcolorbox}
  \begin{tcolorbox}[enhanced,colback=white,colframe=popink,boxrule=2.2pt,arc=10pt,
      drop shadow=popink,left=16pt,right=16pt,top=10pt,bottom=10pt,after skip=8pt]
    \poppill{Dans cette leçon}{poppurple}{white}\par\vspace{5pt}
    \popsemi\fontsize{9.3}{12.6}\selectfont\color{popink}#7
  \end{tcolorbox}
  \noindent\begin{minipage}[t]{0.487\linewidth}
    \begin{tcolorbox}[enhanced,colback=popgreen,colframe=popink,boxrule=2.2pt,arc=10pt,
        drop shadow=popink,left=11pt,right=11pt,top=10pt,bottom=10pt,height=3.1cm,valign=center]
      \tcbox[colback=white,colframe=popink,boxrule=1.3pt,arc=5pt,boxsep=0pt,nobeforeafter,
        left=3pt,right=3pt,top=3pt,bottom=3pt,valign=center]{\qrcode[height=1.9cm]{https://play.google.com/store/apps/details?id=com.luvvix.onbuch&hl=fr}}%
      \hspace{8pt}\begin{minipage}[c]{0.5\linewidth}
        {\popdisplay\fontsize{11}{12.5}\selectfont\color{white}\MakeUppercase{Télécharge OnBuch}}\par\vspace{4pt}
        {\raggedright\popsemi\fontsize{8.4}{11}\selectfont\color{white}Gratuit \textbullet\ Google Play\\Tuteur IA, annales, concours, résultats\par}
      \end{minipage}
    \end{tcolorbox}
  \end{minipage}\hfill
  \begin{minipage}[t]{0.487\linewidth}
    \begin{tcolorbox}[enhanced,colback=poppurple,colframe=popink,boxrule=2.2pt,arc=10pt,
        drop shadow=popink,left=11pt,right=11pt,top=10pt,bottom=10pt,height=3.1cm,valign=center]
      \tikz[baseline=-0.7ex]\node[circle,fill=white,draw=popink,line width=1.3pt,minimum size=1.6cm,inner sep=0]{\popdisplay\fontsize{24}{24}\selectfont\color{poppurple}+};%
      \hspace{8pt}\begin{minipage}[c]{0.6\linewidth}
        {\popdisplay\fontsize{11}{12.5}\selectfont\color{white}\MakeUppercase{Passe à OnBuch+}}\par\vspace{4pt}
        {\raggedright\popsemi\fontsize{8.4}{11}\selectfont\color{white}Tous les cours signés, Tuteur IA illimité, fiches PDF — onglet OnBuch+ de l'app\par}
      \end{minipage}
    \end{tcolorbox}
  \end{minipage}
  \vspace{8pt}
  \popcontenu
  \vfill
  \signatureonbuch
  \restoregeometry
  \newpage
}

% Rangée « Ce cours contient » (page de garde)
\newcommand{\poptile}[3]{% #1 couleur #2 gros texte #3 légende
  \begin{tcolorbox}[enhanced,colback=#1,colframe=popink,boxrule=2pt,arc=9pt,shadow={2.5pt}{-2.5pt}{0pt}{popink},
      left=6pt,right=6pt,top=9pt,bottom=9pt,halign=center,valign=center,height=1.75cm,width=\linewidth,nobeforeafter]
    {\popdisplay\fontsize{12.5}{13}\selectfont\color{white}\MakeUppercase{#2}}\par\vspace{3pt}
    {\centering\popsemi\fontsize{7.8}{9.5}\selectfont\color{white}#3\par}
  \end{tcolorbox}}
\newcommand{\popcontenu}{%
  \par\noindent{\popxboldls\fontsize{7.5}{7.5}\selectfont\color{popmuted}CE COURS SIGNÉ CONTIENT}\par\vspace{7pt}
  \noindent\begin{minipage}[t]{0.235\linewidth}\poptile{popblue}{Cours}{complet, illustré, pas à pas}\end{minipage}\hfill
  \begin{minipage}[t]{0.235\linewidth}\poptile{poporange}{Méthodes}{et exercices résolus}\end{minipage}\hfill
  \begin{minipage}[t]{0.235\linewidth}\poptile{poppink}{Exercices}{type examen + corrigés}\end{minipage}\hfill
  \begin{minipage}[t]{0.235\linewidth}\poptile{popgreen}{Fiche bilan}{l'essentiel à retenir}\end{minipage}\par}

% Sommaire
\newtcolorbox{sommaire}{enhanced,colback=white,colframe=popink,boxrule=2.2pt,arc=10pt,
  drop shadow=popink,left=18pt,right=18pt,top=13pt,bottom=13pt,before skip=6pt,after skip=20pt,
  before upper={\poppill{Sommaire}{poppurple}{white}\par\vspace{9pt}\popsemi\fontsize{10.6}{14.5}\selectfont\color{popink}}}

% Signature de fin de cours — poussée tout en bas de la dernière page
\newcommand{\findecours}{%
  \par\vfill\nopagebreak
  \begin{center}\popsquiggle{poporange}\end{center}\vspace{4pt}
  \signatureonbuch
}
\AtBeginDocument{\def\llbracket{\mathopen{[\mkern-3.5mu[}}\def\rrbracket{\mathclose{]\mkern-3.5mu]}}} % intervalles d'entiers (glyphe ⟦ absent de la police)
% Glyphes absents de Fira Math : redéfinis à partir de glyphes disponibles
\AtBeginDocument{%
  \def\setminus{\mathbin{\backslash}}\let\smallsetminus\setminus
  \def\vdots{\mathord{\vbox{\baselineskip3.2pt\lineskiplimit0pt\kern4pt\hbox{.}\hbox{.}\hbox{.}}}}%
  \def\ddots{\mathinner{\mkern1mu\raise6.5pt\hbox{.}\mkern2mu\raise3.6pt\hbox{.}\mkern2mu\raise0.7pt\hbox{.}\mkern1mu}}%
  \def\triangle{\mathord{\scalebox{0.9}{$\symup{\Delta}$}}}%
  \def\nabla{\mathord{\raisebox{\depth}{\scalebox{1}[-1]{$\symup{\Delta}$}}}}%
  \def\checkmark{\popcheck{popdark}}%
  \def\star{\mathbin{\ast}}\def\bigstar{\mathord{\ast}}%
  \def\diamond{\mathbin{\scalebox{0.6}{$\Diamond$}}}%
  \def\bigcup{\mathop{\mathchoice{\raisebox{-0.3em}{\scalebox{1.7}{$\cup$}}}{\scalebox{1.25}{$\cup$}}{\cup}{\cup}}\displaylimits}%
  \def\bigcap{\mathop{\mathchoice{\raisebox{-0.3em}{\scalebox{1.7}{$\cap$}}}{\scalebox{1.25}{$\cap$}}{\cap}{\cap}}\displaylimits}%
  \def\lesseqqgtr{\mathrel{\lessgtr}}\def\bigoplus{\mathop{\oplus}\displaylimits}%
}
\providecommand{\ssi}{si et seulement si} % abréviation fréquente
\AtBeginDocument{\providecommand{\wideparen}{\overparen}} % arc de cercle

%% FIGFIX-BEGIN v4 — correctifs globaux des figures (docs/onbuch-generation/tools/figfix)
\usepackage{adjustbox}\usepackage{etoolbox}
% toute figure TikZ/pgfplots plus large que la ligne est réduite à la largeur disponible
\BeforeBeginEnvironment{tikzpicture}{\begin{adjustbox}{max width=\linewidth}}
\AfterEndEnvironment{tikzpicture}{\end{adjustbox}}
% légendes pgfplots lisibles par-dessus les courbes
\pgfplotsset{every axis/.append style={legend style={fill=white,fill opacity=0.92,text opacity=1,draw=black!15,inner sep=3pt}}}
% étiquettes d'arêtes (syntaxe edge["..."]) avec fond blanc
\tikzset{every edge quotes/.append style={fill=white,inner sep=1.2pt}}
%% FIGFIX-END
\begin{document}

\pagegarde{Leçon 16 — Grammaire : Les compl.de durée ; redoublement des adjectifs}{Chinois}{1ère A}{Module 2 : Environnement et santé}{ZHA16}{2 h}{Cette leçon de grammaire présente deux structures fondamentales du chinois : le complément de durée (exprimer « pendant/depuis combien de temps ») et le redoublement des adjectifs (insister, adoucir ou modifier un adjectif). Schémas, comparaisons avec le français et exercices de transformation permettent d'éviter les erreurs typiques des francophones.
\vspace{4pt}
\begin{multicols}{2}\small
\begin{itemize}
\item À la fin de cette leçon, je sais construire le complément de durée avec un verbe simple (V + durée)
\item je sais placer correctement le complément de durée quand le verbe a un objet (V + O + V + 了 + durée)
\item je sais distinguer « durée d'une action terminée » et « durée qui continue » (avec 了 en fin de phrase)
\item je sais former la négation du complément de durée avec 没(有)
\item je sais redoubler les adjectifs monosyllabiques (AA) et bisyllabiques (AABB)
\item je sais utiliser le redoublement adjectival comme complément circonstanciel avec 地
\end{itemize}\end{multicols}}

\begin{sommaire}
\begin{multicols}{2}
\begin{enumerate}
\item Observation : découvrir le complément de durée
\item Schéma et emplois du complément de durée
\item Négation, exceptions et pièges des francophones
\item Observation : le redoublement des adjectifs
\item Schémas, emplois et limites du redoublement
\item Entraînement intégré : phrases et mini-texte
\item Activité d'intégration
\item Méthodes et exercices résolus
\item Exercices
\item Corrigés des exercices
\item Fiche bilan
\end{enumerate}
\end{multicols}
\end{sommaire}

\begin{objectifs}
\begin{itemize}
\item À la fin de cette leçon, je sais construire le complément de durée avec un verbe simple (V + durée)
\item je sais placer correctement le complément de durée quand le verbe a un objet (V + O + V + 了 + durée)
\item je sais distinguer « durée d'une action terminée » et « durée qui continue » (avec 了 en fin de phrase)
\item je sais former la négation du complément de durée avec 没(有)
\item je sais redoubler les adjectifs monosyllabiques (AA) et bisyllabiques (AABB)
\item je sais utiliser le redoublement adjectival comme complément circonstanciel avec 地
\item je sais comparer ces structures avec le français et éviter les calques
\item je sais employer les deux structures dans des phrases et un court texte sur ma vie quotidienne
\end{itemize}
\end{objectifs}

\begin{prerequis}
\begin{itemize}
\item La particule 了 d'accomplissement (Seconde)
\item Les mots de mesure du temps : 年、月、天、小时、分钟、星期
\item La structure de base Sujet + Verbe + Objet
\item Les adjectifs courants : 好、快、慢、高兴、漂亮、认真
\item L'adverbe 已经 et la négation 没(有)
\end{itemize}
\end{prerequis}

\newpage

\coursec{Observation : découvrir le complément de durée}

\courssub{Situation d'accroche : « Depuis combien de temps ? »}

Imagine la scène : tu es dans la cour du lycée à Yaoundé. Ton ami te demande : « Tu apprends le chinois depuis combien de temps ? » Tu veux répondre fièrement : « J'apprends le chinois depuis deux ans. » En français, c'est facile : on utilise le mot \emph{depuis} et on place la durée à la fin de la phrase. Mais en chinois, impossible de traduire mot à mot ! Il faut une structure spéciale : le \cle{complément de durée}, placé \textbf{après le verbe} :

\zh{我学中文学了两年了。} \emph{(Wǒ xué Zhōngwén xué le liǎng nián le.)} « J'apprends le chinois depuis deux ans. »

Et quand ta mère te dit « Écris bien proprement ! », le chinois, lui, \textbf{redouble l'adjectif} : \zh{你要好好地写字。} \emph{(Nǐ yào hǎohāo de xiězì.)} « Tu dois écrire bien soigneusement. »

Deux structures très fréquentes à l'examen, que nous allons maîtriser aujourd'hui. Dans cette première partie, nous partons de quatre phrases réelles pour \textbf{découvrir par l'observation} ce qu'est le complément de durée, où il se place, et pourquoi le verbe se répète parfois.

\textbf{Problématique :} comment dire en chinois « pendant combien de temps » ou « depuis combien de temps » une action dure ?

\courssub{Lecture guidée : quatre phrases à observer}

Lis attentivement les quatre phrases suivantes. Pour chacune, souligne mentalement le groupe de mots qui indique \textbf{combien de temps} dure l'action.

\begin{textebox}[Phrase modèle n° 1]
我学中文学了两年了。\\
\emph{Wǒ xué Zhōngwén xué le liǎng nián le.}\\
« J'apprends le chinois depuis deux ans. »
\end{textebox}

\begin{exemplebox}[Trois autres phrases à observer]
\zh{他睡了八个小时。} \emph{(Tā shuì le bā ge xiǎoshí.)} « Il a dormi huit heures. »\\[2pt]
\zh{我们上课上了四十分钟。} \emph{(Wǒmen shàngkè shàng le sìshí fēnzhōng.)} « Nous avons fait cours pendant quarante minutes. »\\[2pt]
\zh{她在中国住了三个月。} \emph{(Tā zài Zhōngguó zhù le sān ge yuè.)} « Elle a habité en Chine pendant trois mois. »
\end{exemplebox}

\textbf{Questions d'observation (à faire en classe) :}
\begin{enumerate}
\item Dans chaque phrase, quel groupe de mots exprime la durée ? (\zh{两年}, \zh{八个小时}, \zh{四十分钟}, \zh{三个月})
\item Où se trouve ce groupe : avant ou après le verbe ?
\item Dans les phrases 1 et 3, le verbe apparaît-il une ou deux fois ? Pourquoi, à ton avis ?
\end{enumerate}

\courssub{Premier constat : la durée se place APRÈS le verbe}

En français, la durée se place souvent \emph{avant} le verbe ou en toute fin de phrase, introduite par \emph{depuis}, \emph{pendant}, \emph{ça fait... que} : « \textbf{Depuis deux ans}, j'apprends le chinois » ou « Il a dormi \textbf{huit heures} ». En chinois, la règle est différente et sans exception : le complément de durée se place \textbf{immédiatement après le verbe} (et après la particule \zh{了} quand l'action est accomplie ou en cours depuis un certain temps).

\begin{definition}[Le complément de durée (时量补语)]
Le \cle{complément de durée} indique \textbf{pendant combien de temps} une action s'est déroulée ou se déroule encore. Il est formé d'un \textbf{nombre + une unité de temps} (\zh{两年}, \zh{八个小时}, \zh{四十分钟}, \zh{三个月}) et se place \textbf{après le verbe}, généralement après la particule \zh{了}.
\end{definition}

Observons la correspondance entre les deux langues :

\begin{center}
\begin{tabularx}{\linewidth}{|X|X|X|}
\hline
\rowcolor{popblueL} Français & Chinois (caractères + pinyin) & Place de la durée \\
\hline
J'apprends le chinois depuis deux ans. & 我学中文学了两年了。Wǒ xué Zhōngwén xué le liǎng nián le. & après le verbe répété \\
\hline
Il a dormi huit heures. & 他睡了八个小时。Tā shuì le bā ge xiǎoshí. & après 睡了 \\
\hline
Nous avons fait cours pendant quarante minutes. & 我们上课上了四十分钟。Wǒmen shàngkè shàng le sìshí fēnzhōng. & après le verbe répété \\
\hline
Elle a habité en Chine trois mois. & 她在中国住了三个月。Tā zài Zhōngguó zhù le sān ge yuè. & après 住了 \\
\hline
\end{tabularx}
\end{center}

\begin{aretenir}[L'ordre des mots à retenir]
En chinois, on ne dit jamais « depuis deux ans » au début de la phrase comme en français. La durée vient \textbf{après} le verbe : Sujet + Verbe + 了 + Durée. Le mot français \emph{depuis} n'a pas d'équivalent : il est rendu par la structure elle-même (verbe + \zh{了} + durée + \zh{了} final si l'action continue).
\end{aretenir}

\courssub{Deuxième constat : le verbe se répète quand il y a un objet}

Comparons deux de nos phrases :

\begin{itemize}
\item \zh{他睡了八个小时。} — le verbe \zh{睡} (dormir) n'a \textbf{pas d'objet} : il apparaît une seule fois, suivi directement de la durée.
\item \zh{我学中文学了两年了。} — le verbe \zh{学} (apprendre) a un \textbf{objet} : \zh{中文} (le chinois). Le verbe est alors \textbf{répété} : on dit d'abord \zh{学中文} (verbe + objet), puis on reprend le verbe avec la durée : \zh{学了两年}.
\end{itemize}

C'est une règle fondamentale du chinois : un verbe ne peut pas être suivi à la fois d'un objet et d'un complément de durée. On ne peut donc pas dire directement « apprendre le chinois deux ans » dans cet ordre. La solution : \textbf{répéter le verbe}.

\begin{propriete}[Répétition du verbe devant le complément de durée]
Quand le verbe a un \textbf{objet}, la structure complète est :\\[4pt]
\textbf{Sujet + Verbe + Objet + Verbe + 了 + Durée (+ 了)}\\[4pt]
\zh{我学中文学了两年了。} = je + apprendre + le chinois + apprendre + 了 + deux ans + 了.\\
Quand le verbe n'a \textbf{pas d'objet} (dormir, habiter, attendre seul...), pas de répétition :\\
\textbf{Sujet + Verbe + 了 + Durée} : \zh{他睡了八个小时。}
\end{propriete}

\begin{popfigure}
\begin{tikzpicture}[
  bloc/.style={rectangle, rounded corners=3pt, draw=popblue, fill=popblueL, align=center, minimum height=0.9cm, font=\small},
  blocv/.style={rectangle, rounded corners=3pt, draw=poporange, fill=poporangeL, align=center, minimum height=0.9cm, font=\small},
  fleche/.style={-{Stealth[length=2.5mm]}, thick, popdark}
]
\node[bloc, align=center] (s) at (0,0){Sujet\\我 Wǒ};
\node[blocv, right=0.5cm of s, align=center] (v1){Verbe\\学 xué};
\node[bloc, right=0.5cm of v1, align=center] (o){Objet\\中文 Zhōngwén};
\node[blocv, right=0.5cm of o, align=center] (v2){Verbe répété\\学 xué};
\node[bloc, right=0.5cm of v2] (le) {了 le};
\node[bloc, right=0.5cm of le, fill=poppurpleL, draw=poppurple, align=center] (d){Durée\\两年 liǎng nián};
\node[bloc, right=0.5cm of d] (le2) {了 le};
\draw[fleche, poporange, bend left=45] (v1.north) to node[above, font=\footnotesize\itshape, popdark]{le verbe se répète} (v2.north);
\end{tikzpicture}
\legende{Schéma de la structure S + V + O + V + 了 + Durée (+ 了) : la flèche montre la répétition du verbe.}
\end{popfigure}

\courssub{Troisième constat : une durée qui va du passé jusqu'à maintenant}

Le complément de durée décrit une action qui a commencé dans le passé et qui dure pendant un certain temps, souvent \textbf{jusqu'au moment où l'on parle}. C'est exactement le sens du français « \emph{depuis} deux ans ». Le \zh{了} final, dans \zh{我学中文学了两年了}, signale que l'action \textbf{continue encore} aujourd'hui.

\begin{popfigure}
\begin{tikzpicture}[fleche/.style={-{Stealth[length=2.5mm]}, thick}]
\draw[fleche, popdark] (0,0) -- (10.5,0);
\fill[popgreen] (1,0) circle (2.5pt);
\node[below, align=center, font=\small] at (1,-0.1) {début de l'action\\il y a deux ans};
\draw[very thick, poporange, decorate, decoration={brace, amplitude=5pt, mirror}] (1,-1.2) -- (8.5,-1.2) node[midway, below=7pt, font=\small\itshape]{durée : 两年 liǎng nián};
\fill[popblue] (8.5,0) circle (2.5pt);
\node[below, align=center, font=\small] at (8.5,-0.1) {maintenant\\(l'action continue : 了 final)};
\end{tikzpicture}
\legende{Frise des temps : l'action commence dans le passé, dure deux ans, et continue au moment de la parole.}
\end{popfigure}

\courssub{La question qui appelle cette structure}

Pour demander « depuis / pendant combien de temps ? », le chinois utilise le mot interrogatif \zh{多长时间} (duō cháng shíjiān, « combien de temps »), placé lui aussi \textbf{après le verbe} :

\zh{你学中文学了多长时间？} \emph{(Nǐ xué Zhōngwén xué le duō cháng shíjiān ?)} « Tu apprends le chinois depuis combien de temps ? »

Réponse : \zh{我学中文学了两年了。} « J'apprends le chinois depuis deux ans. »

Autre exemple : \zh{他睡了多长时间？} « Il a dormi combien de temps ? » — \zh{他睡了八个小时。} « Il a dormi huit heures. »

\begin{center}
\begin{tabularx}{\linewidth}{|X|X|X|X|}
\hline
\rowcolor{popblueL} Caractères & Pinyin & Nature & Traduction \\
\hline
多长时间 & duō cháng shíjiān & expression interrogative & combien de temps \\
\hline
年 & nián & nom (unité de temps) & année \\
\hline
月 & yuè & nom (unité de temps) & mois \\
\hline
天 & tiān & nom (unité de temps) & jour \\
\hline
个小时 & ge xiǎoshí & classificateur + nom & heure (soixante minutes) \\
\hline
分钟 & fēnzhōng & nom & minute \\
\hline
住 & zhù & verbe & habiter, loger \\
\hline
睡 & shuì & verbe & dormir \\
\hline
上课 & shàngkè & verbe + objet & suivre un cours, faire cours \\
\hline
\end{tabularx}
\end{center}

\begin{attention}[L'erreur classique du francophone]
Ne dis jamais ✱\zh{我学了两年中文了} à ton niveau : placer l'objet \zh{中文} après la durée est une structure avancée qui obéit à des conditions particulières. La forme \textbf{sûre à l'examen} est celle avec le verbe répété : \zh{我学中文学了两年了。} De même, ne traduis pas \emph{depuis} par un mot séparé : \emph{depuis} est rendu par toute la structure Verbe + \zh{了} + durée + \zh{了}.
\end{attention}

\begin{exoresolu}[Repérer le complément de durée]
Énoncé : dans la phrase \zh{她在中国住了三个月。}, (1) souligne le complément de durée, (2) explique pourquoi le verbe n'est pas répété, (3) traduis en français.
\tcblower
Solution : (1) Le complément de durée est \zh{三个月} (sān ge yuè, « trois mois »), placé après \zh{住了}. (2) Le verbe \zh{住} (habiter) n'a pas d'objet : \zh{在中国} est un complément de lieu introduit par \zh{在}, placé avant le verbe, donc aucune répétition n'est nécessaire. (3) Traduction : « Elle a habité en Chine pendant trois mois. »
\end{exoresolu}

\begin{exercice}[À toi de jouer]
1. Souligne le complément de durée dans chaque phrase : a) \zh{他睡了八个小时。} b) \zh{我们上课上了四十分钟。}\\
2. Remets dans l'ordre : 了 / 我 / 两年 / 学 / 中文 / 学 / 了。\\
3. Traduis en chinois : « Il a dormi huit heures. »
\end{exercice}

\begin{corrige}[Exercice « À toi de jouer »]
1. a) \zh{八个小时} ; b) \zh{四十分钟}. 2. \zh{我学中文学了两年了。} (verbe répété car \zh{中文} est l'objet). 3. \zh{他睡了八个小时。} (Tā shuì le bā ge xiǎoshí.) — pas de répétition car \zh{睡} n'a pas d'objet.
\end{corrige}

\coursec{Leçon 16 — Grammaire : Les compléments de durée ; redoublement des adjectifs}

\courssub{Observation : découvrir le complément de durée}

\begin{textebox}[Dialogue : Au lycée de Yaoundé]
\zh{— 王伟，你学中文学了多长时间？} \emph{Wáng Wěi, nǐ xué Zhōngwén xué le duō cháng shíjiān ?}\\
\zh{— 我学中文学了三年了。你呢？} \emph{Wǒ xué Zhōngwén xué le sān nián le. Nǐ ne ?}\\
\zh{— 我只学了一个学期。昨天我复习复习了两个小时，累死了！} \emph{Wǒ zhǐ xué le yí ge xuéqī. Zuótiān wǒ fùxí fùxí le liǎng ge xiǎoshí, lèi sǐ le !}\\
\zh{— 别急，慢慢来。我们一起去操场跑跑步吧？} \emph{Bié jí, màn man lái. Wǒmen yìqǐ qù cāochǎng pǎo pǎo bù ba ?}
\end{textebox}

\begin{center}
\begin{tabularx}{\linewidth}{|X|X|X|X|}\hline
\rowcolor{popblueL} Caractères & Pinyin & Nature & Traduction \\ \hline
\zh{学期} & xuéqī & nom & semestre \\ \hline
\zh{复习} & fùxí & verbe & réviser \\ \hline
\zh{累死} & lèisǐ & verbe résultatif & être mort de fatigue \\ \hline
\zh{别急} & bié jí & impératif négatif & ne t'inquiète pas / ne te presse pas \\ \hline
\zh{慢慢来} & màn man lái & expression & allons-y doucement, pas à pas \\ \hline
\zh{操场} & cāochǎng & nom & terrain de sport \\ \hline
\zh{跑步} & pǎobù & verbe (VO) & courir, faire du jogging \\ \hline
\end{tabularx}
\end{center}

\begin{aretenir}[Observation]
Dans le dialogue, repérez comment s'exprime la durée : \zh{三年了}, \zh{一个学期}, \zh{两个小时}. La durée suit \emph{toujours} le verbe (éventuellement répété). C'est le \cle{complément de durée}.
\end{aretenir}

\courssub{Schéma et emplois du complément de durée}

\courssub{Schéma 1 : verbe sans objet (verbes intransitifs ou utilisés absolument)}

\begin{exemplebox}[Observation]
\zh{他等了半个小时。} \emph{Tā děng le bàn gè xiǎoshí.} « Il a attendu une demi-heure. »\\
\zh{我们走了两个小时。} \emph{Wǒmen zǒu le liǎng gè xiǎoshí.} « Nous avons marché deux heures. »\\
\zh{爷爷睡了一个下午。} \emph{Yéye shuì le yí gè xiàwǔ.} « Grand-père a dormi tout un après-midi. »\\
\zh{我学了三年中文。} \emph{Wǒ xué le sān nián Zhōngwén.} « J'ai appris le chinois pendant trois ans. » (\zh{中文} est objet, mais structure acceptée pour verbes d'étude).
\end{exemplebox}

\begin{propriete}[Schéma 1 : S + V + 了 + Durée]
\textbf{Sujet + Verbe + 了 + Durée}\\
Le mot-outil \zh{了} (aspect accompli) se place immédiatement après le verbe. La durée (nombre + classificateur + nom de temps) suit.
\begin{itemize}
\item Classificateur \zh{个} obligatoire : \zh{小时 gè xiǎoshí}, \zh{星期 gè xīngqī}, \zh{月 gè yuè}, \zh{分钟 fēnzhōng} (souvent sans \zh{个} à l'oral), \zh{天 tiān} (sans \zh{个}), \zh{年 nián} (sans \zh{个}).
\item \zh{半} « demi » précède le classificateur : \zh{半个小时}, \zh{半年}.
\item \zh{两 liǎng} (et non \zh{二 èr}) pour « deux » devant classificateur.
\end{itemize}
\end{propriete}

\begin{center}
\begin{tabularx}{\linewidth}{|X|X|X|}\hline
\rowcolor{popblueL} Structure & Exemple (caractères + pinyin) & Traduction \\ \hline
S + V + 了 + Durée & \zh{他等了半个小时。} Tā děng le bàn gè xiǎoshí. & Il a attendu une demi-heure. \\ \hline
S + V + 了 + Durée & \zh{我睡了八个小时。} Wǒ shuì le bā gè xiǎoshí. & J'ai dormi huit heures. \\ \hline
S + V + 了 + Durée & \zh{他们住了两年。} Tāmen zhù le liǎng nián. & Ils ont habité deux ans. \\ \hline
\end{tabularx}
\end{center}

\begin{attention}[Piège n°1 : La durée ne se place JAMAIS avant le verbe]
En français : « Pendant trois ans, j'ai appris le chinois » (durée en tête possible).\\
En chinois : \textbf{Impossible}. \zh{三年我学了中文} $\rightarrow$ \textbf{Incorrect}.\\
Seule structure : \zh{我学了三年中文。} (Durée \emph{post-verbal}).
\end{attention}

\courssub{Schéma 2 : verbe transitif avec objet (ou verbe VO / 离合词)}

\begin{exemplebox}[Observation]
\zh{她看电视看了一个小时。} \emph{Tā kàn diànshì kàn le yí gè xiǎoshí.} « Elle a regardé la télé une heure. »\\
\zh{他学中文学了两年。} \emph{Tā xué Zhōngwén xué le liǎng nián.} « Il a appris le chinois deux ans. »\\
\zh{我做作业做了三个小时。} \emph{Wǒ zuò zuòyè zuò le sān gè xiǎoshí.} « J'ai fait mes devoirs trois heures. »\\
\zh{他睡觉睡了八个小时。} \emph{Tā shuìjiào shuì le bā gè xiǎoshí.} « Il a dormi huit heures. » (Verbe VO \zh{睡觉} : on répète \zh{睡} seulement).
\end{exemplebox}

\begin{propriete}[Schéma 2 : S + V + O + V + 了 + Durée]
Si le verbe a un objet (nom ou verbe VO), on \cle{répète le verbe} (seul, sans l'objet) avant \zh{了} + durée.\\
\textbf{Sujet + V + Objet + V + 了 + Durée}
\end{propriete}

\begin{attention}[Piège n°2 : Erreur typique du francophone]
$\times$ \zh{她看了电视一个小时。} (Durée collée après l'objet : incorrect).\\
$\times$ \zh{他睡觉了八个小时。} (\zh{了} après le verbe VO complet : incorrect).\\
$\checkmark$ \zh{她看电视看了一个小时。} $\checkmark$ \zh{他睡觉睡了八个小时。}\\
\textbf{Règle d'or :} « Verbe + Objet, puis Verbe + 了 + Durée ».
\end{attention}

\begin{propriete}[Variante : Durée comme attribut (的-structure)]
Quand l'objet est court et très lié au verbe, on peut dire :\\
\textbf{S + V + 了 + Durée + 的 + Objet}\\
\zh{我看了一个小时的电视。} \emph{Wǒ kàn le yí gè xiǎoshí de diànshì.}\\
\zh{他学了两年的中文。} \emph{Tā xué le liǎng nián de Zhōngwén.}\\
\emph{Au lycée, maîtrisez d'abord la répétition du verbe (schéma 2), forme canonique à l'écrit et à l'oral formel.}
\end{propriete}

\courssub{Schéma 3 : Action terminée vs Action qui continue (le double 了)}

\begin{exemplebox}[Observation]
\zh{她在中国住了三个月。} \emph{Tā zài Zhōngguó zhù le sān gè yuè.} « Elle a habité en Chine 3 mois » (fini, elle est partie).\\
\zh{她在中国住了三个月了。} \emph{Tā zài Zhōngguó zhù le sān gè yuè le.} « Cela fait 3 mois qu'elle habite en Chine » (elle y est encore).\\
\zh{他已经工作了一天了。} \emph{Tā yǐjing gōngzuò le yì tiān le.} « Il travaille déjà depuis une journée entière. »
\end{exemplebox}

\begin{propriete}[Le double 了 : marqueur d'état actuel]
\begin{itemize}
\item \textbf{Un seul 了} (après V) : Action \emph{terminée}, durée close. \zh{他等了半个小时。}
\item \textbf{了...了} (après V + fin de phrase) : Action \emph{en cours}, durée écoulée depuis le début jusqu'à maintenant. \zh{他等了半个小时了。}
\item \textbf{已经 yǐjing} + ... \zh{了} final : Insiste sur « déjà », renforce l'aspect continu. \zh{我已经学中文学了三年了。}
\end{itemize}
\end{propriete}

\begin{center}
\begin{tabularx}{\linewidth}{|X|X|}\hline
\rowcolor{popblueL} Action terminée (un 了) & Action qui continue (了...了 / 已经...了) \\ \hline
\zh{他等了半个小时。} Il a attendu 30 min (fini). & \zh{他等了半个小时了。} Il attend depuis 30 min (encore). \\ \hline
\zh{我学了两年中文。} J'ai appris le chinois 2 ans. & \zh{我学中文学了两年了。} J'apprends le chinois depuis 2 ans. \\ \hline
\zh{她睡了一个下午。} Elle a dormi un aprèm. & \zh{她已经睡了一个下午了。} Elle dort depuis un aprèm. \\ \hline
\end{tabularx}
\end{center}

\begin{savaistu}[La logique chinoise : Action d'abord, mesure ensuite]
Le français anticipe : « \emph{Depuis} deux ans, j'apprends... ». Le chinois raconte : « J'apprends le chinois, [action accomplie] deux ans, [état actuel] \zh{了} ». Pensez : \zh{我 学 中文 学 了 两年 了} (Sujet - V - O - V - 了 - Durée - 了).
\end{savaistu}

\courssub{Négation du complément de durée}

\begin{propriete}[Négation : 没(有) + V (+ O) + V + Durée]
Pour nier une durée passée (action non réalisée ou durée non atteinte) :\\
\zh{我没有等他。} \emph{Wǒ méiyǒu děng tā.} « Je ne l'ai pas attendu. »\\
\zh{我没等他等那么久。} \emph{Wǒ méi děng tā děng nàme jiǔ.} « Je ne l'ai pas attendu aussi longtemps. »\\
\zh{他没学中文学两年。} \emph{Tā méi xué Zhōngwén xué liǎng nián.} « Il n'a pas appris le chinois pendant deux ans. »
\end{propriete}

\begin{propriete}[Négation de l'action continue : 还没 + V (+ O) + V + Durée + 了]
« Pas encore [durée] que... » (l'action continue mais n'a pas atteint cette durée).\\
\zh{我还没学中文学一年了。} \emph{Wǒ hái méi xué Zhōngwén xué yì nián le.} « Cela ne fait pas encore un an que j'apprends le chinois. »\\
\zh{他还没睡觉睡八个小时了。} \emph{Tā hái méi shuìjiào shuì bā gè xiǎoshí le.} « Il ne dort pas encore depuis 8 heures. »
\end{propriete}

\begin{attention}[Piège n°3 : Ne pas confondre 没 et 不]
$\times$ \zh{我不等了半个小时。} (\zh{不} nie l'habitude ou le futur, pas le passé accompli avec durée).\\
$\checkmark$ \zh{我没等半个小时。} / \zh{我没有等半个小时。}
\end{attention}

\courssub{Poser la question : 多长时间 / 多久}

\begin{methode}[Transformer la durée en question]
\begin{enumerate}
\item Phrase affirmative : \zh{你等他等了半个小时。}
\item Repérer la durée : \zh{半个小时}.
\item Remplacer par \zh{多长时间} (plus formel, « quelle longueur de temps ») ou \zh{多久} (court, « combien de temps ») : \zh{你等他等了多久？}
\item Réponse courte : durée seule (+ \zh{了} si action continue). \zh{两年了。} / \zh{三天。}
\end{enumerate}
\end{methode}

\begin{exemplebox}[Questions / Réponses]
\zh{你学中文学了多长时间？} \emph{Nǐ xué Zhōngwén xué le duō cháng shíjiān ?} — \zh{两年了。} \emph{Liǎng nián le.}\\
\zh{你睡觉睡了多久？} \emph{Nǐ shuìjiào shuì le duōjiǔ ?} — \zh{八个小时。} \emph{Bā gè xiǎoshí.}\\
\zh{他在雅温得住了多长时间？} \emph{Tā zài Yǎwēndé zhù le duō cháng shíjiān ?} — \zh{五年。} \emph{Wǔ nián.}
\end{exemplebox}

\courssub{Organigramme de décision : construire sa phrase}

\begin{popfigure}
\begin{tikzpicture}[
  node distance=0.9cm and 1.6cm,
  question/.style={diamond, draw=popblue, fill=popblueL, align=center, inner sep=1.5pt, aspect=2.2, font=\small},
  reponse/.style={rectangle, rounded corners, draw=popgreen, fill=popgreenL, align=center, font=\small, text width=4.8cm},
  fleche/.style={-{Stealth}, thick, popdark}]
\node[question, align=center] (q){Le verbe a-t-il\\un objet (ou est-ce un verbe VO) ?};
\node[reponse, below left=of q, align=center] (non){\textbf{NON (V seul)}\\S + V + 了 + Durée\\\zh{他等了半个小时。}};
\node[reponse, below right=of q, align=center] (oui){\textbf{OUI (V + O / VO)}\\Répéter le verbe seul\\S + V + O + V + 了 + Durée\\\zh{她看电视看了一个小时。}\\\zh{他睡觉睡了八个小时。}};
\draw[fleche] (q) -- node[above left, font=\small]{non} (non);
\draw[fleche] (q) -- node[above right, font=\small]{oui} (oui);
\node[rectangle, rounded corners, draw=poporange, fill=poporangeL, align=center, font=\small, below=3.0cm of q, text width=8cm] (suite) {L'action continue jusqu'à maintenant ?\\Ajouter \zh{了} final (ou \zh{已经}... \zh{了}) : \zh{我学中文学了两年了。}};
\draw[fleche, dashed] (non.south) |- (suite.west);
\draw[fleche, dashed] (oui.south) |- (suite.east);
\end{tikzpicture}
\legende{Organigramme : construire une phrase avec complément de durée.}
\end{popfigure}

\begin{center}
\begin{tabularx}{\linewidth}{|X|X|X|X|}\hline
\rowcolor{popblueL} Caractères & Pinyin & Nature & Traduction \\ \hline
\zh{多长时间} & duō cháng shíjiān & expr. interrogative & combien de temps \\ \hline
\zh{多久} & duōjiǔ & pron. interrogatif & combien de temps \\ \hline
\zh{已经} & yǐjing & adverbe & déjà \\ \hline
\zh{半个小时} & bàn gè xiǎoshí & groupe de durée & une demi-heure \\ \hline
\zh{等} & děng & verbe & attendre \\ \hline
\zh{住} & zhù & verbe & habiter \\ \hline
\zh{睡觉} & shuìjiào & verbe (VO) & dormir \\ \hline
\zh{跑步} & pǎobù & verbe (VO) & courir \\ \hline
\zh{复习} & fùxí & verbe & réviser \\ \hline
\end{tabularx}
\end{center}

\begin{exoresolu}[Application : Construire et interroger]
Énoncé. (1) Traduisez : « Elle a habité en Chine trois mois. » (2) Traduisez : « Depuis combien de temps apprends-tu le chinois ? » Répondez : « Deux ans (et je continue). » (3) Niez : « Il n'a pas couru pendant une heure. » (Verbe \zh{跑步}).
\tcblower
Solution détaillée.\\
(1) \zh{住} « habiter » : \zh{在中国} est compl. de lieu (avant V). Pas d'objet. Schéma 1. \zh{她在中国住了三个月。} \emph{Tā zài Zhōngguó zhù le sān gè yuè.} (\zh{月} prend \zh{个}).\\
(2) \zh{学} + objet \zh{中文} $\rightarrow$ répétition V. Question : \zh{你学中文学了多长时间？} \emph{Nǐ xué Zhōngwén xué le duō cháng shíjiān ?} Réponse (continue) : \zh{两年了。} \emph{Liǎng nián le.}\\
(3) \zh{跑步} est verbe VO. Négation passé : \zh{没(有)} + V + O + V + Durée. \zh{他没跑步跑一个小时。} \emph{Tā méi pǎobù pǎo yí gè xiǎoshí.} (Ou \zh{他没有跑步跑一个小时。})
\end{exoresolu}

\begin{exercice}[Entraînement : Compléments de durée]
Traduisez en chinois :\\
(1) « Nous avons marché deux heures. »\\
(2) « Il étudie le chinois depuis trois ans (il continue). »\\
(3) « Combien de temps as-tu regardé la télévision ? »\\
(4) Répondez à (3) : « Une heure. »\\
(5) « Je n'ai pas révisé pendant deux heures hier. » (\zh{复习}, \zh{昨天} en début de phrase).\\
(6) « Cela ne fait pas encore six mois qu'il habite à Douala. » (\zh{住}, \zh{杜阿拉 Dù'ālā}).
\end{exercice}

\begin{corrige}[Exercice « Entraînement : Compléments de durée »]
(1) \zh{我们走了两个小时。} \emph{Wǒmen zǒu le liǎng gè xiǎoshí.} (V seul, Schéma 1).\\
(2) \zh{他学中文学了三年了。} \emph{Tā xué Zhōngwén xué le sān nián le.} (Objet \zh{中文}, répétition V ; \zh{了} final = continue).\\
(3) \zh{你看电视看了多长时间？} \emph{Nǐ kàn diànshì kàn le duō cháng shíjiān ?} (Répétition \zh{看}, \zh{多长时间}).\\
(4) \zh{一个小时。} \emph{Yí gè xiǎoshí.} (Réponse courte).\\
(5) \zh{我昨天没复习复习两个小时。} \emph{Wǒ zuótiān méi fùxí fùxí liǎng gè xiǎoshí.} (Négation \zh{没}, répétition V \zh{复习}).\\
(6) \zh{他还没在杜阿拉住六个月了。} \emph{Tā hái méi zài Dù'ālā zhù liù gè yuè le.} (Négation continue \zh{还没...了}, lieu avant V). \textit{Note : \zh{六个月} prend \zh{个}.}
\end{corrige}

\courssub{Observation : le redoublement des adjectifs (形容词重叠)}

\begin{textebox}[Description : Le temps à Douala et à Kunming]
\zh{杜阿拉的天气热热的，湿湿的，让人觉得很不舒服。} \emph{Dù'ālā de tiānqì rè re de, shī shi de, ràng rén juéde hěn bù shūfu.}\\
\zh{昆明的天气凉凉快快的，干干净净的，非常舒服。} \emph{Kūnmíng de tiānqì liáng liang kuai kuai de, gān gān jìng jìng de, fēicháng shūfu.}
\end{textebox}

\begin{aretenir}[Observation]
En chinois, on peut \textbf{redoubler} un adjectif pour en nuancer le sens : intensité atténuée, aspect visuel, impression subjective, ou valeur adverbiale. On distingue les adjectifs \textbf{monosyllabiques} (1 caractère) et \textbf{disyllabiques} (2 caractères).
\end{aretenir}

\courssub{Schémas et emplois du redoublement}

\courssub{Adjectifs monosyllabiques (A) : forme AA ou AABB}

\begin{propriete}[Redoublement monosyllabique : AA / AABB + 的 / 地]
\begin{itemize}
\item \textbf{Forme AA + 的} (attributif / épithète) : « un peu A », « assez A », aspect visuel. \zh{红红的苹果} (pommes rougeâtres/rouges), \zh{热热的天气} (temps chaud).
\item \textbf{Forme AA + 地} (adverbial) : « d'une manière A ». \zh{他慢慢地走。} \emph{Tā màn man de zǒu.} « Il marche lentement. »
\item \textbf{Forme AABB + 的} (plus expressif, imagé) : \zh{高高大大的} (grand et costaud), \zh{白白胖胖的} (blanc et dodou).
\item \textbf{Forme AABB} (sans 的/地, prédicatif souvent) : \zh{天气冷冷清清。} « Le temps est froid et désert. »
\end{itemize}
\end{propriete}

\begin{center}
\begin{tabularx}{\linewidth}{|X|X|X|}\hline
\rowcolor{popblueL} Adjectif de base & Redoublement (AA / AABB) & Exemple + Traduction \\ \hline
\zh{大 dà} & \zh{大大的 dà da de} & \zh{大大的房子} une grande maison (imposante) \\ \hline
\zh{小 xiǎo} & \zh{小小的 xiǎo xiao de} & \zh{小小的鸟} un tout petit oiseau \\ \hline
\zh{红 hóng} & \zh{红红的 hóng hóng de} & \zh{红红的脸} un visage tout rouge \\ \hline
\zh{慢 màn} & \zh{慢慢地 màn man de} & \zh{慢慢地吃} manger lentement \\ \hline
\zh{高 gāo / 大 dà} & \zh{高高大大 gāo gāo dà dà} & \zh{高高大大的树} un arbre grand et costaud \\ \hline
\zh{干净 gānjìng} & \zh{干干净净 gān gān jìng jìng} & \zh{干干净净的屋子} une pièce propre nette \\ \hline
\end{tabularx}
\end{center}

\courssub{Adjectifs disyllabiques (AB) : forme AABB}

\begin{propriete}[Redoublement disyllabique : AABB (+ 的)]
Seuls certains adjectifs disyllabiques (souvent d'état, de qualité, ou onomatopéiques) se redoublent en AABB. Structure : \textbf{AABB + 的} (attributif).
\begin{itemize}
\item \zh{干净 gānjìng} $\rightarrow$ \zh{干干净净 gān gān jìng jìng} (propre, net, bien rangé).
\item \zh{清楚 qīngchu} $\rightarrow$ \zh{清清楚楚 qīng qīng chu chu} (clair, distinct).
\item \zh{高兴 gāoxìng} $\rightarrow$ \zh{高高兴兴 gāo gāo xìng xìng} (joyeusement, de bon cœur).
\item \zh{安静 ānjìng} $\rightarrow$ \zh{安安静静 ān ān jìng jìng} (calme, tranquillement).
\item \zh{漂亮 piàoliang} $\rightarrow$ \zh{漂漂亮亮 piào piào liang liang} (joli, bien fait).
\end{itemize}
\emph{Attention :} Tous les adjectifs disyllabiques ne se redoublent pas (ex: \zh{重要}, \zh{必要}, \zh{正确} ne se redoublent pas).
\end{propriete}

\begin{exemplebox}[Exemples en contexte]
\zh{请把字写工工整整的。} \emph{Qǐng bǎ zì xiě gōng gōng zhěng zhěng de.} « Écris les caractères proprement/soigneusement. »\\
\zh{他高高兴兴地去踢足球了。} \emph{Tā gāo gāo xìng xìng de qù tī zúqiú le.} « Il est allé jouer au foot joyeusement. »\\
\zh{雅温得的天气热热的，不像昆明凉凉快快的。} \emph{Yǎwēndé de tiānqì rè re de, bù xiàng Kūnmíng liáng liang kuai kuai de.}
\end{exemplebox}

\begin{attention}[Pièges du redoublement pour francophones]
\begin{itemize}
\item \textbf{Tons :} Le redoublement suit les règles de \cle{changement de ton (变调)}. Ex: \zh{大 dà (4e)} $\rightarrow$ \zh{大大 dà da} (souvent 4e + ton neutre/léger). \zh{小 xiǎo (3e)} $\rightarrow$ \zh{小小 xiǎo xiao} (2e + 3e ou neutre). À l'oral, le 2e caractère tend vers le ton neutre.
\item \textbf{的 vs 地 :} \zh{红红的苹果} (adj/épithète = \zh{的}) ; \zh{慢慢地走} (adverbe = \zh{地}). Ne les confondez pas.
\item \textbf{Pas de "très" (很) :} On ne dit pas \zh{很红红的}. Le redoublement \emph{remplace} l'intensificateur.
\item \textbf{Limites :} Verbes d'état non qualitatifs (\zh{是}, \zh{有}, \zh{在}), adjectifs abstraits/formels (\zh{必要}, \zh{可能}) ne se redoublent pas.
\end{itemize}
\end{attention}

\courssub{Entraînement intégré : phrases et mini-texte}

\begin{exoresolu}[Exercices mixtes : Durée + Redoublement]
Énoncé. Traduisez en chinois :\\
(1) « Il a couru pendant 40 minutes. » (\zh{跑步}, \zh{分钟 fēnzhōng} sans \zh{个}).\\
(2) « La petite fille est toute rouge (visage). » (\zh{小女孩}, \zh{脸}, \zh{红}).\\
(3) « Depuis combien de temps dors-tu ? » (\zh{睡觉}). Réponse : « 20 minutes (et je continue). »\\
(4) « Écrivez les caractères proprement (工整). » (\zh{写}, \zh{字}, \zh{工整}).
\tcblower
Solution détaillée.\\
(1) \zh{跑步} verbe VO $\rightarrow$ répétition V \zh{跑}. \zh{他跑步跑了四十分钟。} \emph{Tā pǎobù pǎo le sìshí fēnzhōng.}\\
(2) Adj. monosyllabique \zh{红} attributif $\rightarrow$ AA + 的. \zh{小女孩的脸红红的。} \emph{Xiǎo nǚhái de liǎn hóng hóng de.}\\
(3) Question : \zh{你睡觉睡了多久？} \emph{Nǐ shuìjiào shuì le duōjiǔ ?} Réponse (continue) : \zh{二十分钟了。} \emph{Èrshí fēnzhōng le.}\\
(4) Adj. disyllabique \zh{工整} $\rightarrow$ AABB + 地 (adverbial). \zh{请把字写工工整整的。} \emph{Qǐng bǎ zì xiě gōng gōng zhěng zhěng de.} (Souvent \zh{的} à l'oral pour adverbe de manière, \zh{地} formel. Les deux acceptés).
\end{exoresolu}

\begin{exercice}[Production écrite guidée : Mon environnement]
Rédigez un court paragraphe (5-6 phrases) en chinois pour décrire votre ville/village (Yaoundé, Douala, Buea, etc.) et vos habitudes d'étude. Utilisez \textbf{au moins 2 compléments de durée} (dont un avec action continue \zh{了...了}) et \textbf{2 adjectifs redoublés} (monosyllabique et disyllabique).
\vspace{0.5cm}
\textbf{Mots-outils suggérés :} \zh{天气}, \zh{热/冷}, \zh{干净}, \zh{学习}, \zh{复习}, \zh{跑步}, \zh{睡觉}, \zh{已经}, \zh{还没}.
\end{exercice}

\begin{corrige}[Exercice « Production écrite guidée » — Proposition de correction]
\zh{我住在雅温得。这里的天气热热的，湿湿的，不像昆明凉凉快快的。我学中文学了两年了，每天复习复习一个小时。我喜欢跑步，昨天跑步跑了四十分钟。晚上我睡觉睡了八个小时了，很舒服。}\\
\emph{Wǒ zhù zài Yǎwēndé. Zhèlǐ de tiānqì rè re de, shī shi de, bù xiàng Kūnmíng liáng liang kuai kuai de. Wǒ xué Zhōngwén xué le liǎng nián le, měitiān fùxí fùxí yí gè xiǎoshí. Wǒ xǐhuan pǎobù, zuótiān pǎobù pǎo le sìshí fēnzhōng. Wǎnshang wǒ shuìjiào shuì le bā gè xiǎoshí le, hěn shūfu.}
\end{corrige}

\begin{aretenir}[L'essentiel de la Leçon 16]
\begin{itemize}
\item \textbf{Complément de durée :} Toujours \emph{après} le verbe.
\item \textbf{Sans objet :} S + V + 了 + Durée. \zh{他等了半个小时。}
\item \textbf{Avec objet / Verbe VO :} S + V + O + V + 了 + Durée. \zh{她看电视看了一个小时。} / \zh{他睡觉睡了八个小时。}
\item \textbf{Action continue :} 了...了 / 已经...了. \zh{我学中文学了两年了。}
\item \textbf{Négation passée :} 没(有) + V (+ O) + V + Durée.
\item \textbf{Négation continue :} 还没 + V (+ O) + V + Durée + 了.
\item \textbf{Question :} \zh{多长时间} / \zh{多久} remplace la durée.
\item \textbf{Redoublement Adj. Mono (A) :} AA + 的/地 (\zh{红红的}, \zh{慢慢地}) ; AABB + 的 (\zh{高高大大的}).
\item \textbf{Redoublement Adj. Di (AB) :} AABB + 的/地 (\zh{干干净净的}, \zh{高高兴兴地}). Pas tous les adj. disyllabiques.
\item \textbf{Pas de 很} devant adjectif redoublé.
\end{itemize}
\end{aretenir}

\coursec{Négation, exceptions et pièges des francophones}

Dans la section précédente, nous avons découvert le complément de durée : la durée se place \emph{après} le verbe, souvent précédée de \zh{了} (le), et le verbe se répète lorsqu'il a un objet (\zh{我学中文学了两年。} Wǒ xué Zhōngwén xué le liǎng nián. « J'ai étudié le chinois pendant deux ans. »). Mais comment dire « il n'a \emph{pas} dormi huit heures » ? Et que faire des verbes comme \zh{睡觉} (shuìjiào, dormir) qui se « coupent » en deux ? C'est l'objet de cette section : la négation, les cas particuliers, et surtout les trois pièges dans lesquels tombent presque tous les élèves francophones.

\courssub{La négation avec \zh{没(有)} : la règle d'or}

\textbf{Observation.} Lisez et comparez les paires de phrases suivantes :

\begin{exemplebox}[Affirmation et négation : observez la différence]
\zh{他睡了八个小时。} \emph{Tā shuì le bā ge xiǎoshí.} « Il a dormi huit heures. »\\
\zh{他没睡八个小时。} \emph{Tā méi shuì bā ge xiǎoshí.} « Il n'a pas dormi huit heures. »\\[4pt]
\zh{我学了两年中文。} \emph{Wǒ xué le liǎng nián Zhōngwén.} « J'ai étudié le chinois pendant deux ans. »\\
\zh{我没学两年中文。} \emph{Wǒ méi xué liǎng nián Zhōngwén.} « Je n'ai pas étudié le chinois pendant deux ans. »
\end{exemplebox}

Que remarquez-vous ? À la forme négative, deux changements se produisent : \zh{没} (méi) apparaît devant le verbe, et \zh{了} \emph{disparaît}. C'est une règle générale du chinois : \zh{没(有)} nie un accompli, et \zh{了} marque l'accompli ; les deux ne peuvent donc pas cohabiter.

\begin{propriete}[Négation du complément de durée]
Structure : \textbf{Sujet + 没(有) + Verbe + durée (+ objet)} — \emph{sans} \zh{了}.\\
\zh{他没(有)睡八个小时。} \emph{Tā méi(yǒu) shuì bā ge xiǎoshí.} « Il n'a pas dormi huit heures. »\\
La particule \zh{了} est toujours supprimée après \zh{没(有)}. Dire ✱\zh{他没睡了八个小时} est une faute grave.
\end{propriete}

\begin{attention}[Avec \zh{没(有)}, on supprime \zh{了}]
\zh{他没睡八个小时。} \emph{Tā méi shuì bā ge xiǎoshí.} « Il n'a pas dormi huit heures. »\\
Ne dites jamais ✱\zh{他没睡了八个小时。} — \zh{没} et \zh{了} sont incompatibles : l'un nie l'accomplissement, l'autre l'affirme.
\end{attention}

\courssub{La structure spéciale \zh{有 + durée + 没 + V + 了}}

Le chinois possède une tournure très fréquente pour dire « cela fait X (temps) que je n'ai pas… ». Observez :

\begin{exemplebox}[« Cela fait… que je n'ai pas… »]
\zh{我有两年没学中文了。} \emph{Wǒ yǒu liǎng nián méi xué Zhōngwén le.} « Cela fait deux ans que je n'ai pas étudié le chinois. »\\
\zh{我们有三天没上课了。} \emph{Wǒmen yǒu sān tiān méi shàngkè le.} « Cela fait trois jours que nous n'avons pas eu cours. »\\
\zh{她有一个月没回家了。} \emph{Tā yǒu yí ge yuè méi huí jiā le.} « Cela fait un mois qu'elle n'est pas rentrée à la maison. »
\end{exemplebox}

\begin{propriete}[Structure 有 + durée + 没 + Verbe + 了]
Schéma : \textbf{Sujet + 有 + durée + 没 + Verbe + 了}.\\
Sens : « cela fait (durée) que (sujet) n'a pas (verbe) ». Ici, le \zh{了} final est le \zh{了} de changement de situation (nouvelle situation constatée au moment de la parole) ; il est autorisé car il ne se colle pas au verbe nié. La durée, portée par \zh{有}, se trouve \emph{avant} \zh{没}, mais toujours \emph{après} un verbe (\zh{有}) : la règle « la durée suit un verbe » reste respectée.
\end{propriete}

\courssub{Cas particulier : les verbes séparables (离合词)}

Certains « verbes » chinois sont en réalité des blocs verbe + objet soudés, appelés \cle{verbes séparables} (\zh{离合词} líhécí) : \zh{睡觉} (shuìjiào, dormir = « dormir un sommeil »), \zh{上课} (shàngkè, avoir cours), \zh{吃饭} (chīfàn, manger = « manger du riz »). Avec une durée, on peut \emph{insérer} le complément au milieu :

\begin{exemplebox}[Durée insérée dans le verbe séparable]
\zh{他睡了八个小时的觉。} \emph{Tā shuì le bā ge xiǎoshí de jiào.} « Il a dormi huit heures. » (litt. « il a dormi huit heures de sommeil »)\\
\zh{我们上了两个小时的课。} \emph{Wǒmen shàng le liǎng ge xiǎoshí de kè.} « Nous avons eu deux heures de cours. »
\end{exemplebox}

\begin{aretenir}[Forme avancée, hors programme strict]
La forme \zh{睡了八个小时的觉} (verbe + 了 + durée + 的 + objet) est correcte et très naturelle, mais elle est \textbf{hors programme strict} en Première : il suffit de la reconnaître à la lecture. À l'écrit, utilisez la forme simple : \zh{他睡了八个小时。} ou la répétition du verbe : \zh{他睡觉睡了八个小时。}
\end{aretenir}

\courssub{Les trois pièges des francophones}

\begin{propriete}[Piège n°1 : la durée ne se place JAMAIS avant le verbe]
En français, on dit « \emph{depuis deux ans}, j'étudie le chinois » : la durée peut ouvrir la phrase. En chinois, c'est impossible : la durée est un \emph{complément}, elle suit le verbe.\\
✱\zh{两年我学中文。} → \zh{我学中文学了两年了。} \emph{Wǒ xué Zhōngwén xué le liǎng nián le.} « Cela fait deux ans que j'étudie le chinois. »\\
C'est l'erreur n°1 des francophones : ne calquez jamais l'ordre français.
\end{propriete}

\begin{attention}[Piège n°2 : oublier de répéter le verbe devant un objet]
✱\zh{我学中文了两年了。} est à éviter au stade de cette leçon : avec un objet (\zh{中文}), la forme sûre est la répétition du verbe :\\
\zh{我学中文学了两年了。} \emph{Wǒ xué Zhōngwén xué le liǎng nián le.} « J'étudie le chinois depuis deux ans. »\\
Autre solution correcte : placer l'objet après la durée avec \zh{的} : \zh{我学了两年的中文。}
\end{attention}

\begin{attention}[Piège n°3 : garder \zh{了} avec la négation \zh{没}]
✱\zh{我没学了两年。} est fautif : après \zh{没(有)}, on supprime \zh{了}.\\
Forme correcte : \zh{我没学两年中文。} \emph{Wǒ méi xué liǎng nián Zhōngwén.} « Je n'ai pas étudié le chinois pendant deux ans. »
\end{attention}

\begin{popfigure}
\begin{tikzpicture}[mindmap, grow cyclic, every node/.style={concept, align=center, font=\small}, concept color=popblue!40, level 1 concept/.append style={level distance=3.6cm, sibling angle=120}]
\node[concept color=poporange!50, align=center]{\textbf{Complément}\\\textbf{de durée}\\pièges}
  child[concept color=poppink!50]{ node[align=center]{Piège 1\\durée avant\\le verbe\\✱ 两年我学中文} }
  child[concept color=poppink!50]{ node[align=center]{Piège 2\\verbe non répété\\avec objet\\✱ 我学中文了两年了} }
  child[concept color=poppink!50]{ node[align=center]{Piège 3\\了 gardé après 没\\✱ 我没学了两年} };
\end{tikzpicture}
\legende{Carte mentale des trois pièges : chaque branche porte une phrase fautive à bannir.}
\end{popfigure}

\courssub{Comparaison systématique français / chinois}

En français, la durée est introduite par une \emph{préposition} (« pendant », « depuis », « cela fait… que ») et sa place est libre. En chinois, la durée est un \emph{complément de verbe} : sa place est fixe, après le verbe.

\begin{center}
\begin{tabularx}{\linewidth}{|X|X|X|}
\hline
\rowcolor{popblueL} \textbf{Français} & \textbf{Chinois (caractères + pinyin)} & \textbf{Remarque} \\
\hline
Il a dormi huit heures. & \zh{他睡了八个小时。} Tā shuì le bā ge xiǎoshí. & durée après le verbe, avec \zh{了} \\
\hline
Il n'a pas dormi huit heures. & \zh{他没睡八个小时。} Tā méi shuì bā ge xiǎoshí. & \zh{没} remplace et chasse \zh{了} \\
\hline
J'étudie le chinois depuis deux ans. & \zh{我学中文学了两年了。} Wǒ xué Zhōngwén xué le liǎng nián le. & verbe répété devant l'objet \\
\hline
Cela fait deux ans que je n'ai pas étudié le chinois. & \zh{我有两年没学中文了。} Wǒ yǒu liǎng nián méi xué Zhōngwén le. & structure 有 + durée + 没 + V + 了 \\
\hline
Nous avons eu deux heures de cours. & \zh{我们上了两个小时的课。} Wǒmen shàng le liǎng ge xiǎoshí de kè. & verbe séparable, durée insérée \\
\hline
\end{tabularx}
\end{center}

\begin{popfigure}
\begin{tikzpicture}[node distance=0.6cm and 2.2cm, every node/.style={align=center, font=\small}]
\node[draw=popblue, fill=popblueL, rounded corners, text width=4.6cm, minimum height=1.4cm] (fr) {\textbf{Français}\\« Je dors \emph{depuis} 8 heures. »\\préposition + durée\\(place libre)};
\node[draw=popgreen, fill=popgreenL, rounded corners, text width=5.2cm, minimum height=1.4cm, right=of fr] (zh) {\textbf{Chinois}\\\zh{我睡了八个小时(了)。}\\Wǒ shuì le bā ge xiǎoshí (le).\\verbe + 了 + durée (place fixe)};
\draw[-{Stealth[length=3mm]}, very thick, poporange] (fr) -- node[above, font=\footnotesize]{traduire = déplacer la durée après le verbe} (zh);
\end{tikzpicture}
\legende{Du français au chinois : la durée passe du statut de prépositionnel (place libre) à celui de complément postverbal (place fixe).}
\end{popfigure}

\courssub{Méthode et entraînement}

\begin{methode}[Le réflexe en 3 étapes]
Pour construire une phrase avec complément de durée :
\begin{enumerate}
\item \textbf{Verbe seul ou avec objet ?} Identifiez si le verbe a un objet (\zh{学中文} : oui ; \zh{睡} : non).
\item \textbf{Répétez le verbe s'il y a un objet} : \zh{我学中文\textbf{学}了…} ; sinon, gardez le verbe simple : \zh{他睡了…}
\item \textbf{Placez 了 + durée après le verbe} : \zh{…学了两年。} À la négative : remplacez par \zh{没} + verbe + durée, sans \zh{了} : \zh{他没睡八个小时。}
\end{enumerate}
\end{methode}

\begin{exoresolu}[Corriger les phrases fautives]
Énoncé : les phrases suivantes, produites par des élèves de Première à Yaoundé, contiennent chacune une erreur. Corrigez-les et traduisez.\\
1. ✱\zh{三年他踢足球了。} 2. ✱\zh{我没睡了八个小时。} 3. ✱\zh{她学汉语了三年了。}
\tcblower
Solution détaillée :\\
1. La durée \zh{三年} est placée avant le verbe (calque du français « depuis trois ans »). Correction : \zh{他踢足球踢了三年了。} \emph{Tā tī zúqiú tī le sān nián le.} « Il joue au football depuis trois ans. » (verbe répété car il y a un objet \zh{足球}).\\
2. \zh{了} est gardé après \zh{没}. Correction : \zh{我没睡八个小时。} \emph{Wǒ méi shuì bā ge xiǎoshí.} « Je n'ai pas dormi huit heures. »\\
3. Le verbe n'est pas répété devant l'objet \zh{汉语}. Correction : \zh{她学汉语学了三年了。} \emph{Tā xué Hànyǔ xué le sān nián le.} « Elle étudie le chinois depuis trois ans. » (autre forme correcte : \zh{她学了三年的汉语。}).
\end{exoresolu}

\begin{exercice}[À vous de jouer]
Traduisez en chinois (caractères et pinyin) :\\
1. « Je n'ai pas mangé depuis dix heures. » (utiliser 有 + durée + 没 + V + 了)\\
2. « Nous avons eu quatre heures de cours de chinois. »\\
3. « Cela fait cinq jours qu'il n'est pas allé au marché de Douala. »\\
4. Corrigez : ✱\zh{两天我没有上课了。}
\end{exercice}

\begin{corrige}[Exercice « À vous de jouer »]
1. \zh{我有十个小时没吃饭了。} \emph{Wǒ yǒu shí ge xiǎoshí méi chīfàn le.}\\
2. \zh{我们上了四个小时的中文课。} \emph{Wǒmen shàng le sì ge xiǎoshí de Zhōngwén kè.} (verbe séparable \zh{上课}, durée insérée) ; forme simple acceptée : \zh{我们上课上了四个小时。}\\
3. \zh{他有五天没去杜阿拉的市场了。} \emph{Tā yǒu wǔ tiān méi qù Dù'ālā de shìchǎng le.}\\
4. La durée ne précède jamais le verbe ; de plus \zh{了} final est ici le \zh{了} de changement de situation, acceptable : \zh{我有两天没上课了。} \emph{Wǒ yǒu liǎng tiān méi shàngkè le.} « Cela fait deux jours que je n'ai pas eu cours. »
\end{corrige}

\begin{aretenir}[L'essentiel de la section]
\begin{itemize}
\item Négation : \zh{没(有)} + verbe + durée, \textbf{sans} \zh{了} : \zh{他没睡八个小时。}
\item « Cela fait X que je n'ai pas… » : \zh{有 + durée + 没 + V + 了} : \zh{我有两年没学中文了。}
\item Verbes séparables : la durée peut s'insérer (\zh{睡了八个小时的觉}), forme à reconnaître seulement.
\item Jamais de durée avant le verbe ; verbe répété devant un objet ; jamais \zh{了} après \zh{没}.
\end{itemize}
\end{aretenir}

\coursec{Observation : le redoublement des adjectifs}

Dans la section précédente, nous avons vu comment nier un complément de durée et quels pièges attendent les francophones (\zh{他没来过中国。} et non \zh{*他没来过中国了。}). Nous quittons maintenant le verbe pour observer un phénomène qui concerne les \cle{adjectifs} : en chinois, on peut les \cle{redoubler} pour rendre la description plus vivante ou pour exprimer la manière. Observons d'abord, réglons ensuite.

\courssub{Lecture guidée : quatre phrases à observer}

Lisons attentivement les quatre phrases suivantes. Soulignez mentalement ce qui se répète.

\begin{textebox}[Quatre phrases du quotidien]
她的眼睛大大的。\\
\emph{Tā de yǎnjing dàdà de.}\\
« Elle a de grands yeux (tout grands). »\\[4pt]
你要好好地学习。\\
\emph{Nǐ yào hǎohǎo de xuéxí.}\\
« Tu dois étudier bien, sérieusement. »\\[4pt]
教室里干干净净的。\\
\emph{Jiàoshì lǐ gāngānjìngjìng de.}\\
« La salle de classe est toute propre, impeccable. »\\[4pt]
他高高兴兴地回家了。\\
\emph{Tā gāogāoxìngxìng de huí jiā le.}\\
« Il est rentré à la maison tout joyeux. »
\end{textebox}

\textbf{Questions d'observation}
\begin{enumerate}
\item Dans la phrase 1, quel mot est répété ? Combien de syllabes a l'adjectif de départ \zh{大} ?
\item Dans les phrases 3 et 4, quels sont les adjectifs de départ ? Combien de syllabes ont-ils ? Comment sont-ils redoublés ?
\item Quel petit mot suit le redoublement dans les phrases 1, 2 et 3 ? Quel mot le suit dans la phrase 4 ?
\item Le redoublement décrit-il un nom (épithète ou attribut) ou la manière de faire une action ?
\end{enumerate}

\textbf{Constats attendus}
\begin{itemize}
\item \zh{大} (une syllabe) devient \zh{大大} : forme \cle{AA}.
\item \zh{干净} (deux syllabes) devient \zh{干干净净} ; \zh{高兴} devient \zh{高高兴兴} : forme \cle{AABB} (chaque syllabe est redoublée séparément, jamais ABAB).
\item Après le redoublement, on trouve \zh{的} quand on décrit un nom ou un état, et \zh{地} quand on décrit la manière d'accomplir une action.
\end{itemize}

\courssub{La règle : formes AA et AABB}

\begin{definition}[Le redoublement des adjectifs]
Le \cle{redoublement} d'un adjectif consiste à répéter ses syllabes pour \textbf{renforcer} ou \textbf{vivifier} la description, ou pour exprimer la \textbf{manière}.\\
\textbf{Adjectif monosyllabique} : A $\rightarrow$ \textbf{AA} (的) : \zh{大} $\rightarrow$ \zh{大大的}.\\
\textbf{Adjectif bisyllabique} : AB $\rightarrow$ \textbf{AABB} (的/地) : \zh{干净} $\rightarrow$ \zh{干干净净的} ; \zh{高兴} $\rightarrow$ \zh{高高兴兴地}.\\
Le redoublement s'emploie dans les descriptions vivantes, les récits, et devant un verbe pour indiquer la manière. Il est incompatible avec \zh{很} et avec les comparatifs.
\end{definition}

\begin{propriete}[Effet de sens du redoublement]
Le redoublement produit trois effets principaux :
\begin{enumerate}
\item \textbf{Insistance affectueuse} (attribut ou épithète) : \zh{她的眼睛大大的。} insiste sur la grandeur des yeux, avec une nuance attendrie, là où \zh{她的眼睛很大。} est une constatation neutre.
\item \textbf{Description vivante} : \zh{教室里干干净净的。} peint la salle « impeccable de partout » ; le redoublement a une valeur d'intensité et de totalité.
\item \textbf{Manière soignée} (avec \zh{地} devant un verbe) : \zh{你要好好地学习。} signifie « étudie comme il faut, sérieusement » ; \zh{他高高兴兴地回家了。} décrit la manière joyeuse de rentrer.
\end{enumerate}
\end{propriete}

\begin{attention}[Deux interdits absolus]
\begin{itemize}
\item \textbf{Jamais de} \zh{很} \textbf{devant un adjectif redoublé} : le redoublement porte déjà l'intensité. On dit \zh{她的眼睛大大的。}, jamais \zh{*她的眼睛很大大的。}
\item \textbf{Jamais de forme ABAB} pour les adjectifs : \zh{*干净干净} est une erreur. On redouble syllabe par syllabe : \zh{干干净净}. (La forme ABAB existe, mais pour les \emph{verbes} : \zh{研究研究} « examiner un peu ».)
\end{itemize}
\end{attention}

\courssub{Schéma de formation}

\begin{popfigure}
\begin{tikzpicture}[
  node distance=0.9cm and 1.6cm,
  base/.style={draw=popblue, fill=popblueL, rounded corners=2pt, minimum height=0.8cm, align=center, font=\small},
  res/.style={draw=poporange, fill=poporangeL, rounded corners=2pt, minimum height=0.8cm, align=center, font=\small},
  fl/.style={-{Stealth[length=2.5mm]}, thick, popdark}
]
\node[base, align=center] (a){A\\ 大 \emph{dà}};
\node[res, right=of a, align=center] (aa){AA + 的\\ 大大的 \emph{dàdà de}};
\node[base, below=of a, align=center] (ab){AB\\ 干净 \emph{gānjìng}};
\node[res, right=of ab, align=center] (aabb){AABB + 的/地\\ 干干净净的 \emph{gāngānjìngjìng de}};
\draw[fl] (a) -- (aa) node[midway, above, font=\footnotesize\itshape, popmuted]{monosyllabe};
\draw[fl] (ab) -- (aabb) node[midway, above, font=\footnotesize\itshape, popmuted]{bisyllabe};
\end{tikzpicture}
\legende{Formation du redoublement : A devient AA, AB devient AABB.}
\end{popfigure}

\courssub{Deux emplois : épithète avec 的, manière avec 地}

\begin{propriete}[Redoublement + 的 : épithète ou attribut]
Devant un nom, le redoublement sert d'\cle{épithète} : \zh{AA/AABB + 的 + Nom}.\\
\zh{大大的眼睛} \emph{dàdà de yǎnjing} « de grands yeux (tout grands) »\\
\zh{红红的脸} \emph{hónghóng de liǎn} « un visage tout rouge »\\
\zh{干干净净的教室} \emph{gāngānjìngjìng de jiàoshì} « une salle impeccable »\\
En fin de phrase, \zh{的} marque l'attribut descriptif : \zh{教室里干干净净的。}
\end{propriete}

\begin{propriete}[Redoublement + 地 : complément circonstanciel de manière]
Devant un verbe, le redoublement décrit la \cle{manière} : \zh{AA/AABB + 地 + Verbe}.\\
\zh{认认真真地写字} \emph{rènrènzhēnzhēn de xiě zì} « écrire avec beaucoup de soin »\\
\zh{慢慢地走} \emph{mànmàn de zǒu} « marcher lentement »\\
\zh{高高兴兴地回家} \emph{gāogāoxìngxìng de huí jiā} « rentrer tout joyeux »
\end{propriete}

\begin{attention}[的 ou 地 ? Le piège des mots-outils]
Les deux se prononcent \emph{de}, mais ne s'écrivent pas pareil :
\begin{itemize}
\item devant un \textbf{nom} $\rightarrow$ \zh{的} : \zh{大大的眼睛} ;
\item devant un \textbf{verbe} $\rightarrow$ \zh{地} : \zh{好好地学习}.
\end{itemize}
Erreur typique du francophone : écrire \zh{*他高高兴兴的回家了} avec \zh{的} devant le verbe. Demandez-vous toujours : « le mot qui suit est-il un nom ou un verbe ? »
\end{attention}

\begin{popfigure}
\begin{tikzpicture}[
  grow=down,
  level distance=1.2cm,
  level 1/.style={sibling distance=6.5cm},
  level 2/.style={sibling distance=3.2cm},
  every node/.style={draw=poppurple, fill=poppurpleL, rounded corners=2pt, align=center, font=\small, minimum height=0.7cm},
  edge from parent/.style={draw=popdark, thick, -{Stealth[length=2mm]}}
]
\node[align=center]{Redoublement\\ AA / AABB}
  child { node[align=center]{+ 的 : épithète / attribut\\ (devant un nom)}
    child { node[align=center]{大大的眼睛\\ \emph{de grands yeux}} }
    child { node[align=center]{红红的脸\\ \emph{un visage tout rouge}} }
  }
  child { node[align=center]{+ 地 : manière\\ (devant un verbe)}
    child { node[align=center]{认认真真地写字\\ \emph{écrire avec soin}} }
    child { node[align=center]{慢慢地走\\ \emph{marcher lentement}} }
  };
\end{tikzpicture}
\legende{Arbre de choix : 的 devant un nom, 地 devant un verbe.}
\end{popfigure}

\courssub{Tableau récapitulatif et comparaison avec le français}

\begin{center}
\begin{tabularx}{\linewidth}{|X|X|X|X|}
\hline
\rowcolor{popblueL} Base & Forme redoublée & Exemple + pinyin & Traduction \\
\hline
大 \emph{dà} (grand) & AA : 大大的 & 大大的眼睛 \emph{dàdà de yǎnjing} & de grands yeux \\
\hline
红 \emph{hóng} (rouge) & AA : 红红的 & 红红的脸 \emph{hónghóng de liǎn} & un visage tout rouge \\
\hline
慢 \emph{màn} (lent) & AA : 慢慢地 & 慢慢地走 \emph{mànmàn de zǒu} & marcher lentement \\
\hline
好 \emph{hǎo} (bien) & AA : 好好地 & 好好地学习 \emph{hǎohǎo de xuéxí} & étudier sérieusement \\
\hline
干净 \emph{gānjìng} (propre) & AABB : 干干净净 & 干干净净的教室 \emph{gāngānjìngjìng de jiàoshì} & une salle impeccable \\
\hline
高兴 \emph{gāoxìng} (joyeux) & AABB : 高高兴兴 & 高高兴兴地回家 \emph{gāogāoxìngxìng de huí jiā} & rentrer tout joyeux \\
\hline
认真 \emph{rènzhēn} (sérieux) & AABB : 认认真真 & 认认真真地写字 \emph{rènrènzhēnzhēn de xiě zì} & écrire très soigneusement \\
\hline
\end{tabularx}
\end{center}

\begin{aretenir}[Comparaison avec le français]
Le français n'a pas de redoublement grammatical de l'adjectif ; il utilise d'autres moyens que le chinois rend par le redoublement :
\begin{itemize}
\item l'adverbe d'intensité : « de \emph{tout} grands yeux » = \zh{大大的眼睛} ;
\item l'adverbe de manière : « marcher \emph{lentement} » = \zh{慢慢地走} ;
\item le ton affectueux ou enfantin : « tout propre, tout joyeux » = \zh{干干净净的、高高兴兴地}.
\end{itemize}
Attention : ne traduisez jamais le redoublement par une répétition en français (« grand grand ») ; cherchez l'équivalent naturel (« tout grand », « bien grand », « impeccable »).
\end{aretenir}

\begin{exemplebox}[Exemples en contexte camerounais]
\zh{雅温得的街道干干净净的。} \emph{Yǎwēndé de jiēdào gāngānjìngjìng de.} « Les rues de Yaoundé sont toutes propres. »\\[3pt]
\zh{孩子们高高兴兴地去上学。} \emph{Háizi men gāogāoxìngxìng de qù shàngxué.} « Les enfants vont à l'école tout joyeux. »\\[3pt]
\zh{请你慢慢地说法语。} \emph{Qǐng nǐ mànmàn de shuō Fǎyǔ.} « Parle français lentement, s'il te plaît. »\\[3pt]
\zh{妈妈做了香香的恩多莱。} \emph{Māma zuò le xiāngxiāng de ndolé.} « Maman a préparé un ndolé tout parfumé. »
\end{exemplebox}

\begin{savaistu}[Un ton vivant et affectueux]
Le redoublement adjectival est extrêmement fréquent dans les récits, les comptines et les descriptions destinées aux enfants : il donne un ton vivant, tendre et imagé. Dire à un enfant \zh{小手白白的} « les petites mains toutes blanches » sonne bien plus doux que \zh{手很白}. Les grands écrivains chinois l'utilisent aussi pour peindre une scène en quelques coups de pinceau. Au Cameroun, on retrouve un procédé voisin dans les contes en langues locales, où la répétition crée le rythme et l'émotion : les deux cultures aiment faire « danser » les mots.
\end{savaistu}

\courssub{Exercice résolu et piège final}

\begin{exoresolu}[Transformez avec le redoublement]
Énoncé — Réécrivez chaque phrase en redoublant l'adjectif entre parenthèses, avec \zh{的} ou \zh{地} selon le cas.\\
1. 她的眼睛很大。（大）\quad 2. 他认真地学习。（认真）\quad 3. 教室很干净。（干净）
\tcblower
Solution détaillée —\\
1. L'adjectif \zh{大} est monosyllabique $\rightarrow$ AA ; il décrit un état en fin de phrase $\rightarrow$ \zh{的}. On supprime \zh{很}, interdit devant un redoublement : \zh{她的眼睛大大的。} \emph{Tā de yǎnjing dàdà de.} « Elle a de grands yeux. »\\
2. \zh{认真} est bisyllabique $\rightarrow$ AABB ; il précède le verbe \zh{学习} $\rightarrow$ \zh{地} : \zh{他认认真真地学习。} \emph{Tā rènrènzhēnzhēn de xuéxí.} « Il étudie très sérieusement. »\\
3. \zh{干净} $\rightarrow$ \zh{干干净净} ; attribut en fin de phrase $\rightarrow$ \zh{的}, sans \zh{很} : \zh{教室里干干净净的。} \emph{Jiàoshì lǐ gāngānjìngjìng de.} « La salle est impeccable. »
\end{exoresolu}

\begin{attention}[Les trois erreurs à bannir]
\begin{enumerate}
\item Garder \zh{很} : \zh{*很大大的} $\rightarrow$ \zh{大大的}.
\item Redoubler en bloc : \zh{*高兴高兴地} $\rightarrow$ \zh{高高兴兴地} (syllabe par syllabe).
\item Confondre \zh{的} et \zh{地} : devant un nom \zh{的}, devant un verbe \zh{地}.
\end{enumerate}
\end{attention}

\begin{exercice}[À vous de jouer]
Traduisez en chinois avec un redoublement : 1. « La petite fille a un visage tout rouge. » 2. « Écoutez attentivement le professeur. » (écouter = \zh{听}) 3. « Le marché de Douala est très animé. » (animé = \zh{热闹} \emph{rènao}). Puis décrivez en deux phrases votre salle de classe en utilisant au moins un redoublement.
\end{exercice}

\begin{corrige}[Exercice : Redoublement des adjectifs]
1. \zh{小女孩的脸红红的。} \emph{Xiǎo nǚhái de liǎn hónghóng de.} (Adj. monosyllabique \zh{红} $\rightarrow$ AA + \zh{的} attribut).\\[4pt]
2. \zh{请认认真真地听老师的话。} \emph{Qǐng rènrènzhēnzhēn de tīng lǎoshī de huà.} (Adj. bisyllabique \zh{认真} $\rightarrow$ AABB + \zh{地} devant verbe \zh{听}).\\[4pt]
3. \zh{杜阿拉的市场热热闹闹的。} \emph{Dùālā de shìchǎng rènàorènao de.} (Adj. bisyllabique \zh{热闹} $\rightarrow$ AABB + \zh{的} attribut ; on supprime \zh{很}).\\[4pt]
Production personnelle (exemple) : \zh{我们的教室亮亮堂堂的。} \emph{Wǒmen de jiàoshì liàngliàngtángtáng de.} « Notre salle de classe est toute lumineuse. » / \zh{同学们高高兴兴地读课文。} \emph{Tóngxué men gāogāoxìngxìng de dú kèwén.} « Les camarades lisent le texte tout joyeux. »
\end{corrige}

Nous savons désormais former et employer le redoublement des adjectifs. La prochaine section réunira complément de durée et redoublement dans un entraînement intégré : phrases à transformer et mini-texte à rédiger.

\coursec{Schémas, emplois et limites du redoublement}

Dans la section précédente, nous avons \emph{observé} des phrases comme \zh{她高高兴兴地去上学。} et \zh{妈妈的手白白的。} Nous avons remarqué que les adjectifs \zh{高兴} et \zh{白} y apparaissent sous une forme \cle{redoublée}, et que ce redoublement donne à la phrase un ton vivant, affectueux ou insistant. Il est temps maintenant de dégager les \textbf{schémas exacts} de ce redoublement, ses \textbf{emplois} précis et surtout ses \textbf{limites}, car le redoublement des adjectifs obéit à des règles strictes que le francophone doit connaître pour ne pas tomber dans les pièges de la traduction mot à mot.

\courssub{Les deux schémas de redoublement : AA et AABB}

\begin{propriete}[Schéma AA : adjectifs monosyllabiques]
Un adjectif d'\textbf{une seule syllabe} se redouble en répétant le caractère : \textbf{A $\rightarrow$ AA}. Le second caractère prend généralement le \textbf{ton léger (neutre)} à l'oral (pas de marque de ton en pinyin).
\begin{itemize}
\item \zh{大} (dà, grand) $\rightarrow$ \zh{大大} (dàda)
\item \zh{慢} (màn, lent) $\rightarrow$ \zh{慢慢} (mànman)
\item \zh{红} (hóng, rouge) $\rightarrow$ \zh{红红} (hónghong)
\item \zh{白} (bái, blanc) $\rightarrow$ \zh{白白} (báibai)
\end{itemize}
\end{propriete}

\begin{propriete}[Schéma AABB : adjectifs dissyllabiques]
Un adjectif de \textbf{deux syllabes} AB se redouble syllabe par syllabe : \textbf{AB $\rightarrow$ AABB}. Attention : on ne dit pas \zh{高兴高兴} (ABAB), forme réservée aux verbes. Les 1ʳᵉ et 3ᵉ syllabes gardent souvent leur ton d'origine (ou passent au 1ᵉʳ ton), les 2ᵉ et 4ᵉ syllabes prennent souvent le ton léger.
\begin{itemize}
\item \zh{高兴} (gāoxìng, content) $\rightarrow$ \zh{高高兴兴} (gāogāoxìngxing)
\item \zh{干净} (gānjìng, propre) $\rightarrow$ \zh{干干净净} (gāngānjìngjing)
\item \zh{漂亮} (piàoliang, joli) $\rightarrow$ \zh{漂漂亮亮} (piàopiàoliangliang)
\item \zh{清楚} (qīngchu, clair) $\rightarrow$ \zh{清清楚楚} (qīngqīngchǔchu)
\end{itemize}
\end{propriete}

\begin{attention}[Ne pas confondre ABAB et AABB]
Le schéma \textbf{ABAB} (\zh{看看} kànkan, \zh{休息休息} xiūxi xiūxi) est le redoublement des \emph{verbes} : il exprime une action brève ou un essai. Le schéma \textbf{AABB} est le redoublement des \emph{adjectifs} : il décrit un état de façon vivante. Dire \zh{高兴高兴} pour « être un peu content » est une erreur fréquente des francophones qui calquent le redoublement verbal sur les adjectifs.
\end{attention}

\begin{popfigure}
\begin{tikzpicture}[
box/.style={draw=popblue, fill=popblueL, rounded corners=2pt, align=center, font=\small, inner sep=4pt},
res/.style={draw=poporange, fill=poporangeL, rounded corners=2pt, align=center, font=\small, inner sep=4pt},
use/.style={draw=popgreen, fill=popgreenL, rounded corners=2pt, align=center, font=\small, inner sep=4pt},
fleche/.style={-{Stealth[length=2mm]}, thick, popdark}]
\node[box, align=center] (a1) at (0,2.2){A\\ \zh{大} dà};
\node[box, align=center] (a2) at (0,1.1){A\\ \zh{慢} màn};
\node[box, align=center] (b1) at (0,-0.3){AB\\ \zh{高兴} gāoxìng};
\node[box, align=center] (b2) at (0,-1.4){AB\\ \zh{干净} gānjìng};
\node[res, align=center] (r1) at (3.4,1.65){AA\\ \zh{大大} dàda / \zh{慢慢} mànman};
\node[res, align=center] (r2) at (3.4,-0.85){AABB\\ \zh{高高兴兴} gāogāoxìngxing / \zh{干干净净} gāngānjìngjing};
\node[use, align=center] (u1) at (7.2,1.65){+ \zh{的} : description\\ \zh{红红的苹果} hónghong de píngguǒ};
\node[use, align=center] (u2) at (7.2,0.4){+ \zh{地} : manière\\ \zh{慢慢地说} mànman de shuō};
\draw[fleche] (a1) -- (r1);
\draw[fleche] (a2) -- (r1);
\draw[fleche] (b1) -- (r2);
\draw[fleche] (b2) -- (r2);
\draw[fleche] (r1) -- (u1);
\draw[fleche] (r1) -- (u2);
\draw[fleche] (r2) -- (u1);
\draw[fleche] (r2) -- (u2);
\end{tikzpicture}
\legende{Schéma de transformation : adjectif $\rightarrow$ AA/AABB $\rightarrow$ + 的 (description) ou + 地 (manière)}
\end{popfigure}

\courssub{Emploi 1 : épithète ou attribut descriptif, avec 的}

Le premier grand emploi du redoublement est la \textbf{description} : l'adjectif redoublé, suivi de \zh{的}, qualifie un nom (épithète) ou sert d'attribut après le sujet. Il peint l'état avec vivacité, souvent avec une nuance affective ou admirative, comme le français « tout + adjectif » (\emph{tout rouge}, \emph{toute propre}).

\begin{exemplebox}[Description avec AA/AABB + 的]
\zh{红红的苹果很好吃。} \emph{Hónghong de píngguǒ hěn hǎochī.} « Les pommes bien rouges sont bonnes. »\\
\zh{她的手白白的。} \emph{Tā de shǒu báibai de.} « Ses mains sont toutes blanches. »\\
\zh{教室里干干净净的。} \emph{Jiàoshì lǐ gāngānjìngjing de.} « La salle de classe est toute propre. »\\
\zh{妹妹穿着一件漂漂亮亮的衣服。} \emph{Mèimei chuānzhe yí jiàn piàopiàoliangliang de yīfu.} « La petite sœur porte un très joli vêtement. »\\
\zh{杜阿拉的市场热热闹闹的。} \emph{Dù'ālā de shìchǎng rèrènàonao de.} « Le marché de Douala est tout animé. »
\end{exemplebox}

Remarquez la construction attributive fréquente \textbf{Sujet + AA/AABB + 的} (\zh{她的手白白的。}) : le \zh{的} final est presque obligatoire ici ; sans lui, la phrase sonne incomplète.

\courssub{Emploi 2 : complément de manière, avec 地}

Le deuxième emploi est \textbf{adverbial} : l'adjectif redoublé, suivi de \zh{地}, se place \emph{avant le verbe} pour dire \emph{comment} l'action est faite. C'est l'équivalent du français « bien + adverbe » ou d'un adverbe en \emph{-ment} appuyé.

\begin{exemplebox}[Manière avec AA/AABB + 地 + verbe]
\zh{请你慢慢地说。} \emph{Qǐng nǐ mànman de shuō.} « Parle lentement, s'il te plaît. »\\
\zh{我们要好好地学习汉语。} \emph{Wǒmen yào hǎohao de xuéxí Hànyǔ.} « Nous devons étudier le chinois sérieusement (bien). »\\
\zh{她高高兴兴地去上学。} \emph{Tā gāogāoxìngxing de qù shàngxué.} « Elle va à l'école joyeusement. »\\
\zh{老师把汉字清清楚楚地写在黑板上。} \emph{Lǎoshī bǎ Hànzì qīngqīngchǔchu de xiě zài hēibǎn shang.} « Le professeur écrit les caractères très clairement au tableau. »
\end{exemplebox}

\begin{attention}[的 ou 地 ? Un choix d'écriture décisif]
À l'oral, \zh{的} et \zh{地} se prononcent tous les deux \emph{de} (ton léger) : impossible de les distinguer en écoutant. À l'\textbf{écrit}, la confusion est une faute grave :
\begin{itemize}
\item \zh{的} introduit la \textbf{description} d'un nom : \zh{红红的苹果} « la pomme toute rouge » ;
\item \zh{地} introduit la \textbf{manière} d'un verbe : \zh{慢慢地说} « parler lentement ».
\end{itemize}
Test simple : si le mot qui suit est un \emph{nom}, écrivez \zh{的} ; si c'est un \emph{verbe}, écrivez \zh{地}.
\end{attention}

\begin{methode}[Choisir la bonne construction]
\begin{enumerate}
\item Voulez-vous \textbf{décrire} une personne ou une chose avec vivacité ? $\rightarrow$ AA/AABB + \zh{的} + nom, ou Sujet + AA/AABB + \zh{的}. Exemple : \zh{干干净净的教室} « une salle toute propre ».
\item Voulez-vous dire \textbf{comment} se fait une action ? $\rightarrow$ AA/AABB + \zh{地} + verbe. Exemple : \zh{好好地休息} « bien se reposer ».
\item Voulez-vous simplement constater un degré (« très », « bien ») sans effet descriptif ? $\rightarrow$ Pas de redoublement : utilisez \zh{很} + adjectif : \zh{很大} « très grand ».
\end{enumerate}
\end{methode}

\courssub{Les limites du redoublement : ce qu'il est interdit de faire}

\begin{propriete}[Incompatibilité avec 很]
Un adjectif redoublé \textbf{ne se combine jamais avec} \zh{很} (hěn, très) ni avec les autres adverbes de degré (\zh{太}, \zh{非常}…). Le redoublement contient déjà l'idée d'intensité descriptive. On dit \zh{很大} \emph{ou} \zh{大大的}, jamais \zh{很大大的}.
\end{propriete}

\begin{attention}[Le piège du mot à mot : « très grand »]
Le francophone veut traduire « très grand » et produit \zh{大大的} ou pire \zh{很大大的}. Or :
\begin{itemize}
\item « Cette maison est \emph{très grande} » (constat neutre) $\rightarrow$ \zh{这个房子很大。} (\emph{Zhège fángzi hěn dà.})
\item « Une \emph{grande, grande} maison » (description vivante, affective) $\rightarrow$ \zh{一个大大的房子} (\emph{yí ge dàda de fángzi})
\end{itemize}
Le redoublement n'est pas un simple renforcement : c'est un \emph{effet de style} descriptif. Dans un exposé neutre ou une phrase informative, préférez \zh{很} + adjectif.
\end{attention}

\begin{propriete}[Adjectifs qui ne se redoublent pas]
Certains adjectifs monosyllabiques fréquents \textbf{ne se redoublent pas} dans cet usage descriptif : \zh{对} (duì, correct), \zh{错} (cuò, faux), \zh{多} (duō, nombreux), \zh{少} (shǎo, peu nombreux). On dit \zh{很对} « tout à fait correct », \zh{很多} « beaucoup », jamais \zh{对对的} ni \zh{多多的} (sauf \zh{多多} dans des expressions figées comme \zh{多多关照}, formule de politesse). De même, les adjectifs « absolus » (\zh{男} masculin, \zh{女} féminin) et les adjectifs dissyllabiques formels ou à connotation négative (\zh{伟大} wěidà, \zh{难看} nánkàn) ne se redoublent pas.
\end{propriete}

\courssub{Comparaison avec le français et tableau récapitulatif}

Le français ne redouble pas ses adjectifs ; il utilise d'autres procédés que le chinois rend par le redoublement :

\begin{center}
\begin{tabularx}{\linewidth}{|X|X|X|}
\hline
\rowcolor{popblueL} Français & Chinois (redoublement) & Pinyin \\
\hline
tout rouge (épithète) & \zh{红红的苹果} & hónghong de píngguǒ \\
\hline
toute propre (attribut) & \zh{干干净净的} & gāngānjìngjing de \\
\hline
bien, sérieusement & \zh{好好地学习} & hǎohao de xuéxí \\
\hline
lentement, doucement & \zh{慢慢地说} & mànman de shuō \\
\hline
joyeusement & \zh{高高兴兴地去} & gāogāoxìngxing de qù \\
\hline
très grand (constat neutre) & \zh{很大} (sans redoublement) & hěn dà \\
\hline
\end{tabularx}
\end{center}

\begin{center}
\begin{tabularx}{\linewidth}{|X|X|X|X|}
\hline
\rowcolor{popblueL} Caractères & Pinyin & Nature & Traduction \\
\hline
\zh{好好} & hǎohao & adj. redoublé (AA) & bien, sérieusement \\
\hline
\zh{慢慢} & mànman & adj. redoublé (AA) & lentement, doucement \\
\hline
\zh{高高兴兴} & gāogāoxìngxing & adj. redoublé (AABB) & joyeux, joyeusement \\
\hline
\zh{干干净净} & gāngānjìngjing & adj. redoublé (AABB) & tout propre, impeccable \\
\hline
\zh{漂漂亮亮} & piàopiàoliangliang & adj. redoublé (AABB) & très joli, élégant \\
\hline
\zh{清清楚楚} & qīngqīngchǔchu & adj. redoublé (AABB) & très clair, clairement \\
\hline
\zh{热闹} & rènao & adjectif (base) & animé, vivant \\
\hline
\end{tabularx}
\end{center}

\begin{popfigure}
\begin{tikzpicture}[mindmap, grow cyclic, every node/.style={concept, align=center, font=\small},
root concept/.append style={concept color=popblue, font=\small\bfseries},
level 1/.append style={level distance=3.4cm, sibling angle=90},
level 1 concept/.append style={concept color=poporangeL}]
\node[root concept, align=center]{Redoublement\\ des adjectifs}
child { node[align=center]{\zh{好好} hǎohao\\ « bien,\\ sérieusement »} }
child { node[align=center]{\zh{慢慢} mànman\\ « lentement,\\ doucement »} }
child { node[align=center]{\zh{高高兴兴}\\ gāogāoxìngxing\\ « joyeusement »} }
child { node[align=center]{\zh{干干净净}\\ gāngānjìngjing\\ « tout propre »} };
\end{tikzpicture}
\legende{Carte mentale : quatre redoublements fréquents et leurs traductions}
\end{popfigure}

\begin{exoresolu}[Transformer et corriger]
\textbf{Énoncé.} 1) Transformez \zh{她很高兴地去学校。} en utilisant le redoublement AABB. 2) La phrase \zh{这个苹果很大大的。} est-elle correcte ? Justifiez et corrigez. 3) Traduisez : « Écoute bien (sérieusement), s'il te plaît. »
\tcblower
\textbf{Solution détaillée.}
\begin{enumerate}
\item L'adjectif dissyllabique \zh{高兴} se redouble en AABB : \zh{高高兴兴}. Devant le verbe \zh{去}, il faut la marque de manière \zh{地} : \zh{她高高兴兴地去学校。} \emph{Tā gāogāoxìngxing de qù xuéxiào.} « Elle va à l'école joyeusement. » (On supprime \zh{很}, incompatible avec le redoublement.)
\item La phrase est \textbf{incorrecte} : un adjectif redoublé ne se combine pas avec \zh{很}. Deux corrections possibles selon le sens : constat neutre $\rightarrow$ \zh{这个苹果很大。} « Cette pomme est très grande » ; description vivante $\rightarrow$ \zh{这个苹果大大的。} « Cette pomme est bien grosse. »
\item « Écouter sérieusement » = manière $\rightarrow$ AA + \zh{地} + verbe : \zh{请你好好地听。} \emph{Qǐng nǐ hǎohao de tīng.}
\end{enumerate}
\end{exoresolu}

\begin{exercice}[À vous de jouer]
1) Mettez l'adjectif entre parenthèses à la forme redoublée correcte : \zh{妹妹（漂亮）地穿上了新衣服。} 2) Choisissez \zh{的} ou \zh{地} : \zh{他慢慢（ ）走回家。} / \zh{红红（ ）花儿真好看。} 3) Traduisez en chinois : « La cour de notre lycée de Yaoundé est toute propre. » 4) Corrigez la faute : \zh{我们要很好好地学习。}
\end{exercice}

\begin{corrige}[Exercice « À vous de jouer »]
1) \zh{漂漂亮亮} (AABB) : \zh{妹妹漂漂亮亮地穿上了新衣服。} 2) \zh{他慢慢地走回家。} (manière, devant le verbe \zh{走}) ; \zh{红红的花儿真好看。} (description, devant le nom \zh{花儿}). 3) \zh{我们雅温得中学的校园干干净净的。} \emph{Wǒmen Yǎwēndé zhōngxué de xiàoyuán gāngānjìngjing de.} 4) Supprimer \zh{很} : \zh{我们要好好地学习。} « Nous devons étudier sérieusement. »
\end{corrige}

\begin{aretenir}[L'essentiel de la section]
\begin{itemize}
\item Deux schémas : \textbf{AA} pour les adjectifs d'une syllabe (\zh{大} $\rightarrow$ \zh{大大}), \textbf{AABB} pour ceux de deux syllabes (\zh{高兴} $\rightarrow$ \zh{高高兴兴}).
\item Deux emplois : \textbf{description} avec \zh{的} (\zh{红红的苹果}) ; \textbf{manière} avec \zh{地} + verbe (\zh{慢慢地说}).
\item Trois interdits : jamais de \zh{很} devant un adjectif redoublé ; pas de redoublement pour \zh{对}, \zh{错}, \zh{多}, \zh{少} ; ne pas utiliser le redoublement pour traduire mécaniquement « très ».
\end{itemize}
\end{aretenir}

\coursec{Entraînement intégré : phrases et mini-texte}

Après avoir analysé séparément le complément de durée et le redoublement des adjectifs, nous allons maintenant les mettre en pratique de façon combinée. C'est dans l'enchaînement naturel du discours — raconter sa journée, décrire une activité, exprimer un ressenti — que ces deux structures révèlent toute leur puissance expressive. Le complément de durée ancre le fait dans le temps (combien de temps ?), tandis que le redoublement colore l'action ou l'état d'une nuance affective, visuelle ou rythmique (comment ?).

\courssub{Exercices de transformation : ajout du complément de durée}

\begin{methode}[Transformer une phrase simple en phrase avec durée]
Pour ajouter une durée à une phrase existante, suivez ces étapes :
\begin{enumerate}
    \item Identifiez le \cle{verbe} principal (ex. : \zh{学习} xuéxí, \zh{睡觉} shuìjiào).
    \item Identifiez l'\cle{objet} éventuel (ex. : \zh{中文} Zhōngwén, \zh{功课} gōngkè).
    \item Déterminez la \cle{durée} à insérer (ex. : \zh{三个小时} sān ge xiǎoshí).
    \item Appliquez la règle de l'\cle{ordre des mots} :
    \begin{itemize}
        \item Verbe intransitif ou verbe sans objet exprimé : Sujet + Verbe + \textbf{了} + Durée + \textbf{led} (aspect + modalité).
        \item Verbe transitif avec objet (ou verbe à structure VO \emph{séparable} type \zh{睡觉}, \zh{打球}) : Sujet + Verbe + Objet + \textbf{Verbe} + \textbf{led} + Durée + \textbf{led} (répétition du verbe \emph{seul} obligatoire).
    \end{itemize}
    \item Vérifiez la place de \textbf{led} (après le verbe répété, jamais après l'objet seul).
\end{enumerate}
\end{methode}

\begin{exemplebox}[Transformations types]
\begin{itemize}
    \item \textbf{Base} : \zh{他学习。} (Tā xuéxí.) — Il étudie.
    \item \textbf{+ Durée (3h)} : \zh{他学习学了三个小时了。} (Tā xuéxí xué le sān ge xiǎoshí le.) — Il étudie depuis trois heures / Il a étudié pendant trois heures.
    \item \textbf{Base} : \zh{我学中文。} (Wǒ xué Zhōngwén.) — J'apprends le chinois.
    \item \textbf{+ Durée (2 ans)} : \zh{我学中文学了两年了。} (Wǒ xué Zhōngwén xué le liǎng nián le.) — J'apprends le chinois depuis deux ans.
    \item \textbf{Base (verbe VO séparable)} : \zh{她睡觉。} (Tā shuìjiào.) — Elle dort.
    \item \textbf{+ Durée (8h)} : \zh{她睡觉睡了八个小时了。} (Tā shuìjiào shuì le bā ge xiǎoshí le.) — Elle dort depuis 8 heures / Elle a dormi 8 heures. \emph{(Répétition du morphème verbal \zh{睡})}
\end{itemize}
\end{exemplebox}

\begin{exoresolu}[Transformation guidée]
\textbf{Énoncé :} Transformez les phrases suivantes en y insérant la durée indiquée entre parenthèses. Respectez la structure du complément de durée (répétition du verbe si objet).

1. \zh{她睡觉。} (八个小时) \\
2. \zh{我们打球。} (两个小时) — \emph{Objet : 球} \\
3. \zh{我写汉字。} (半个小时) — \emph{Objet : 汉字} \\
4. \zh{小李看书。} (一年) — \emph{Objet : 书} \\

\tcblower
\textbf{Solution détaillée :}

1. \zh{她睡觉睡了八个小时了。} (Tā shuìjiào shuì le bā ge xiǎoshí le.) \\
\emph{Verbe séparable \zh{睡觉} (VO) : on répète le premier morphème \zh{睡}. Traduction : Elle a dormi pendant 8 heures / Elle dort depuis 8 heures.}

2. \zh{我们打球打了两个小时了。} (Wǒmen dǎqiú dǎ le liǎng ge xiǎoshí le.) \\
\emph{Structure VO \zh{打球} : répétition de \zh{打}. Traduction : Nous jouons au ballon depuis 2 heures.}

3. \zh{我写汉字写了半个小时了。} (Wǒ xiě Hànzì xiě le bàn ge xiǎoshí le.) \\
\emph{Verbe transitif standard : répétition de \zh{写}. Traduction : J'écris des caractères chinois depuis une demi-heure.}

4. \zh{小李看书看了一年了。} (Xiǎo Lǐ kànshū kàn le yī nián le.) \\
\emph{Structure VO \zh{看书} : répétition de \zh{看}. Traduction : Petit Li lit des livres / étudie depuis un an.}
\end{exoresolu}

\begin{attention}[Piège n°1 : L'ordre des mots « Durée avant verbe »]
En français, on dit « \emph{J'étudie \textbf{depuis deux ans} le chinois} » ou « \emph{Il a dormi \textbf{pendant huit heures}} » (durée en fin de phrase ou devant le verbe avec « depuis »).
En chinois, la durée \textbf{ne se place jamais avant le verbe principal}.
\begin{itemize}
    \item \textbf{Faux} : \zh{*我两年学中文学了。} / \zh{*我学了两年中文。}
    \item \textbf{Vrai} : \zh{我学中文学了两年了。}
\end{itemize}
La durée suit \emph{obligatoirement} le verbe (éventuellement répété).
\end{attention}

\courssub{Exercices de mise en ordre : remettre les mots dans l'ordre}

\begin{exercice}[Mise en ordre — Complément de durée]
Remettez les mots dans l'ordre correct pour former une phrase grammaticale. Attention à la double occurrence de \zh{了} et à la répétition du verbe \emph{uniquement si le verbe a un objet}.

1. \zh{了 / 我 / 两年 / 学 / 中文 / 学 / 了} \\
2. \zh{三个 / 小时 / 了 / 他 / 写作业 / 写 / 了} \\
3. \zh{看 / 电视 / 了 / 一晚上 / 妹妹 / 看 / 了} \\
4. \zh{了 / 等 / 你 / 等 / 了 / 半天 / 我} \\
\end{exercice}

\begin{corrige}[Exercice : Mise en ordre]
1. \zh{我学中文学了两年了。} (Wǒ xué Zhōngwén xué le liǎng nián le.) \\
2. \zh{他写作业写了三个小时了。} (Tā xiě zuòyè xiě le sān ge xiǎoshí le.) \\
\emph{Note : \zh{写作业} (VO) $\rightarrow$ répétition de \zh{写}. \zh{工作} est intransitif, pas de répétition : \zh{他工作了三个小时了。}} \\
3. \zh{妹妹看电视看了一晚上了。} (Mèimei kàn diànshì kàn le yī wǎnshàng le.) \\
4. \zh{我等你等了半天了。} (Wǒ děng nǐ děng le bàn tiān le.) \\
\emph{Note : Pour « 等你 », le verbe \zh{等} prend l'objet \zh{你}, donc on répète \zh{等}.}
\end{corrige}

\courssub{Exercices de redoublement des adjectifs : formes AA et AABB}

\begin{propriete}[Règles de formation du redoublement]
\begin{itemize}
    \item \textbf{Adjectifs monosyllabiques} (ex. \zh{高} gāo, \zh{慢} màn, \zh{好} hǎo, \zh{早} zǎo) $\rightarrow$ forme \textbf{AA} (\zh{高高} gāogāo, \zh{慢慢} mànman, \zh{好好} hǎohao, \zh{早早} zǎozǎo).
    \item \textbf{Adjectifs bisyllabiques} (ex. \zh{干净} gānjìng, \zh{高兴} gāoxìng, \zh{认真} rènzhēn, \zh{舒服} shūfu) $\rightarrow$ forme \textbf{AABB} (\zh{干干净净} gāngānjìngjìng, \zh{高高兴兴} gāogāoxìngxìng, \zh{认认真真} rènrènzhēnzhēn, \zh{舒舒服服} shūshūfúfú).
    \item \textbf{Fonction adverbiale (manière)} : forme redoublée + \zh{地} (de) devant le verbe. Ex: \zh{慢慢走} / \zh{慢慢地走}, \zh{高高兴兴地来}.
    \item \textbf{Fonction complémentaire (résultat/degré)} : verbe + \zh{得} (de) + forme redoublée. Ex: \zh{写得工工整整}, \zh{学得扎扎实实}.
    \item \textbf{Interdiction formelle} : Ne mettez \textbf{JAMAIS} \zh{很} (hěn) devant un adjectif redoublé. Le redoublement exprime déjà l'intensité, la délimitation ou la manière.
\end{itemize}
\end{propriete}

\begin{exemplebox}[Redoublement : transformation et insertion]
\begin{center}
\begin{tabularx}{\linewidth}{|X|X|X|X|}
\hline
\rowcolor{popblueL} Adjectif de base & Type & Forme redoublée & Phrase exemple (Caractères + Pinyin + Traduction) \\
\hline
\zh{高} (gāo) & Mono. & \zh{高高} (gāogāo) & \zh{树长得高高的。} (Shù zhǎng de gāogāo de.) — Les arbres poussent bien hauts. \\
\hline
\zh{慢} (màn) & Mono. & \zh{慢慢} (mànman) & \zh{请慢慢吃。} (Qǐng mànman chī.) — Mangez lentement. \\
\hline
\zh{好} (hǎo) & Mono. & \zh{好好} (hǎohao) & \zh{你要好好学习。} (Nǐ yào hǎohao xuéxí.) — Tu dois bien étudier. \\
\hline
\zh{高兴} (gāoxìng) & Bisyll. & \zh{高高兴兴} (gāogāoxìngxìng) & \zh{他高高兴兴地回家了。} (Tā gāogāoxìngxìng de huí jiā le.) — Il est rentré tout joyeux. \\
\hline
\zh{干净} (gānjìng) & Bisyll. & \zh{干干净净} (gāngānjìngjìng) & \zh{请把屋子收拾得干干净净。} (Qǐng bǎ wūzi shōushi de gāngānjìngjìng.) — Range pour qu'elle soit bien propre. \\
\hline
\zh{认真} (rènzhēn) & Bisyll. & \zh{认认真真} (rènrènzhēnzhēn) & \zh{我们要认认真真地复习。} (Wǒmen yào rènrènzhēnzhēn de fùxí.) — Nous devons réviser sérieusement. \\
\hline
\end{tabularx}
\end{center}
\end{exemplebox}

\begin{exercice}[Redoublement et insertion]
Transformez les adjectifs entre parenthèses en forme redoublée (AA ou AABB) et insérez-les correctement dans la phrase (ajoutez \zh{地} si nécessaire pour la manière, \zh{得} pour le résultat).

1. Il a nettoyé la table \_\_\_\_\_. (\zh{干净} gānjìng) \\
2. Elle est rentrée \_\_\_\_\_. (\zh{高兴} gāoxìng) \\
3. Conduis \_\_\_\_\_ ! (\zh{慢} màn) \\
4. Nous devons réviser \_\_\_\_\_. (\zh{好} hǎo) \\
5. Les caractères sont écrits \_\_\_\_\_. (\zh{工整} gōngzhěng) \\
\end{exercice}

\begin{corrige}[Exercice : Redoublement]
1. \zh{他把桌子擦得干干净净。} (Tā bǎ zhuōzi cā de gāngānjìngjìng.) — \emph{Structure \textbf{把} + V + \zh{得} + Adj redoublé AABB (résultat).} \\
2. \zh{她高高兴兴地回家了。} (Tā gāogāoxìngxìng de huí jiā le.) — \emph{Adverbe de manière (AABB) + \zh{地}.} \\
3. \zh{请慢慢开车！} (Qǐng mànman kāichē!) — \emph{Forme AA \zh{慢慢} + \zh{地} souvent omis pour monosyllabes en injonction.} \\
4. \zh{我们要好好复习。} (Wǒmen yào hǎohao fùxí.) — \emph{\zh{好好} (hǎohao) = sérieusement, bien. \zh{地} omis.} \\
5. \zh{汉字写得工工整整。} (Hànzì xiě de gōnggōngzhěngzhěng.) — \emph{V + \zh{得} + AABB (résultat/degré).}
\end{corrige}

\begin{attention}[Piège n°2 : \zh{很} + Redoublement = ERREUR MAJEURE]
C'est l'erreur n°1 des francophones qui veulent « renforcer » l'adjectif.
\begin{itemize}
    \item \textbf{Faux} : \zh{*他很高高兴兴地来了。} / \zh{*请很慢慢走。} / \zh{*写得很工工整整。}
    \item \textbf{Vrai} : \zh{他高高兴兴地来了。} / \zh{请慢慢走。} / \zh{写得工工整整。}
\end{itemize}
\emph{Explication :} Le redoublement \textbf{est déjà} une marque d'intensité, de délimitation ou de manière. \zh{很} (hěn) est un marqueur de degré pour l'adjectif \textbf{simple} (ex. \zh{很高} hěn gāo). Ils sont en distribution complémentaire.
\end{attention}

\courssub{Correction de phrases fautives : erreurs types francophones}

\begin{exoresolu}[Correction d'erreurs — Diagnostic et réparation]
\textbf{Énoncé :} Chaque phrase ci-dessous contient une ou deux erreurs typiques (ordre durée/verbe, double \zh{了}, \zh{没} + \zh{了}, \zh{很} + redoublement, place de \zh{地}/\zh{得}). Corrigez-les et expliquez l'erreur.

1. \zh{我学了两年了中文。} \\
2. \zh{他没有来了两天。} \\
3. \zh{她很高高兴兴地跳舞。} \\
4. \zh{我们吃饭吃两个小时了。} (contexte : action terminée/accomplie) \\
5. \zh{请好好地写字。} (contexte : injonction standard) \\

\tcblower
\textbf{Solutions et explications :}

1. \textbf{Correction :} \zh{我学中文学了两年了。} \\
\textbf{Erreur :} Place de la durée (\zh{两年}) avant l'objet (\zh{中文}) et absence de répétition du verbe \zh{学}. Règle : S + V + O + V + \zh{了} + Durée + \zh{led}.

2. \textbf{Correction :} \zh{他两天没来了。} (Il n'est pas venu depuis deux jours / Ça fait deux jours qu'il ne vient plus). \\
\textbf{Erreur :} \textbf{\zh{没} (méi) et \zh{了} (led aspectuel) sont incompatibles.} \zh{没} nie le passé/accompli ; \zh{了} (aspectuel) marque l'accompli. On ne dit jamais \zh{*没有...了}. Le \zh{了} final (modalité) est possible avec \zh{没} pour marquer l'état actuel (\zh{他两天没来了}), mais pas \zh{了} aspectuel après le verbe.

3. \textbf{Correction :} \zh{她高高兴兴地跳舞。} \\
\textbf{Erreur :} \textbf{\zh{很} + Redoublement.} Supprimer \zh{很}.

4. \textbf{Correction :} \zh{我们吃饭吃了两个小时了。} \\
\textbf{Erreur :} Manque du \zh{led} \textbf{aspectuel} (le premier \zh{led}, après le verbe répété). Sans ce \zh{led}, la phrase sonne comme une habitude présente (« Nous mangeons pendant 2h »), or le contexte et le \zh{led} final (modalité) exigent l'aspect accompli.

5. \textbf{Correction :} \zh{请好好写字。} (ou \zh{请把字写好。}) \\
\textbf{Erreur :} \zh{地} (de) souvent omis après \zh{好好} (hǎohao) en style injonctif/parlé courant, mais sa présence n'est pas \emph{fausse} (juste plus formel/lent). L'erreur principale attendue ici est souvent l'ajout de \zh{很} (ex: \zh{*请很好好写字}) ou la confusion \zh{得}/\zh{地}. \emph{Note : \zh{好好写字} est la norme orale.}
\end{exoresolu}

\courssub{Mini-texte guidé : « Ma journée » (我的一天)}

Voici un texte modèle combinant les deux structures. Observez comment la durée structure la chronologie factuelle, tandis que le redoublement exprime l'attitude du locuteur.

\begin{textebox}[Ma journée — 我的一天]
\zh{我今天早上七点起床，七点半出门。} \\
\zh{我上课上了六个小时了，中午只休息了四十分钟。} \\
\zh{下午放学后，我好好地复习了功课，} \\
\zh{把生词抄得工工整整，把课文读得熟熟练练。} \\
\zh{晚上六点半，我高高兴兴地回家吃饭。} \\
\zh{饭后，我慢悠悠地散了二十分钟步，} \\
\zh{然后开开心心地睡觉了。}
\end{textebox}

\begin{center}
\begin{tabularx}{\linewidth}{|X|X|X|}
\hline
\rowcolor{popblueL} \textbf{Phrase chinoise + Pinyin} & \textbf{Structure ciblée} & \textbf{Traduction} \\
\hline
\zh{我上课上了六个小时了} (Wǒ shàngkè shàng le liù ge xiǎoshí le) & Complément de durée (VO + V + 了 + Durée + 了) & J'ai eu cours pendant 6 heures (jusqu'à maintenant). \\
\hline
\zh{中午只休息了四十分钟} (Zhōngwǔ zhǐ xiūxi le sìshí fēnzhōng) & Complément de durée simple (V + 了 + Durée) & À midi, je me suis reposé seulement 40 min. \\
\hline
\zh{我好好地复习了功课} (Wǒ hǎohao de fùxí le gōngkè) & Redoublement \zh{好好} + \zh{地} (manière) & J'ai bien / sérieusement révisé mes leçons. \\
\hline
\zh{把生词抄得工工整整} (Bǎ shēngcí chāo de gōnggōngzhěngzhěng) & \textbf{把} + V + \zh{得} + Redoublement AABB (résultat) & J'ai copié les mots nouveaux bien proprement. \\
\hline
\zh{把课文读得熟熟练练} (Bǎ kèwén dú de shúshúliànliàn) & \textbf{把} + V + \zh{得} + Redoublement AABB (degré) & J'ai lu le texte jusqu'à le connaître sur le bout des doigts. \\
\hline
\zh{我高高兴兴地回家} (Wǒ gāogāoxìngxìng de huí jiā) & Redoublement AABB + \zh{地} (manière/état d'esprit) & Je suis rentré chez moi tout joyeux. \\
\hline
\zh{我慢悠悠地散了二十分钟步} (Wǒ mányōuyōu de sàn le èrshí fēnzhōng bù) & Redoublement \zh{慢悠悠} (manière) + Durée (VO séparé) & Je me suis promené tranquillement pendant 20 min. \\
\hline
\zh{开开心心地睡觉了} (Kāikāixīnxīn de shuìjiào le) & Redoublement AABB + \zh{地} + \zh{led} (modalité) & Je me suis endormi tout content. \\
\hline
\end{tabularx}
\end{center}

\begin{methode}[Rédiger un texte combinant Durée + Redoublement]
Pour écrire un court récit (journal, mail, exposé oral) riche et naturel :
\begin{enumerate}
    \item \textbf{Structurez le temps} : Utilisez le \textbf{complément de durée} pour les actions factuelles, mesurables, les « blocs » de temps (cours, sommeil, trajet, travail). \emph{Ex: 我睡了八个小时觉。}
    \item \textbf{Colorez l'expérience} : Utilisez le \textbf{redoublement} pour décrire \textbf{comment} vous avez vécu l'action (votre état d'esprit, la qualité du résultat, la manière). \emph{Ex: 我舒舒服服地睡了八个小时觉。}
    \item \textbf{Alternez} : Ne mettez pas du redoublement partout (effet « surchargé »). Réservez-le aux moments clés : transition importante, réussite, plaisir, effort conscient.
    \item \textbf{Vérifiez la syntaxe} :
    \begin{itemize}
        \item Durée : Après le verbe (répété si objet). Double \zh{led} si action terminée/état actuel.
        \item Redoublement : Devant le verbe + \zh{地} (manière) ; Après \zh{得} (résultat/degré) ; \textbf{Sans \zh{很}}.
    \end{itemize}
\end{enumerate}
\end{methode}

\courssub{Production guidée : Votre journée}

\begin{exercice}[Expression écrite — Mon emploi du temps idéal]
Rédigez un court paragraphe (80-100 caractères chinois) intitulé \zh{我的理想一天} (Mon jour idéal). Contraintes obligatoires :
\begin{enumerate}
    \item Utilisez \textbf{au moins 3 compléments de durée} différents (heures, minutes, mois, ans...).
    \item Utilisez \textbf{au moins 3 adjectifs redoublés} (formes AA ou AABB) dont au moins un après \zh{得} (structure \textbf{把} ou V+\zh{得}) et un devant le verbe + \zh{地}.
    \item Respectez l'ordre des mots (Durée après V ; Redoublement sans \zh{很}).
    \item Utilisez le vocabulaire du module (Environnement/Santé : \zh{运动} yùndòng, \zh{散步} sànbù, \zh{呼吸} hūxī, \zh{新鲜} xīnxiān, \zh{放松} fàngsōng).
\end{enumerate}
\end{exercice}

\begin{corrige}[Exercice : Mon jour idéal — Proposition de corrigé type]
\zh{我的理想一天} \\
\zh{早上六点，我早早地起床，在公园慢悠悠地跑了四十分钟步，} \\
\zh{呼吸着清清爽爽的空气。} \\
\zh{早饭后，我认认真真地学习了两个小时汉语，} \\
\zh{把课文背得滚瓜烂熟。} \\
\zh{下午高高兴兴地和朋友打了两个小时球，} \\
\zh{晚上十点舒舒服服地睡觉，睡了整整八个小时。} \\
\zh{真棒！}

\emph{Analyse des structures utilisées :}
\begin{itemize}
    \item \textbf{Durée} : \zh{跑了四十分钟步}, \zh{学习了两个小时}, \zh{打了两个小时球}, \zh{睡了整整八个小时}.
    \item \textbf{Redoublement + \zh{地}} : \zh{早早地}, \zh{慢悠悠地}, \zh{认认真真地}, \zh{高高兴兴地}, \zh{舒舒服服地}.
    \item \textbf{Redoublement + \zh{得}} (Résultat) : \zh{背得滚瓜烂熟} (idiome AABB signifiant « par cœur / sur le bout des doigts »).
    \item \textbf{Redoublement attributif} : \zh{清清爽爽的空气} (air frais et pur).
\end{itemize}
\end{corrige}

\courssub{Préparation à l'activité d'intégration (Oral / Écrit)}

\begin{experience}[Activité de classe : « Enquête : Nos rythmes de vie »]
\textbf{Objectif :} Réinvestir le complément de durée (questions \zh{多久} / \zh{多长时间}) et le redoublement (description de l'état d'esprit) à l'oral.

\textbf{Déroulement (15 min) :}
\begin{enumerate}
    \item \textbf{Préparation individuelle (3 min) :} Chaque élève remplit un mini-tableau pour 3 activités de sa semaine type (ex. : Sport, Devoirs, Téléphone, Sommeil, Trajet).
    \begin{center}
    \begin{tabularx}{\linewidth}{|X|X|X|X|}
    \hline
    \rowcolor{popblueL} Activité & Durée (结构) & Redoublement (Comment ?) & Phrase complète \\
    \hline
    睡觉 & 八个小时 & 舒舒服服地 & 我每天舒舒服服地睡八个小时觉。 \\
    \hline
    学习 & 两个小时 & 认认真真地 & 我每天认认真真地学习两个小时。 \\
    \hline
    看手机 & 三个小时 & 开开心心地 & 我每天开开心心地看三个小时手机。 \\
    \hline
    \end{tabularx}
    \end{center}
    \item \textbf{Interview en binôme (7 min) :} L'élève A interroge B : \zh{「你每天睡觉睡多久？感觉怎么样？」} B répond avec la structure complète. A note les réponses de B.
    \item \textbf{Restitution orale (5 min) :} 3-4 binômes présentent le rythme de leur partenaire à la classe : \zh{「小王每天认认真真地学习两个小时，但是开开心心地看三个小时手机！」}
\end{enumerate}
\textbf{Critères d'évaluation :} Exactitude de la structure durée (ordre, double \zh{led}), exactitude du redoublement (forme, pas de \zh{很}, place \zh{地}), fluidité, prononciation des tons (notamment 3e ton \zh{好好} hǎohao $\rightarrow$ háo hǎo sandhi).
\end{experience}

\begin{savaistu}[Le Bac Blanc : Les pièges grammaticaux « Durée \& Redoublement »]
Au Bac blanc (et au Baccalauréat), les items de grammaire sur ces deux points sont très codifiés. Trois types d'exercices reviennent chaque année :
\begin{enumerate}
    \item \textbf{Transformation (Durée) :} « Transformez la phrase en y ajoutant la durée. » \\
    \emph{Ex: 我学习法语。 (两年) $\rightarrow$ 我学法文学了两年了。} \textbf{Piège :} Oublier de répéter le verbe si objet présent.
    \item \textbf{Correction d'erreur (Ordre / 了) :} « Corrigez la phrase. » \\
    \emph{Ex: *他没吃了饭。 $\rightarrow$ 他没吃饭。 / 他没吃饭吃。} \textbf{Piège :} \zh{没} + \zh{led} (aspectuel).
    \item \textbf{Choix multiple / Complétion (Redoublement) :} « Choisissez la forme correcte. » \\
    \emph{Ex: 请把屋子收拾 \_\_\_\_\_ (干净). Options: A. 干净 B. 很干净 C. 干干净净 D. 干净干净.} \textbf{Bonne réponse : C.} \textbf{Piège :} Option B (\zh{很} + adj simple ne va pas après \zh{得} pour un résultat), Option D (forme inexistante pour bisyllabique).
\end{enumerate}
\textbf{Astuce de révision :} Faites une fiche « 3 colonnes » : \textbf{Structure} | \textbf{Exemple canonique} | \textbf{Erreur interdite}. Révisez-la la veille de l'épreuve.
\end{savaistu}

% --- FIGURE TIKZ : Frise « Ma journée » ---
\begin{popfigure}
\begin{tikzpicture}[
    node distance=1.2cm and 0.5cm,
    event/.style={rectangle, rounded corners=4pt, draw=popblue, thick, fill=popblueL, text width=2.8cm, align=center, font=\small, minimum height=1.8cm},
    arrow/.style={-Stealth, thick, popdark},
    time/.style={font=\bfseries\small, popdark, anchor=east},
    duration/.style={font=\tiny, poporange, anchor=west},
    labelstruct/.style={font=\tiny, anchor=west, align=left, text width=2.5cm}
]
% Nœuds temps (colonne gauche)
\node[time] (t1) at (0,0) {07:00};
\node[time] (t2) at (0,-2.5) {08:00};
\node[time] (t3) at (0,-5.0) {14:00};
\node[time] (t4) at (0,-7.5) {18:30};
\node[time] (t5) at (0,-10.0) {22:00};

% Nœuds événements (colonne droite)
\node[event, right=of t1, align=center] (e1){\textbf{起床 qǐchuáng}\\\emph{Se lever}\\\zh{早早地} (tôt/gaîment)};
\node[event, right=of t2, align=center] (e2){\textbf{上课 shàngkè}\\\emph{Avoir cours}\\\zh{上了六个小时了} (6h)};
\node[event, right=of t3, align=center] (e3){\textbf{复习 fùxí}\\\emph{Réviser}\\\zh{好好地 / 工工整整} (bien/proprement)};
\node[event, right=of t4, align=center] (e4){\textbf{回家 huí jiā}\\\emph{Rentrer}\\\zh{高高兴兴地} (joyeusement)\\\zh{散步二十分钟} (prom. 20min)};
\node[event, right=of t5, align=center] (e5){\textbf{睡觉 shuìjiào}\\\emph{Dormir}\\\zh{舒舒服服地}\\\zh{睡了八个小时觉} (8h)};

% Flèches
\draw[arrow] (e1) -- (e2);
\draw[arrow] (e2) -- (e3);
\draw[arrow] (e3) -- (e4);
\draw[arrow] (e4) -- (e5);

% Étiquettes structure
\node[labelstruct, poppurple, right=0.2cm of e1] {Redoublement +\zh{地}};
\node[labelstruct, popgreen, right=0.2cm of e2] {Durée (VO+V+了+了)};
\node[labelstruct, poppurple, right=0.2cm of e3] {Redoublement +\zh{地} / +\zh{得}};
\node[labelstruct, poppurple, right=0.2cm of e4] {Redoublement +\zh{地}};
\node[labelstruct, popgreen, right=0.2cm of e4, yshift=-0.8cm] {Durée (VO séparé)};
\node[labelstruct, poppurple, right=0.2cm of e5] {Redoublement +\zh{地}};
\node[labelstruct, popgreen, right=0.2cm of e5, yshift=-0.8cm] {Durée (VO+V+了+了)};

% Titre
\node[font=\bfseries\normalsize, popdark, above=0.3cm of e1] {Frise chronologique : Ma journée (我的一天)};
\end{tikzpicture}
\legende{Visualisation de l'alternance Durée (vert) / Redoublement (violet) dans un récit journalier.}
\end{popfigure}

% --- TABLEAU RÉCAPITULATIF FINAL ---
\begin{center}
\begin{tabularx}{\linewidth}{|X|X|X|X|}
\hline
\rowcolor{popblueL} \textbf{Structure} & \textbf{Schéma canonique} & \textbf{Exemple clé (Car. + Pin. + Trad.)} & \textbf{Points de vigilance (Francophones)} \\
\hline
\textbf{Complément de durée} &
S + V + \zh{led} + Durée + \zh{led} \\
(Si Objet / VO : S + V + O + V + \zh{led} + Durée + \zh{led}) &
\zh{我学中文学了两年了。} \\
(Wǒ xué Zhōngwén xué le liǎng nián le.) \\
J'apprends le chinois depuis 2 ans. &
\begin{itemize}
\item Durée \textbf{après} le verbe (jamais avant).
\item Répétition du V \textbf{obligatoire} si Objet (verbe transitif ou VO séparable).
\item Double \zh{led} (aspect + modalité) pour action terminée/état actuel.
\item \textbf{Interdit :} \zh{没} + \zh{led} (aspectuel).
\end{itemize} \\
\hline
\textbf{Redoublement Adj} &
\textbf{Mono.} (ex. \zh{高、慢、好}) $\rightarrow$ \textbf{AA} (\zh{高高、慢慢、好好}) \\
\textbf{Bisyll.} (ex. \zh{干净、高兴}) $\rightarrow$ \textbf{AABB} (\zh{干干净净、高高兴兴}) \\
+ \zh{地} (devant V) / + \zh{得} (après V/\textbf{把}) &
\zh{他高高兴兴地回家。} \\
(Tā gāogāoxìngxìng de huí jiā.) \\
Il rentre tout joyeux. &
\begin{itemize}
\item \textbf{JAMAIS \zh{很} + Redoublement.}
\item \zh{地} devant V (manière), \zh{得} après V (résultat/degré).
\item Attributif : Redoublement + \zh{的} + Nom (\zh{清清爽爽的空气}).
\item Tons : sandhi fréquent (3e ton + 3e ton $\rightarrow$ 2e + 3e : \zh{好好} háo hǎo).
\end{itemize} \\
\hline
\end{tabularx}
\end{center}

\begin{aretenir}[Synthèse de la leçon ZHA16]
\begin{itemize}
    \item Le \textbf{complément de durée} quantifie l'action dans le temps. Sa syntaxe rigoureuse (V + O $\rightarrow$ V + \zh{led} + Durée + \zh{led}) est la source n°1 d'erreurs de syntaxe (ordre, double \zh{led}, répétition V).
    \item Le \textbf{redoublement des adjectifs} qualifie la manière ou le résultat. Il remplace l'adverbe d'intensité \zh{很} et structure l'expression affective (AABB pour bisyllabiques) ou visuelle/résultative (AA pour monosyllabiques, AABB pour bisyllabiques).
    \item \textbf{Combinaison gagnante :} \zh{Subject + [Redoublement+地] + Verb + [Durée] + [Redoublement+得]}. Ex : \zh{我认认真真地学习了两个小时，学得扎扎实实。}
    \item Pour le Bac : Maîtrisez la transformation (ajout durée), la correction (ordre/\zh{没}/\zh{led}/\zh{很}) et la production courte (récit journaliers).
\end{itemize}
\end{aretenir}

\newpage

\coursec{Activité d'intégration}

\courssub{Objectifs de l'activité}

À l'issue de cette activité d'intégration, tu seras capable de :
\begin{itemize}
\item poser des questions sur la durée d'une action ou d'une habitude en utilisant correctement le \cle{complément de durée} ;
\item répondre à l'oral en plaçant le complément de durée au bon endroit dans la phrase ;
\item employer le \cle{redoublement des adjectifs} pour décrire la manière ou l'état d'une personne ;
\item rédiger un court portrait écrit (6 à 8 phrases) combinant au moins trois compléments de durée et deux adjectifs redoublés ;
\item repérer et corriger les erreurs d'ordre des mots dans les productions de tes camarades.
\end{itemize}

\courssub{Situation}

Dans le cadre de la semaine des langues du lycée de Nkol-Eton (Yaoundé), le club de chinois prépare un petit journal mural intitulé \zh{我们的同学} (\emph{Wǒmen de tóngxué}, « Nos camarades »). Chaque élève doit interviewer un camarade sur son quotidien : depuis combien de temps il apprend le chinois, combien d'heures il dort, combien de temps il fait du sport, depuis combien de temps il joue au football, etc. Les réponses serviront à rédiger un court portrait qui sera affiché au tableau du club.

Voici la fiche d'interview distribuée par l'enseignant :

\begin{textebox}[Fiche d'interview du club de chinois]
采访问题 \emph{Cǎifǎng wèntí} — Questions de l'interview\\
1. 你学中文学了多长时间？\emph{Nǐ xué Zhōngwén xué le duō cháng shíjiān ?} — Depuis combien de temps apprends-tu le chinois ?\\
2. 你每天睡觉睡几个小时？\emph{Nǐ měi tiān shuìjiào shuì jǐ ge xiǎoshí ?} — Combien d'heures dors-tu chaque jour ?\\
3. 你每天运动多长时间？\emph{Nǐ měi tiān yùndong duō cháng shíjiān ?} — Combien de temps fais-tu du sport chaque jour ?\\
4. 你踢足球踢了几年了？\emph{Nǐ tī zúqiú tī le jǐ nián le ?} — Depuis combien d'années joues-tu au football ?\\
5. 你学习认真吗？\emph{Nǐ xuéxí rènzhēn ma ?} — Étudies-tu sérieusement ?
\end{textebox}

\courssub{Consignes}

\begin{enumerate}
\item \textbf{Préparation (5 min).} Lis la fiche d'interview à voix basse. Souligne dans chaque question le verbe qui porte le complément de durée (structure : verbe + objet + verbe + durée). Vérifie le vocabulaire dans le tableau des ressources.
\item \textbf{Interview orale (10 min).} Par binômes, pose les cinq questions à ton camarade. Il répond par des phrases complètes. Note ses réponses en chinois. Inversez ensuite les rôles. Attention : la question 5 doit recevoir une réponse avec un adjectif redoublé (par exemple \zh{认认真真}).
\item \textbf{Vérification grammaticale (5 min).} Relis tes notes : chaque complément de durée est-il bien placé après le verbe (ou le verbe redoublé) ? As-tu au moins une phrase avec la structure « verbe + objet + verbe + 了 + durée » ? Corrige avec ton camarade.
\item \textbf{Rédaction du portrait (15 min).} Rédige un portrait de 6 à 8 phrases commençant par \zh{我的同学叫……} (\emph{Wǒ de tóngxué jiào...}, « Mon camarade s'appelle... »). Ton texte doit contenir \textbf{au moins trois compléments de durée} et \textbf{au moins deux adjectifs redoublés}. Souligne-les dans ta copie.
\item \textbf{Lecture et repérage collectif (10 min).} Lis ton portrait à voix haute. La classe lève la main à chaque complément de durée entendu et tape dans les mains à chaque adjectif redoublé. L'enseignant note au tableau les erreurs d'ordre des mots et la classe les corrige ensemble.
\end{enumerate}

\courssub{Ressources à mobiliser}

\begin{aretenir}[Rappel : le complément de durée — schémas de base]
\begin{itemize}
\item \textbf{Verbe sans objet (ou objet omis) :} Sujet + Verbe + \zh{了} + Durée.\\
\zh{他睡了八个小时。} \emph{Tā shuì le bā ge xiǎoshí.} « Il a dormi huit heures. »
\item \textbf{Verbe avec objet (structure VO) :} \textbf{Redoublement du verbe} — Sujet + Verbe + Objet + Verbe + \zh{了} + Durée.\\
\zh{我学中文学了一年了。} \emph{Wǒ xué Zhōngwén xué le yī nián le.} « J'apprends le chinois depuis un an (et j	continue). »
\item \textbf{Verbes séparables (VO lexicalisés : 睡觉, 运动, 跳舞, 见面…)} : deux possibilités équivalentes\\
\quad 1. \textbf{Séparation} : Verbe + Durée + Objet-nominal. \zh{睡了八个小时觉} (peu usité à l'oral).\\
\quad 2. \textbf{Redoublement (standard)} : Verbe-Objet + Verbe + Durée. \zh{睡觉睡了八个小时} / \zh{运动运动了一个小时}.
\item \textbf{Action continue jusqu'à présent :} ajouter \zh{了} final (aspectuel). \zh{他踢足球踢了三年了。} « Il joue au football depuis trois ans (et continue). »
\item \textbf{Négation :} \zh{没(有)} + Verbe + Durée (sans \zh{了}). \zh{我两天没运动了。} « Je n'ai pas fait de sport depuis deux jours. »
\end{itemize}
\end{aretenir}

\begin{aretenir}[Rappel : le redoublement des adjectifs — schémas et emplois]
\begin{itemize}
\item \textbf{Adjectif monosyllabique (AA) :} \zh{高} → \zh{高高的} (\emph{gāogāo de}, « tout grand ») ; \zh{壮} → \zh{壮壮的} (\emph{zhuàngzhuàng de}, « bien costaud »).\\
\quad Emploi attributif : AA + \zh{的} + Nom. \zh{高高的个子}.\\
\quad Emploi prédicatif : Sujet + AA + \zh{的}. \zh{他高高的。}
\item \textbf{Adjectif dissyllabique (AABB) :} \zh{认真} → \zh{认认真真} (\emph{rènrènzhēnzhēn}, « très sérieusement ») ; \zh{高兴} → \zh{高高兴兴} (\emph{gāogāoxìngxìng}, « joyeusement ») ; \zh{干净} → \zh{干干净净} (\emph{gāngānjìngjìng}, « bien propre »).\\
\quad Emploi adverbial : AABB + \zh{地} + Verbe. \zh{他认认真真地学习。} « Il étudie très sérieusement. »\\
\quad Emploi attributif : AABB + \zh{的} + Nom. \zh{高高兴兴的孩子} « un enfant joyeux ».
\item \textbf{Restrictions :} Les adjectifs déjà composés de deux caractères sans structure racine-affixe (ex. \zh{漂亮} piàoliang, \zh{年轻} niánqīng, \zh{可爱} kě'ài) ne se redoublent \textbf{pas} en AABB. Les adjectifs de couleur (\zh{红}, \zh{白}…) se redoublent en AA (\zh{红红的}) mais rarement en AABB.
\end{itemize}
\end{aretenir}

\begin{center}
\begin{tabularx}{\linewidth}{|X|X|X|X|}
\hline
\rowcolor{popblueL} Caractères & Pinyin & Nature & Traduction \\
\hline
采访 & cǎifǎng & verbe / nom & interviewer ; interview \\
\hline
多长时间 & duō cháng shíjiān & expression & combien de temps \\
\hline
小时 & xiǎoshí & nom & heure (durée) \\
\hline
睡觉 & shuìjiào & verbe (VO) & dormir \\
\hline
运动 & yùndong & verbe (VO) / nom & faire du sport ; sport \\
\hline
踢足球 & tī zúqiú & locution VO & jouer au football \\
\hline
认真 & rènzhēn & adjectif & sérieux, appliqué \\
\hline
每天 & měi tiān & compl. de temps & chaque jour \\
\hline
已经 & yǐjing & adverbe & déjà \\
\hline
身体 & shēntǐ & nom & corps, santé \\
\hline
个子 & gèzi & nom & taille, stature \\
\hline
\end{tabularx}
\end{center}

\begin{popfigure}
\begin{tikzpicture}[
    node distance=1.2cm and 1.5cm,
    box/.style={draw=popblue, thick, rounded corners, fill=popblueL, text width=3.2cm, align=center, font=\small},
    arrow/.style={-Stealth, thick, popdark}
]
\node[box] (sujet) {Sujet};
\node[box, right=of sujet, align=center] (temps){Compl. de temps\\ (\zh{每天}...)};
\node[box, right=of temps] (verbe1) {Verbe 1};
\node[box, right=of verbe1] (objet) {Objet};
\node[box, right=of objet] (verbe2) {Verbe 2 (redoublé)};
\node[box, right=of verbe2] (le) {\zh{了}};
\node[box, right=of le] (duree) {Durée};
\node[box, right=of duree] (le2) {\zh{了} (final)};

\draw[arrow] (sujet) -- (temps);
\draw[arrow] (temps) -- (verbe1);
\draw[arrow] (verbe1) -- (objet);
\draw[arrow] (objet) -- (verbe2);
\draw[arrow] (verbe2) -- (le);
\draw[arrow] (le) -- (duree);
\draw[arrow] (duree) -- (le2);

\node[draw=poporange, thick, rounded corners, fill=poporangeL, text width=5cm, align=left, font=\small, below=2cm of verbe1.south west, anchor=north west] (ex) {
\textbf{Exemple :} \zh{我 学 中文 学 了 一年 了。}\\
\emph{Wǒ xué Zhōngwén xué le yī nián le.}\\
S \quad V \quad O \quad V \quad \zh{了} \quad Dur \quad \zh{了}
};
\draw[arrow, poporange] (verbe1.south) -- ++(0,-0.5) -| (ex.north);
\end{tikzpicture}
\legende{Schéma de la structure du complément de durée avec objet (redoublement verbal)}
\end{popfigure}

\begin{propriete}[Comparaison Français / Chinois — Ordre des mots]
\begin{center}
\begin{tabularx}{\linewidth}{|X|X|X|}
\hline
\rowcolor{popblueL} Français (SVO + Compléments) & Chinois (Sujet + Temps + Verbe + Durée) & Piège à éviter \\
\hline
J'apprends le chinois \textbf{depuis un an}. & \zh{我 学 中文 学 了 一年 了。} (V + O + V + 了 + Durée) & Ne pas dire : \zh{我 学了 中文 一年} \\
\hline
Il dort \textbf{huit heures} par nuit. & \zh{他 睡觉 睡 八 个 小时。} (VO + V + Durée) & Ne pas dire : \zh{他 八 小时 睡觉} \\
\hline
Il joue au foot \textbf{depuis trois ans}. & \zh{他 踢 足球 踢 了 三年 了。} (V + O + V + 了 + Dur + 了) & Ne pas mettre la durée avant le verbe. \\
\hline
Il étudie \textbf{sérieusement}. & \zh{他 认认真真 地 学习。} (AABB + \zh{地} + V) & Ne pas oublier \zh{地} ; ne pas dire \zh{他 认真 学习} (trop neutre). \\
\hline
\end{tabularx}
\end{center}
\end{propriete}

\begin{attention}[Pièges fréquents des francophones — Check-list de révision]
\begin{itemize}
\item \textbf{Ordre durée / verbe :} En français, la durée est souvent en fin de phrase ou introduite par « depuis ». En chinois, la durée suit \textbf{immédiatement le verbe} (ou le verbe redoublé). \textbf{Interdit :} \zh{我一年学中文}, \zh{我八小时睡觉}.
\item \textbf{Objet + Durée = Redoublement obligatoire :} Si le verbe a un objet exprimé, on \textbf{doit} redoubler le verbe. \textbf{Faux :} \zh{我学了中文一年} ; \zh{我踢了足球三年}. \textbf{Juste :} \zh{我学中文学了一年}, \zh{我踢足球踢了三年}.
\item \textbf{Verbes séparables (睡觉, 运动, 吃饭, 见面…) :} On ne dit pas \zh{睡觉了八个小时}. Choix 1 (courant) : \zh{睡觉睡了八个小时} (redoublement). Choix 2 (objet omis) : \zh{睡了八个小时}.
\item \textbf{\zh{了} double :} Le premier \zh{了} (verbal, aspect accompli) colle au verbe redoublé ; le second \zh{了} (final, modal) marque la continuation ou le changement d'état. \zh{学了一年了} = « cela fait un an que j'apprends (et je continue) ».
\item \textbf{Négation :} \zh{没(有)} + Verbe + Durée. Pas de \zh{了}. \zh{我没睡八个小时} (je n'ai pas dormi 8h) vs \zh{我两天没睡觉了} (ça fait 2 jours que je n'ai pas dormi).
\item \textbf{Redoublement adjectif + \zh{地}/\zh{的} :} \zh{认认真真地} (adverbe), \zh{高高的} (attribut). \textbf{Jamais} \zh{认认真真学习} (manque \zh{地}). \textbf{Jamais} \zh{高高兴兴孩子} (manque \zh{的}).
\item \textbf{Adjectifs non redoublables :} \zh{漂亮} $\not\to$ \zh{漂漂亮亮}, \zh{年轻} $\not\to$ \zh{年年轻轻}. On utilise \zh{很漂亮}, \zh{非常年轻}.
\end{itemize}
\end{attention}

\begin{methode}[Écriture des caractères clés de la leçon]
\textbf{Ordre des traits (règles : haut→bas, gauche→droite, horizontal avant vertical, trait central avant les ailes, cadre avant contenu, trait de fermeture en dernier).}
\begin{itemize}
\item \zh{认} (\emph{rèn}, 14 traits) : Clé \zh{言} (parole, 7 traits) + \zh{忄} (cœur, 3 traits) à gauche. Écrire \zh{言} : point, horizontale, 3 horizontales, verticale, horizontale de fermeture ; puis \zh{忄} : point, point, verticale.
\item \zh{真} (\emph{zhēn}, 10 traits) : Clé \zh{目} (œil, 5 traits) en haut + \zh{十} + \zh{目} (simplifié) en bas. Attention : le bas s'écrit \zh{匕} + \zh{八} (ne pas confondre avec \zh{具}).
\item \zh{运} (\emph{yùn}, 7 traits) : Clé \zh{辶} (marche, 3 traits) en bas à gauche (s'écrit EN DERNIER : point, horizontale crochet, point long). Haut : \zh{云} (nuage, 4 traits) : horizontale, verticale, horizontale, horizontale crochet.
\item \zh{动} (\emph{dòng}, 6 traits) : Clé \zh{力} (force, 2 traits) à droite (s'écrit en dernier : horizontale crochet, point). Haut : \zh{云} (même haut que \zh{运}).
\item \zh{睡} (\emph{shuì}, 13 traits) : Clé \zh{目} (œil) en haut à gauche. Bas : \zh{垂} (pendre) : \zh{土} + \zh{攵} (frapper). \zh{攵} s'écrit : horizontale, verticale, point, horizontale crochet, point.
\item \zh{觉} (\emph{jué}, 14 traits) : Clé \zh{见} (voir, 4 traits) en bas à droite. Haut : \zh{学} (apprendre) sans le point final : \zh{宀} + \zh{子} + \zh{攵} (simplifié).
\end{itemize}
\end{methode}

\courssub{Proposition de production et corrigé commenté}

\begin{exoresolu}[Portrait modèle : « Mon camarade Étienne »]
Rédige un portrait de ton camarade (6 à 8 phrases) avec au moins trois compléments de durée et deux adjectifs redoublés.
\tcblower
\textbf{Production modèle :}

我的同学叫艾田，他是雅温得人。\emph{Wǒ de tóngxué jiào Àitián, tā shì Yǎwēndé rén.} « Mon camarade s'appelle Étienne, il est de Yaoundé. »

他学中文学了一年了。\emph{Tā xué Zhōngwén xué le yī nián le.} « Il apprend le chinois depuis un an (et continue). »

他每天都认认真真地学习。\emph{Tā měi tiān dōu rènrènzhēnzhēn de xuéxí.} « Chaque jour, il étudie très sérieusement. »

他每天晚上睡觉睡八个小时。\emph{Tā měi tiān wǎnshang shuìjiào shuì bā ge xiǎoshí.} « Chaque soir, il dort huit heures. »

他踢足球踢了三年了，踢得很好。\emph{Tā tī zúqiú tī le sān nián le, tī de hěn hǎo.} « Il joue au football depuis trois ans (et continue), et il joue très bien. »

他高高的，壮壮的，身体很健康。\emph{Tā gāogāo de, zhuàngzhuàng de, shēntǐ hěn jiànkāng.} « Il est tout grand, bien costaud, et en très bonne santé. »

周末，他常常高高兴兴地去运动场。\emph{Zhōumò, tā chángcháng gāogāoxìngxìng de qù yùndòngchǎng.} « Le week-end, il va souvent joyeusement au terrain de sport. »

我很喜欢他。\emph{Wǒ hěn xǐhuan tā.} « Je l'aime beaucoup. »

\textbf{Commentaire des choix de langue :}
\begin{itemize}
\item \textbf{Trois compléments de durée} : 1) \zh{学了一年了} (V+O+V+了+Dur+了 final, action continue) ; 2) \zh{睡八个小时} (verbe séparable, redoublement VO+V+Dur) ; 3) \zh{踢了三年了} (même structure que 1). Les trois schémas vus en leçon sont représentés.
\item \textbf{Trois adjectifs redoublés} (dépasse le minimum) : 1) \zh{认认真真地} (AABB + \zh{地}, emploi adverbial) ; 2) \zh{高高的，壮壮的} (AA + \zh{的}, emploi attributif/prédicatif) ; 3) \zh{高高兴兴地} (AABB + \zh{地}). Variété AA / AABB montrée.
\item \textbf{Ordre des mots} : Le complément de temps \zh{每天} / \zh{每天晚上} précède le verbe ; la durée suit toujours le verbe (ou V2). Aucune phrase ne place la durée avant le verbe (erreur typique : *\zh{我一年学中文}).
\item \textbf{Complément de degré} : \zh{踢得很好} (V + \zh{得} + Adv + Adj) révisé ici naturellement.
\item \textbf{Niveau HSK 1-3} : Vocabulaire de fréquence élevée, phrases courtes, aucun caractère hors programme.
\end{itemize}
\end{exoresolu}

\begin{experience}[Jeu de rôle : le journaliste du club]
Un élève joue le journaliste du journal mural et interviewe trois camarades devant la classe avec les questions de la fiche. Les auditeurs complètent une grille : nom, durée d'apprentissage du chinois, heures de sommeil, sport pratiqué et depuis combien de temps. Le journaliste résume ensuite à l'oral : \zh{玛丽学中文学了两年了。她每天睡觉睡七个小时。……} La classe vérifie la place de chaque complément de durée.
\end{experience}

\begin{exercice}[Entraînement individuel — Transformation et correction]
\begin{enumerate}
\item Transforme ces phrases en ajoutant le complément de durée indiqué (utilise le redoublement si nécessaire) :
      \begin{itemize}
      \item 我学法语。 (两年) $\rightarrow$ \dots
      \item 她睡觉。 (九个小时) $\rightarrow$ \dots
      \item 我们打篮球。 (三年) $\rightarrow$ \dots
      \item 他跑步。 (半小时) $\rightarrow$ \dots
      \end{itemize}
\item Corrige les erreurs dans les phrases suivantes (une erreur par phrase) :
      \begin{itemize}
      \item 我学了中文两年。
      \item 他每天运动了一个小时。
      \item 她不高高兴兴地上学。
      \item 我们踢足球了五年了。
      \item 我的同学认真真地学习。
      \end{itemize}
\item Traduis en chinois :
      \begin{itemize}
      \item Je dors sept heures chaque nuit.
      \item Elle apprend le chinois depuis six mois (et continue).
      \item Il court joyeusement dans le parc.
      \item Nous n'avons pas fait de sport depuis une semaine.
      \end{itemize}
\end{enumerate}
\end{exercice}

\begin{corrige}[Exercice Entraînement individuel]
\begin{enumerate}
\item \textbf{Transformation :}
      \begin{itemize}
      \item 我学法文学了两年了。 (Objet \zh{法语} $\rightarrow$ redoublement \zh{学} + \zh{了} final car action continue).
      \item 她睡觉睡九个小时。 / 她睡了九个小时觉。 (Verbe séparable : redoublement standard ou séparation).
      \item 我们打篮球打了三年了。 (VO \zh{打篮球} $\rightarrow$ redoublement \zh{打} + \zh{了} final).
      \item 他跑步跑半小时。 / 他跑了半小时步。 (VO \zh{跑步} $\rightarrow$ redoublement ou séparation ; pas de \zh{了} final si habitude présente sans insistance sur la continuation).
      \end{itemize}
\item \textbf{Correction :}
      \begin{itemize}
      \item \zh{我学中文学了两年了。} (Règle : Objet présent $\rightarrow$ redoublement verbe).
      \item \zh{他每天运动一个小时。} / \zh{他每天运动运动一个小时。} (Habitude \zh{每天} $\rightarrow$ pas de \zh{了} verbal après le verbe redoublé ; \zh{了} final possible si insistance sur la routine établie : \zh{他每天运动运动一个小时了}).
      \item \zh{她不高高兴兴地上学。} $\rightarrow$ \zh{她不高兴地上学。} / \zh{她不认真地上学。} (Règle : Pas de redoublement AABB en phrase négative simple).
      \item \zh{我们踢足球踢了五年了。} (Manquait le verbe redoublé \zh{踢} avant \zh{了}).
      \item \zh{我的同学认认真真地学习。} (Redoublement AABB \zh{认真} $\rightarrow$ \zh{认认真真} + \zh{地} obligatoire).
      \end{itemize}
\item \textbf{Traduction :}
      \begin{itemize}
      \item 我每天晚上睡觉睡七个小时。 / 我每晚睡七个小时。
      \item 她学中文学了六个月了。
      \item 他高高兴兴地在公园跑步。 / 他高高兴兴地跑步。
      \item 我们一周没运动了。 / 我们已经一周没运动了。
      \end{itemize}
\end{enumerate}
\end{corrige}

\courssub{Bilan de l'activité}

Cette activité t'a permis de mobiliser en situation réelle les deux structures de la leçon. Retiens l'essentiel :

\begin{itemize}
\item Le \cle{complément de durée} se place \textbf{après le verbe} ; quand le verbe a un objet, on redouble le verbe : \zh{学中文学了一年}.
\item Le \zh{了} verbal (après V2) marque l'accompli ; le \zh{了} final (en fin de phrase) indique que l'action \textbf{continue encore} ou qu'un nouvel état est atteint.
\item La négation du complément de durée (action non réalisée sur une période) se fait avec \zh{没(有)} + Verbe + Durée (sans \zh{了}).
\item Le \cle{redoublement des adjectifs} (AA pour monosyllabes, AABB pour dissyllabes) exprime l'intensité, l'affect ou une manière agréable ; devant un verbe, il exige \zh{地}, devant un nom, \zh{的}.
\item \textbf{Interdictions :} Pas de redoublement AABB pour \zh{漂亮}, \zh{年轻}, \zh{可爱}… ; pas de redoublement adjectif en phrase négative simple ; pas de durée avant le verbe.
\item À l'oral comme à l'écrit, la place des mots est la clé : \textbf{Sujet + (Compl. temps) + Verbe (+ Objet) + Verbe + 了 + Durée (+ 了 final)}.
\end{itemize}

Pour t'auto-évaluer, vérifie que ton portrait final contient bien : trois durées correctement placées (dont au moins une structure V+O+V+了+Dur), deux adjectifs redoublés avec \zh{地} ou \zh{的} selon l'emploi, et aucune phrase calquée sur l'ordre français. Si oui, tu maîtrises les compétences de la leçon 16.

\newpage

\coursec{Méthodes et exercices résolus}

\coursec{Leçon 16 — Grammaire : Les compléments de durée ; redoublement des adjectifs}

\begin{textebox}[Dialogue d'observation : Au terrain de sport]
\textbf{Contexte :} \emph{Après les cours, Marc et Li Wei discutent au stade de Mfandena à Yaoundé.}
\begin{enumerate}
\item \zh{马克：李伟，你每天跑步吗？} (\emph{Mǎkè: Lǐ Wěi, nǐ měitiān pǎobù ma?})
\item \zh{李伟：是的，我跑步跑了三年了。} (\emph{Lǐ Wěi: Shì de, wǒ pǎobù pǎo le sān nián le.})
\item \zh{马克：三年？哇，你真坚持！每次跑多长时间？} (\emph{Mǎkè: Sān nián? Wa, nǐ zhēn jiānchí! Měicì pǎo duō cháng shíjiān?})
\item \zh{李伟：一般四十分钟。昨天我踢足球踢了一个半小时，累得高高兴兴地回家睡觉。} (\emph{Lǐ Wěi: Yībān sìshí fēnzhōng. Zuótiān wǒ tī zúqiú tī le yí ge bàn xiǎoshí, lèi de gāogāoxìngxìng de huí jiā shuìjiào.})
\item \zh{马克：足球场的空气清清爽爽的，对吧？} (\emph{Mǎkè: Zúqiúchǎng de kōngqì qīngqīngshuǎngshuǎng de, duì ba?})
\item \zh{李伟：当然！环境干干净净的，心情好好的。} (\emph{Lǐ Wěi: Dāngrán! Huánjìng gāngānjìngjìng de, xīnqíng hǎohǎo de.})
\end{enumerate}
\textbf{Traduction :}
1. Marc : Li Wei, tu cours tous les jours ?
2. Li Wei : Oui, cela fait trois ans que je cours (et je continue).
3. Marc : Trois ans ? Waouh, tu tiens le coup ! Combien de temps tu cours à chaque fois ?
4. Li Wei : Généralement quarante minutes. Hier, j'ai joué au foot pendant une heure et demie, fatigué mais joyeux, je suis rentré dormir.
5. Marc : L'air du terrain est bien frais, non ?
6. Li Wei : Bien sûr ! L'environnement est bien propre, l'humeur est au beau fixe.
\end{textebox}

\courssub{Savoir-faire 1 : Observer et construire le complément de durée (verbe simple)}

\begin{propriete}[Structure du complément de durée — Verbe simple]
Le complément de durée (\cle{时量补语}) indique la \textbf{longueur de temps} qu'une action dure. En chinois, il se place \textbf{après le verbe} (et non avant comme en français « pendant deux ans, j'ai... »).

\medskip
\textbf{Schéma de base (action accomplie / passée) :}
\begin{center}
\begin{tabularx}{\linewidth}{|X|X|X|}\hline
\rowcolor{popblueL} \textbf{Sujet} & \textbf{Verbe + \zh{了}} & \textbf{Durée (Nombre + Classif. + Unité)} \\ \hline
\zh{我} & \zh{学了} & \zh{两年} \\ \hline
\zh{他} & \zh{睡了} & \zh{八个小时} \\ \hline
\end{tabularx}
\end{center}
\zh{我学了两年。} (\emph{Wǒ xué le liǎng nián.}) « J'ai étudié pendant deux ans. »

\medskip
\textbf{Schéma « action qui se prolonge jusqu'à présent » (double \zh{了}) :}
\cle{Sujet + Verbe + 了 + Durée + \zh{了}}
\zh{我学了两年了。} (\emph{Wǒ xué le liǎng nián le.}) « Cela fait deux ans que j'étudie (et je continue). »
\begin{itemize}
\item Le premier \zh{了} (après le verbe) = aspect accompli / action réalisée.
\item Le second \zh{了} (fin de phrase) = changement d'état / situation actuelle qui persiste.
\end{itemize}
\end{propriete}

\begin{methode}[Construire une phrase avec complément de durée (verbe simple)]
\begin{enumerate}
\item \textbf{Identifier le verbe} : est-il simple (sans objet) ? Ex. : \zh{睡觉} (dormir), \zh{等} (attendre), \zh{跑步} (courir — analysé comme V+O voir SF2), \zh{工作} (travailler), \zh{学习} (étudier).
\item \textbf{Choisir l'unité de temps et son classificateur} :
\begin{center}
\begin{tabularx}{\linewidth}{|X|X|X|}\hline
\rowcolor{popgreenL} \textbf{Unité} & \textbf{Classificateur} & \textbf{Exemple} \\ \hline
\zh{分钟} (minute) & — (aucun) & \zh{二十分钟} \\ \hline
\zh{小时} (heure) & \zh{个} & \zh{两个小时 / 两个钟头} \\ \hline
\zh{天} (jour) & — (aucun) & \zh{三天} \\ \hline
\zh{星期} (semaine) & \zh{个} (facultatif) & \zh{一个星期 / 一星期} \\ \hline
\zh{月} (mois calendrier) vs \zh{个月} (durée) & \zh{个} & \textbf{Durée :} \zh{三个月} \quad (\textit{pas} *\zh{三月} = mars) \\ \hline
\zh{年} (an) & — (aucun) & \zh{五年} \\ \hline
\end{tabularx}
\end{center}
\item \textbf{Appliquer l'ordre} : Sujet + V + \zh{了} + Durée. (Si action encore en cours : ajouter \zh{了} final).
\item \textbf{Vérifier} : la durée ne précède \textbf{jamais} le verbe.
\end{enumerate}
\end{methode}

\begin{exemplebox}[Exemples variés — Verbes simples]
\begin{itemize}
\item \zh{我们等了二十分钟。} (\emph{Wǒmen děng le èrshí fēnzhōng.}) — Nous avons attendu 20 min.
\item \zh{他在喀麦隆工作了五年了。} (\emph{Tā zài Kāmàilóng gōngzuò le wǔ nián le.}) — Il travaille au Cameroun depuis 5 ans (et y est encore).
\item \zh{妹妹睡了九个小时。} (\emph{Mèimei shuì le jiǔ ge xiǎoshí.}) — Ma petite sœur a dormi 9 heures.
\end{itemize}
\end{exemplebox}

\begin{exoresolu}[Application : Compléter et traduire]
\textbf{Énoncé.} Complétez avec la durée (entre parenthèses) et traduisez.\\
1. 他们住在杜阿拉 \underline{\hspace{2.5cm}} 。(两年, encore maintenant)\\
2. 老师讲课 \underline{\hspace{2.5cm}} 。(quarante-cinq minutes)\\
3. 我等公共汽车 \underline{\hspace{2.5cm}} 。(half an hour / 半小时)
\tcblower
\textbf{Solution détaillée.}\\
\textbf{1.} Verbe \zh{住} (habiter), action continue $\rightarrow$ double \zh{了}. Durée : \zh{两年} (pas de classif.).\\
\zh{他们住在杜阿拉住了两年了。} (\emph{Tāmen zhù zài Dùālā zhù le liǎng nián le.}) « Ils habitent Douala depuis deux ans. »\\
\textbf{2.} Verbe \zh{讲课} (faire cours), action accomplie $\rightarrow$ un \zh{了}. Durée : \zh{四十五分钟} (pas de classif.).\\
\zh{老师讲课讲了四十五分钟。} (\emph{Lǎoshī jiǎngkè jiǎng le sìshíwǔ fēnzhōng.}) « Le prof a fait cours pendant 45 minutes. » (Répétition du verbe car \zh{讲课} = V+O, voir SF2).\\
\textbf{3.} Verbe \zh{等} (attendre) + objet \zh{公共汽车}. Structure V+O $\rightarrow$ répétition du verbe. Durée : \zh{半小时} (pas de classif. pour \zh{小时}? \textbf{Attention} : \zh{半小时} s'emploie sans \zh{个}, mais \zh{一个半小时} ou \zh{半个小时} prennent \zh{个}).\\
\zh{我等公共汽车等了半小时。} (\emph{Wǒ děng gōnggòngqìchē děng le bàn xiǎoshí.}) « J'ai attendu le bus une demi-heure. »
\end{exoresolu}

\begin{tikzpicture}[node distance=1.2cm and 0.8cm, every node/.style={font=\small}]
\node[draw=popblue, thick, rounded corners, fill=popblueL, align=center] (S) {Sujet};
\node[draw=poporange, thick, rounded corners, fill=poporangeL, align=center, right=of S] (V) {Verbe};
\node[draw=poppurple, thick, rounded corners, fill=poppurpleL, align=center, right=of V] (LE1) {\zh{了} (aspect)};
\node[draw=popgreen, thick, rounded corners, fill=popgreenL, align=center, right=of LE1] (Dur) {Durée\\Nombre + Classif. + Unité};
\node[draw=poppink, thick, rounded corners, fill=poppinkL, align=center, right=of Dur] (LE2) {\zh{了} (mode)\\(si action continue)};
\draw[-{Stealth[length=3mm]}, thick, popdark] (S) -- (V);
\draw[-{Stealth[length=3mm]}, thick, popdark] (V) -- (LE1);
\draw[-{Stealth[length=3mm]}, thick, popdark] (LE1) -- (Dur);
\draw[-{Stealth[length=3mm]}, thick, popdark, dashed] (Dur) -- (LE2);
\node[align=center, text=popmuted, font=\footnotesize] at ($(LE2)+(0,-1)$) {Flèche continue = obligatoire\\Flèche pointillée = si action prolonge};
\end{tikzpicture}
\legende{Schéma de l'ordre des mots : Complément de durée (verbe simple)}

\courssub{Savoir-faire 2 : Gérer le complément de durée avec un verbe transitif (Verbe + Objet)}

\begin{propriete}[Contrainte du bloc « Verbe + Objet »]
En chinois, le verbe et son objet forment un \textbf{bloc inséparable} (\cle{动宾结构}). On \textbf{ne peut pas} insérer la durée entre eux.
\begin{center}
\textbf{Interdit :} *\zh{我看了两个小时电视。} \quad \textbf{Interdit :} *\zh{他踢了一个小时足球。}
\end{center}
Trois stratégies grammaticales existent ; la \textbf{répétition du verbe} (Stratégie 1) est la norme attendue au lycée (HSK 3-4).
\end{propriete}

\begin{methode}[Trois solutions pour « Verbe + Objet + Durée »]
\begin{enumerate}
\item \textbf{Répétition du verbe (Standard HSK / Bac)} :\\
\cle{Sujet + [Verbe + Objet] + Verbe + 了 + Durée}\\
\zh{我看电视看了两个小时。} (\emph{Wǒ kàn diànshì kàn le liǎng ge xiǎoshí.}) « J'ai regardé la télé 2h. »\\
$\rightarrow$ Le premier bloc \zh{看电视} pose le sujet ; le second \zh{看} reçoit la durée.
\item \textbf{Thématisation de l'objet (Topicalisation)} :\\
\cle{Objet + ，+ Sujet + Verbe + 了 + Durée}\\
\zh{电视，我看了两个小时。} (\emph{Diànshì, wǒ kàn le liǎng ge xiǎoshí.}) « La télé, je l'ai regardée 2h. »\\
$\rightarrow$ Style plus oral / mis en relief. Virgule obligatoire après l'objet thématisé.
\item \textbf{Structure avec \zh{的} (Lourde, formelle, rare à l'oral)} :\\
\cle{Sujet + Verbe + 了 + Durée + \zh{的} + Objet}\\
\zh{我看了两个小时的电视。} (\emph{Wǒ kàn le liǎng ge xiǎoshí de diànshí.}) « J'ai regardé 2h de télévision. »\\
$\rightarrow$ \zh{的} nominalise la phrase verbale. \textbf{À éviter} en production écrite au lycée.
\end{enumerate}
\textbf{Règle d'or :} En cas de doute, \textbf{répétez le verbe}. Cela fonctionne pour 100\% des cas (verbes simples traités comme V+O comme \zh{跑步}, \zh{睡觉}, \zh{吃饭}, \zh{踢足球}, \zh{学汉语}).
\end{methode}

\begin{exemplebox}[Verbes « V+O » courants à traiter par répétition]
\begin{center}
\begin{tabularx}{\linewidth}{|X|X|X|}\hline
\rowcolor{popblueL} \textbf{Verbe + Objet} & \textbf{Pinyin} & \textbf{Sens} \\ \hline
\zh{跑步} & pǎobù & courir (pied + étape) \\
\zh{睡觉} & shuìjiào & dormir (sommeil + réveil) \\
\zh{吃饭} & chīfàn & manger (manger + riz/repas) \\
\zh{开会} & kāihuì & faire une réunion (ouvrir + réunion) \\
\zh{跳舞} & tiàowǔ & danser (sauter + danse) \\
\zh{见面} & jiànmiàn & se voir / se rencontrer (voir + face) \\
\zh{学习} / \zh{学中文} & xuéxí / xué Zhōngwén & étudier / étudier le chinois \\
\zh{踢足球} & tī zúqiú & jouer au foot \\
\hline
\end{tabularx}
\end{center}
\end{exemplebox}

\begin{exoresolu}[Transformation : Verbe + Objet $\rightarrow$ + Durée]
\textbf{Énoncé.} Transformez en ajoutant la durée (répétition du verbe). Traduisez.\\
1. 弟弟看书。(一个小时)\\
2. 我们在雅温得学法语。(三个月, action continue)\\
3. 妈妈做饭。(deux heures)
\tcblower
\textbf{Solution détaillée.}\\
\textbf{1.} \zh{看书} (lire) $\rightarrow$ \zh{弟弟看书看了一个小时。} (\emph{Dìdi kànshū kàn le yí ge xiǎoshí.}) « Mon frère a lu pendant une heure. »\\
\textbf{2.} \zh{学法语} (apprendre le français), action continue $\rightarrow$ double \zh{了}. \zh{三个月} (classif. \zh{个} obligatoire).\\
\zh{我们在雅温得学法语学了三个月了。} (\emph{Wǒmen zài Yǎwēndé xué Fǎyǔ xué le sān ge yuè le.}) « Cela fait 3 mois que nous apprenons le français à Yaoundé. »\\
\textbf{3.} \zh{做饭} (cuisiner) $\rightarrow$ \zh{妈妈做饭做了两个小时。} (\emph{Māma zuòfàn zuò le liǎng ge xiǎoshí.}) « Maman a cuisiné pendant deux heures. »
\end{exoresolu}

\courssub{Savoir-faire 3 : Négation, interrogation et pièges du complément de durée}

\begin{propriete}[Négation et Interrogation du complément de durée]
\begin{center}
\begin{tabularx}{\linewidth}{|X|X|X|}\hline
\rowcolor{popblueL} \textbf{Type} & \textbf{Structure} & \textbf{Exemple} \\ \hline
\textbf{Affirmative (accompli)} & S + V(+O) + V + \zh{了} + Durée & \zh{他跑步跑了三十分钟。} \\ \hline
\textbf{Négation (accompli)} & S + \zh{没(有)} + V(+O) + V + Durée \textbf{(SANS \zh{了})} & \zh{他跑步没有跑三十分钟。} \\ \hline
\textbf{Question \zh{多长时间} / \zh{多久}} & S + V(+O) + V + \zh{了} + \zh{多长时间 / 多久} ? & \zh{你跑步跑了多长时间？} \\ \hline
\textbf{Question Oui/Non (\zh{吗})} & Phrase affirmative + \zh{吗} ? & \zh{你跑步跑了三十分钟吗？} \\ \hline
\end{tabularx}
\end{center}
\textbf{Rappel crucial :} \zh{不} nie l'habitude / le futur / l'état (\zh{我不跑步} = « je ne cours pas / je ne cours jamais »). \zh{没(有)} nie le \textbf{passé / l'accompli} (\zh{我没跑步} = « je n'ai pas couru »).
\end{propriete}

\begin{attention}[Pièges fréquents des francophones — Négation \& Durée]
\begin{enumerate}
\item \textbf{* \zh{他不跑了三十分钟。}} $\rightarrow$ \zh{不} + \zh{了} est impossible. \textbf{Correction :} \zh{他没有跑三十分钟。} (ou \zh{他跑步没有跑三十分钟。})
\item \textbf{* \zh{他没跑了三十分钟。}} $\rightarrow$ \zh{没} nie le verbe, donc \zh{了} (aspect accompli) disparaît. \textbf{Correction :} \zh{他没跑三十分钟。}
\item \textbf{Confusion \zh{多少时间} vs \zh{多长时间}} : \zh{多少} + nombre comptable (\zh{多少人}, \zh{多少天}) ; \zh{多长时间} / \zh{多久} pour une durée indéfinie.
\item \textbf{Oubli du classif. \zh{个}} : *\zh{两小时} $\rightarrow$ \zh{两个钟头} ou \zh{两个小时}. *\zh{三星期} $\rightarrow$ \zh{三个星期} (recommandé).
\end{enumerate}
\end{attention}

\begin{exoresolu}[Négation et Question]
\textbf{Énoncé.} 1. Nier : \zh{学生们踢足球踢了一个下午。}\\
2. Interroger sur la durée : \zh{她学汉语学了四年。}
\tcblower
\textbf{Solution détaillée.}\\
\textbf{1.} Affirmative : V+O répété + \zh{了}. Négation : \zh{没(有)} devant le verbe répété, \textbf{supprimer \zh{了}}.\\
\zh{学生们踢足球没有踢一个下午。} (\emph{Xuéshengmen tī zúqiú méiyǒu tī yí ge xiàwǔ.}) « Les élèves n'ont pas joué au foot tout l'après-midi. »\\
\textbf{2.} Remplacer \zh{四年} par \zh{多长时间} (ou \zh{多久}).\\
\zh{她学汉语学了多长时间？} (\emph{Tā xué Hànyǔ xué le duō cháng shíjiān?}) « Depuis combien de temps / Combien de temps étudie-t-elle le chinois ? »
\end{exoresolu}

\courssub{Savoir-faire 4 : Maîtriser le redoublement des adjectifs (AA / AABB)}

\begin{propriete}[Règles du redoublement des adjectifs]
Le redoublement (\cle{叠词}) renforce la valeur descriptive, ajoute une nuance \textbf{affective, vivante, esthétique} (« tout... », « bien... », « joliment... »). Il ne s'utilise \textbf{pas} dans la négation, ni avec \zh{比} (comparatif), ni pour les adjectifs « absolus » (\zh{对}, \zh{错}, \zh{真}, \zh{假}, \zh{死}, \zh{活}).

\medskip
\textbf{1. Adjectif monosyllabique $\rightarrow$ Schéma \cle{AA + 的 / 地}}
\begin{center}
\begin{tabularx}{\linewidth}{|X|X|X|X|}\hline
\rowcolor{poppurpleL} \textbf{Base} & \textbf{Redoublé} & \textbf{Épithète (\zh{的})} & \textbf{Manière (\zh{地})} \\ \hline
\zh{大} & \zh{大大} & \zh{大大的眼睛} (de grands yeux) & \zh{大大地改变} (changer radicalement) \\ \hline
\zh{小} & \zh{小小} & \zh{小小的鸟} (un petit oiseau) & — \\ \hline
\zh{高} & \zh{高高} & \zh{高高的山} (haute montagne) & — \\ \hline
\zh{红} & \zh{红红} & \zh{红红的脸} (visage tout rouge) & — \\ \hline
\zh{好} & \zh{好好} & \zh{好好的} (bien / correctement) & \zh{好好地学习} (étudier sérieusement) \\ \hline
\zh{慢} & \zh{慢慢} & — & \zh{慢慢地走} (marcher lentement) \\ \hline
\end{tabularx}
\end{center}

\textbf{2. Adjectif dissyllabique $\rightarrow$ Schéma \cle{AABB + 的 / 地}}
\begin{center}
\begin{tabularx}{\linewidth}{|X|X|X|X|}\hline
\rowcolor{poppurpleL} \textbf{Base} & \textbf{Redoublé (AABB)} & \textbf{Épithète (\zh{的})} & \textbf{Manière (\zh{地})} \\ \hline
\zh{高兴} & \zh{高高兴兴} & \zh{高高兴兴的孩子} & \zh{高高兴兴地回家} \\ \hline
\zh{干净} & \zh{干干净净} & \zh{干干净净的屋子} & \zh{干干净净地打扫} \\ \hline
\zh{舒服} & \zh{舒舒服服} & \zh{舒舒服服的床} & \zh{舒舒服服地睡} \\ \hline
\zh{漂亮} & \zh{漂漂亮亮} & \zh{漂漂亮亮的裙子} & — \\ \hline
\zh{清爽} & \zh{清清爽爽} & \zh{清清爽爽的空气} & — \\ \hline
\zh{认真} & \zh{认认真真} & \zh{认认真真的态度} & \zh{认认真真地听} \\ \hline
\end{tabularx}
\end{center}

\textbf{3. Emplois syntaxiques}
\begin{itemize}
\item \textbf{Épithète (devant nom)} : Adj-redoublé + \zh{的} + Nom. \zh{一个干干净净的教室。}
\item \textbf{Attribut (après sujet, verbe \zh{是} ou nul)} : Sujet + Adj-redoublé + \zh{的}. \zh{教室干干净净的。}
\item \textbf{Complément de manière (devant verbe)} : Adj-redoublé + \zh{地} + Verbe. \zh{他高高兴兴地走了。}
\end{itemize}
\end{propriete}

\begin{savaistu}[Le redoublement dans la culture chinoise]
Le redoublement est omniprésent en chinois (noms : \zh{人人} « tout le monde », \zh{天天} « chaque jour » ; verbes : \zh{看看} « jeter un coup d'œil » ; adjectifs). Pour les adjectifs, il reflète une sensibilité esthétique : décrire non pas une mesure objective (\zh{高} = 1,80m) mais une \textbf{impression subjective} (\zh{高高的} = « qui paraît grand, imposant, beau »). C'est pourquoi on l'évite dans le langage scientifique, juridique ou la négation objective.
\end{savaistu}

\begin{exoresolu}[Redoublement : Forme et emploi]
\textbf{Énoncé.} 1. Redoublez et complétez avec \zh{的} ou \zh{地} :\\
a) 天气 \underline{\hspace{2cm}} 。(冷)\\
b) 妹妹 \underline{\hspace{2cm}} 地笑。(开心)\\
c) 这是一本 \underline{\hspace{2cm}} 书。(有意思 — \textbf{attention, trisyllabique !})\\
2. Traduisez : « La salle de classe est bien rangée / bien propre. » (\zh{整齐} / \zh{干净})
\tcblower
\textbf{Solution détaillée.}\\
\textbf{1.a)} \zh{冷} (monosyll.) $\rightarrow$ \zh{冷冷} + \zh{的} (attribut). \zh{天气冷冷的。} (\emph{Tiānqì lěnglěng de.}) « Il fait bien froid. »\\
\textbf{1.b)} \zh{开心} (dissyll.) $\rightarrow$ \zh{开开心心} + \zh{地} (manière). \zh{妹妹开开心心地笑。} (\emph{Mèimei kāikāixīnxīn de xiào.}) « Ma sœur rit joyeusement. »\\
\textbf{1.c)} \zh{有意思} (3 syllabes) \textbf{ne se redouble pas} (sauf forme rare \zh{有有意思意思} non standard). On utilise \zh{很有意思} ou \zh{挺有意思的}. \textbf{Réponse :} \zh{这是一本很有意思的书。}\\
\textbf{2.} \zh{干净} $\rightarrow$ \zh{干干净净}. Attribut $\rightarrow$ +\zh{的}.\\
\zh{教室干干净净的。} (\emph{Jiàoshì gāngānjìngjìng de.}) « La salle de classe est bien propre. »
\end{exoresolu}

\courssub{Savoir-faire 5 : Production intégrée — Mini-texte « Ma journée santé »}

\begin{methode}[Rédiger un texte cohérent : Durée + Redoublement]
\begin{enumerate}
\item \textbf{Sujet} : Une routine (sport, étude, loisir) ou un événement passé (match, fête, voyage).
\item \textbf{Squelette} : 3 à 5 phrases. 1. Cadre spatio-temporel. 2. Action 1 + Durée (V simple). 3. Action 2 + Durée (V+O). 4. Description sensorielle/humeur (Adj redoublé). 5. Bilan / Sentiment.
\item \textbf{Check-list avant de rendre} :
\begin{itemize}
\item[$\square$] Durées \textbf{après} le verbe (et après \zh{了}).
\item[$\square$] Verbe + Objet $\rightarrow$ \textbf{verbe répété} + \zh{了} + Durée.
\item[$\square$] Action continue $\rightarrow$ \textbf{double \zh{了}}.
\item[$\square$] Adj. redoublé monosyll. = \textbf{AA} ; dissyll. = \textbf{AABB}.
\item[$\square$] \zh{的} (épithète/attribut) ou \zh{地} (manière) \textbf{obligatoire} après l'adj. redoublé.
\item[$\square$] Classificateurs \zh{个} (\zh{小时}, \zh{星期}, \zh{个月}) corrects.
\end{itemize}
\end{enumerate}
\end{methode}

\begin{experience}[Atelier oral : « Enquête : Mes habitudes de vie »]
\textbf{Consigne :} En binôme, interrogez votre partenaire pour remplir le tableau. Utilisez \zh{多长时间} / \zh{多久} et notez la réponse avec complément de durée. Puis présentez un résumé à la classe (30 sec) en utilisant au moins un adjectif redoublé.
\begin{center}
\begin{tabularx}{\linewidth}{|X|X|X|}\hline
\rowcolor{popblueL} \textbf{Activité} & \textbf{Question type} & \textbf{Réponse notée (Durée)} \\ \hline
Dormir & \zh{你每晚睡多久？} & \zh{八个小时} \\ \hline
Étudier le chinois & \zh{你学中文学了多长时间？} & \zh{两年了} \\ \hline
Faire du sport & \zh{你每次跑步跑多久？} & \zh{半个小时} \\ \hline
Regarder téléphone / TV & \zh{你看手机看了多长时间？} & \zh{一个小时} \\ \hline
\end{tabularx}
\end{center}
\textbf{Modèle de restitution :} \zh{我的搭档叫... 他每晚睡八个小时。他学中文学了两年了。他跑步跑半个小时。他看手机看了一个小时。我觉得他生活很规规矩矩的。}
\end{experience}

\begin{exoresolu}[Composition guidée : « Ma matinée de sport » — Version enrichie]
\textbf{Énoncé.} Rédigez 5-6 phrases. Contraintes : (a) 1 durée + V simple ; (b) 1 durée + V+O (répétition) ; (c) 1 adj. redoublé (épithète ou manière) ; (d) 1 phrase avec double \zh{了}.
\tcblower
\textbf{Solution détaillée (Modèle de rédaction).}\\
\textbf{Étape 1 — Plan :} Lever tôt (cadre) $\rightarrow$ Footing (V simple + durée) $\rightarrow$ Basket (V+O + durée) $\rightarrow$ Ambiance (Adj redoublé) $\rightarrow$ Bilan (double \zh{了} pour habitude).\\
\textbf{Étape 2 — Rédaction :}\\
1. \zh{今天周六，我六点半起床。} (\emph{Jīntiān zhōuliù, wǒ liù diǎn bàn qǐchuáng.}) « Aujourd'hui samedi, levé 6h30. »\\
2. \zh{我先跑步，跑了三十分钟。} (\emph{Wǒ xiān pǎobù, pǎo le sānshí fēnzhōng.}) « D'abord footing, couru 30 min. » (V simple \zh{跑} répété car \zh{跑步} = V+O).\\
3. \zh{然后和同学打篮球，打了一个小时。} (\emph{Ránhòu hé tóngxué dǎ lánqiú, dǎ le yí ge xiǎoshí.}) « Ensuite basket avec camarades, joué 1h. » (V+O \zh{打篮球} $\rightarrow$ répète \zh{打}).\\
4. \zh{球场的空气清清爽爽的，心情好好的。} (\emph{Qiúchǎng de kōngqì qīngqīngshuǎngshuǎng de, xīnqíng hǎohǎo de.}) « Air bien frais, humeur au top. » (AABB + \zh{的} attribut ; AA + \zh{的} attribut).\\
5. \zh{我坚持锻炼坚持了两年了。} (\emph{Wǒ jiānchí duànliàn jiānchí le liǎng nián le.}) « Je persévère dans l'exercice depuis 2 ans. » (Double \zh{了} = habitude continue).\\
6. \zh{虽然累，但我舒舒服服地睡了一觉。} (\emph{Suīrán lèi, dàn wǒ shūshūfuwu de shuì le yí giào.}) « Bien que fatigué, j'ai fait un bon sommeil réparateur. » (AABB + \zh{地} manière + \zh{睡了一觉} V+O + durée courte).\\
\textbf{Étape 3 — Vérification :} Contraintes (a) \zh{跑了三十分钟} ✓ (b) \zh{打篮球打了...} ✓ (c) \zh{清清爽爽的}, \zh{好好的}, \zh{舒舒服服地} ✓ (d) \zh{坚持了两年了} ✓. Classificateurs \zh{个} (\zh{小时}) ✓. Tons vérifiés ✓.
\end{exoresolu}

\begin{attention}[Bilan des 6 erreurs fatales pour le Bac / Contrôle]
\begin{enumerate}
\item \textbf{Ordre des mots (Calque français) :} *\zh{两年我学了汉语} $\rightarrow$ \textbf{\zh{我学汉语学了两年。}} (Durée APRÈS verbe).
\item \textbf{Insertion interdite :} *\zh{我看了两个小时电视} $\rightarrow$ \textbf{\zh{我看电视看了两个小时。}} (Répétition verbe).
\item \textbf{Négation \zh{不} vs \zh{没} :} *\zh{他不睡了八个小时} $\rightarrow$ \textbf{\zh{他没有睡八个小时。}} (Pas de \zh{了} avec \zh{没}).
\item \textbf{Classificateurs de durée :} *\zh{三月} (Mars), *\zh{两小时}, *\zh{三星期} $\rightarrow$ \textbf{\zh{三个月}, \zh{两个小时}, \zh{三个星期}}.
\item \textbf{Redoublement sans particule :} *\zh{高高兴兴回家}, *\zh{大大眼睛} $\rightarrow$ \textbf{\zh{高高兴兴地回家}, \zh{大大的眼睛}} (\zh{的}/\zh{地} obligatoires).
\item \textbf{Redoublement inadapté :} *\zh{不对不对的}, *\zh{比高高的} $\rightarrow$ \textbf{Jamais en négation, jamais avec \zh{比}, jamais sur adjectifs absolus.}
\end{enumerate}
\end{attention}

\begin{exercice}[Entraînement autonome — Devoir / Classeur]
\textbf{Exercice 1 — Complément de durée (Verbe simple).} Mettez les mots dans l'ordre et ajoutez \zh{了} si nécessaire. Traduisez.\\
1. 看书 / 我 / 两个小时 / 了\\
2. 在北京 / 住 / 王先生 / 五年 / 了 (action continue)\\
3. 等 / 我们 / 三十分钟\\
\textbf{Exercice 2 — Verbe + Objet + Durée.} Transformez par répétition du verbe.\\
1. 妹妹弹钢琴 (deux heures).\\
2. 他们学游泳 (un mois, continue).\\
3. 爸爸看报纸 (vingt minutes).\\
\textbf{Exercice 3 — Négation et Question.} Transformez la phrase affirmative.\\
Affirmative : \zh{哥哥上网上了三个小时。}\\
a) Forme négative. \quad b) Question sur la durée. \quad c) Question Oui/Non.\\
\textbf{Exercice 4 — Redoublement des adjectifs.} Complétez avec la forme redoublée correcte + \zh{的} ou \zh{地}.\\
1. 孩子们 \underline{\hspace{2cm}} 玩耍。(开心)\\
2. 这条河的水 \underline{\hspace{2cm}} 。(清澈 — qīngchè)\\
3. 请把字写 \underline{\hspace{2cm}} 。(工整 — gōngzhěng)\\
4. 我买了一件 \underline{\hspace{2cm}} 衬衫。(漂亮)\\
\textbf{Exercice 5 — Traduction thème (Français $\rightarrow$ Chinois).}\\
1. J'attends l'autobus depuis vingt minutes (et j'attends encore).\\
2. Ma mère a cuisiné pendant une heure et demie.\\
3. Le terrain de foot est bien vert (vert vert).\\
4. Il a couru joyeusement pendant quarante minutes.\\
5. Cela fait trois ans que j'étudie le chinois au lycée.
\end{exercice}

\begin{corrige}[Exercices 1 à 5]
\textbf{Exercice 1}\\
1. \zh{我看书看了两个小时。} (\emph{Wǒ kànshū kàn le liǎng ge xiǎoshí.}) — J'ai lu deux heures. (V+O \zh{看书} $\rightarrow$ répète \zh{看}).\\
2. \zh{王先生在北京住了五年了。} (\emph{Wáng Xiānsheng zài Běijīng zhù le wǔ nián le.}) — M. Wang habite Pékin depuis 5 ans. (Double \zh{了}).\\
3. \zh{我们等了三十分钟。} (\emph{Wǒmen děng le sānshí fēnzhōng.}) — Nous avons attendu 30 min. (Verbe simple \zh{等}).\\

\textbf{Exercice 2}\\
1. \zh{妹妹弹钢琴弹了两个小时。} (\emph{Mèimei tán gāngqín tán le liǎng ge xiǎoshí.})\\
2. \zh{他们学游泳学了一个月了。} (\emph{Tāmen xué yóuyǒng xué le yí ge yuè le.}) — \zh{个月} obligatoire. Double \zh{了}.\\
3. \zh{爸爸看报纸看了二十分钟。} (\emph{Bàba kàn bàozhǐ kàn le èrshí fēnzhōng.})\\

\textbf{Exercice 3}\\
Affirmative : \zh{哥哥上网上了三个小时。}\\
a) Négation : \zh{哥哥上网没有上三个小时。} (\emph{Gēge shàngwǎng méiyǒu shàng sān ge xiǎoshí.}) — \zh{没(有)} + V répété, SANS \zh{了}.\\
b) Question durée : \zh{哥哥上网上了多长时间？} (\emph{Gēge shàngwǎng shàng le duō cháng shíjiān?})\\
c) Question Oui/Non : \zh{哥哥上网上了三个小时吗？} (\emph{Gēge shàngwǎng shàng le sān ge xiǎoshí ma?})\\

\textbf{Exercice 4}\\
1. \zh{开开心心地} (AABB + \zh{地} manière). \zh{孩子们开开心心地玩耍。}\\
2. \zh{清清澈澈的} (AABB + \zh{的} attribut). \zh{这条河的水清清澈澈的。}\\
3. \zh{工工整整地} (AABB + \zh{地} manière). \zh{请把字写工工整整地。}\\
4. \zh{漂漂亮亮的} (AABB + \zh{的} épithète). \zh{我买了一件漂漂亮亮的衬衫。}\\
\textit{Note :} \zh{工整} et \zh{清澈} sont dissyllabiques $\rightarrow$ AABB.\\

\textbf{Exercice 5}\\
1. \zh{我等公共汽车等了二十分钟了。} (V+O \zh{等公共汽车} $\rightarrow$ répète \zh{等} ; double \zh{了}).\\
2. \zh{妈妈做饭做了一个半小时。} (\zh{做饭} V+O $\rightarrow$ répète \zh{做} ; \zh{一个半小时}).\\
3. \zh{足球场的草绿绿的。} ou \zh{足球场绿绿的。} (Monosyll. \zh{绿} $\rightarrow$ AA + \zh{的} attribut).\\
4. \zh{他高高兴兴地跑了四十分钟。} (AABB + \zh{地} manière ; V simple \zh{跑} + \zh{了} + durée).\\
5. \zh{我在高中学中文学了三年了。} (V+O \zh{学中文} $\rightarrow$ répète \zh{学} ; double \zh{了} ; \zh{三年} sans classif.).
\end{corrige}

\begin{aretenir}[Fiche de révision express — Leçon 16]
\begin{itemize}
\item \textbf{Durée (Verbe simple) :} S + V + \zh{了} + Durée. \quad \textbf{Prolongation :} + \zh{了} final.
\item \textbf{Durée (V + Objet) :} S + V+O + V + \zh{了} + Durée. \quad \textbf{Règle :} Répéter le verbe.
\item \textbf{Négation passée :} \zh{没(有)} + V (+ V répété) + Durée \textbf{sans \zh{了}}.
\item \textbf{Question durée :} \zh{多长时间} / \zh{多久} à la place de la durée.
\item \textbf{Classif. clés :} \zh{个} avec \zh{小时, 星期, 个月} ; \textbf{sans} classif. avec \zh{年, 天, 分钟}.
\item \textbf{Redoublement Mono :} AA + \zh{的}/\zh{地} (\zh{高高的}, \zh{慢慢地}).
\item \textbf{Redoublement Di :} AABB + \zh{的}/\zh{地} (\zh{高高兴兴地}, \zh{干干净净的}).
\item \textbf{Interdits redoublement :} Négation, \zh{比}, adjectifs absolus (\zh{对/错}), trisyllabiques.
\end{itemize}
\end{aretenir}

\newpage

\coursec{Exercices}

\courssub{Je vérifie mes connaissances}

\begin{exercice}[QCM : la bonne phrase]
Pour chaque question, choisis la phrase correcte (une seule réponse).

\begin{enumerate}
\item Pour dire « Elle a vécu trois mois en Chine » :
\begin{itemize}
\item[A.] \zh{她住了三个月在中国。}
\item[B.] \zh{她在中国住了三个月。}
\item[C.] \zh{她三个月住了在中国。}
\end{itemize}
\item Pour dire « Nous avons cours pendant quarante minutes » :
\begin{itemize}
\item[A.] \zh{我们上课上了四十分钟。}
\item[B.] \zh{我们四十分钟上课了。}
\item[C.] \zh{我们上了课四十分钟了。}
\end{itemize}
\item Pour dire « Parle lentement, s'il te plaît » :
\begin{itemize}
\item[A.] \zh{请你慢地说。}
\item[B.] \zh{请你慢慢地说。}
\item[C.] \zh{请你地说慢慢。}
\end{itemize}
\item Pour dire « Il est rentré joyeusement à la maison » :
\begin{itemize}
\item[A.] \zh{他高兴高兴地回家了。}
\item[B.] \zh{他高高兴兴回家了地。}
\item[C.] \zh{他高高兴兴地回家了。}
\end{itemize}
\end{enumerate}
\end{exercice}

\begin{exercice}[Vrai ou faux ? Justifie ta réponse]
Dis si chaque phrase est correcte (V) ou incorrecte (F), puis justifie en français en citant la règle de grammaire concernée.

\begin{enumerate}
\item \zh{我没学了两年的中文。} \emph{(Wǒ méi xué le liǎng nián de Zhōngwén.)}
\item \zh{她看电视看了一个小时。} \emph{(Tā kàn diànshì kàn le yí ge xiǎoshí.)}
\item \zh{很大大大的眼睛。} \emph{(Hěn dà dà dà de yǎnjing.)}
\item \zh{他每天睡觉睡八个小时。} \emph{(Tā měi tiān shuìjiào shuì bā ge xiǎoshí.)}
\end{enumerate}
\end{exercice}

\begin{exercice}[Vocabulaire : durées et adjectifs]
Complète le tableau suivant en donnant le pinyin et la traduction française de chaque expression.

\begin{center}
\begin{tabularx}{\linewidth}{|X|X|X|X|}
\hline
\rowcolor{popblueL} Caractères & Pinyin & Nature & Traduction \\
\hline
三个月 & \ldots\ldots & \ldots\ldots & \ldots\ldots \\
\hline
两年 & \ldots\ldots & \ldots\ldots & \ldots\ldots \\
\hline
四十分钟 & \ldots\ldots & \ldots\ldots & \ldots\ldots \\
\hline
慢慢地 & \ldots\ldots & \ldots\ldots & \ldots\ldots \\
\hline
高高兴兴 & \ldots\ldots & \ldots\ldots & \ldots\ldots \\
\hline
一个小时 & \ldots\ldots & \ldots\ldots & \ldots\ldots \\
\hline
\end{tabularx}
\end{center}
\end{exercice}

\begin{exercice}[Pinyin et caractères : à associer]
Associe chaque transcription pinyin (1 à 6) au bon groupe de caractères (A à F).

\begin{enumerate}
\item \emph{sān ge yuè}
\item \emph{mànman de}
\item \emph{bā ge xiǎoshí}
\item \emph{gāogāoxìngxìng}
\item \emph{liǎng nián}
\item \emph{sìshí fēnzhōng}
\end{enumerate}

\begin{itemize}
\item[A.] \zh{高高兴兴}
\item[B.] \zh{两年}
\item[C.] \zh{三个月}
\item[D.] \zh{四十分钟}
\item[E.] \zh{慢慢地}
\item[F.] \zh{八个小时}
\end{itemize}
\end{exercice}

\courssub{Je m'entraîne}

\begin{exercice}[Remise en ordre]
Remets les mots dans le bon ordre pour former une phrase correcte, puis traduis-la en français.

\begin{enumerate}
\item 了 / 她 / 三个月 / 住 / 在 / 中国
\item 一个小时 / 看 / 了 / 他 / 电视 / 看
\item 两年 / 学 / 我 / 中文 / 了
\item 地 / 请你 / 慢慢 / 说
\end{enumerate}
\end{exercice}

\begin{exercice}[Complète avec la durée]
Complète chaque phrase avec la durée indiquée entre parenthèses, en utilisant correctement le complément de durée (attention au verbe répété si nécessaire).

\begin{enumerate}
\item 我们上课\_\_\_\_\_\_。（四十分钟）
\item 他每天睡觉\_\_\_\_\_\_。（八个小时）
\item 我在雅温得住\_\_\_\_\_\_。（两年）
\item 她学习汉语\_\_\_\_\_\_。（三个月）
\end{enumerate}
\end{exercice}

\begin{exercice}[Transforme les adjectifs]
Transforme l'adjectif entre parenthèses en adjectif redoublé avec \zh{地} et complète la phrase, puis traduis.

\begin{enumerate}
\item 请你\_\_\_\_\_\_说。（慢）
\item 他\_\_\_\_\_\_回家了。（高兴）
\item 老师\_\_\_\_\_\_写字。（认真）
\item 我们\_\_\_\_\_\_走进了教室。（快）
\end{enumerate}
\end{exercice}

\begin{exercice}[Corrige les erreurs]
Chaque phrase contient une erreur typique des francophones. Identifie l'erreur, corrige la phrase et explique la règle en français.

\begin{enumerate}
\item ✱我没学了两年的中文。
\item ✱很大大大的眼睛。
\item ✱她在中国三个月住了。
\item ✱请你慢地说。
\end{enumerate}
\end{exercice}

\courssub{Je me prépare au Bac}

\begin{exercice}[Texte et questions de compréhension (type Bac)]
Lis le texte suivant, puis réponds aux questions en français.

\begin{textebox}[Ma vie à Yaoundé]
我叫玛丽，是雅温得一所中学的学生。\\
\emph{Wǒ jiào Mǎlì, shì Yǎwēndé yì suǒ zhōngxué de xuésheng.}\\
Je m'appelle Marie, je suis élève dans un lycée de Yaoundé.\\[4pt]
我学习汉语学了两年了。每天我们上课上四十分钟。\\
\emph{Wǒ xuéxí Hànyǔ xué le liǎng nián le. Měi tiān wǒmen shàngkè shàng sìshí fēnzhōng.}\\
J'apprends le chinois depuis deux ans. Chaque jour, nous avons cours pendant quarante minutes.\\[4pt]
老师总是慢慢地说，我们高高兴兴地学习。\\
\emph{Lǎoshī zǒngshì mànman de shuō, wǒmen gāogāoxìngxìng de xuéxí.}\\
Le professeur parle toujours lentement, et nous apprenons joyeusement.\\[4pt]
晚上，我看电视看一个小时，然后睡觉睡八个小时。\\
\emph{Wǎnshang, wǒ kàn diànshì kàn yí ge xiǎoshí, ránhòu shuìjiào shuì bā ge xiǎoshí.}\\
Le soir, je regarde la télévision pendant une heure, puis je dors huit heures.
\end{textebox}

\textbf{Questions de compréhension}
\begin{enumerate}
\item Depuis combien de temps Marie apprend-elle le chinois ?
\item Combien de temps dure un cours dans son lycée ?
\item Comment le professeur parle-t-il ? Comment les élèves apprennent-ils ?
\item Combien de temps Marie dort-elle chaque nuit ?
\end{enumerate}

\textbf{Questions de langue}
\begin{enumerate}
\item Releve dans le texte deux compléments de durée et recopie leur schéma (verbe + 了 + durée).
\item Releve dans le texte un adjectif redoublé et explique sa fonction.
\item Transforme : \zh{我们上课。}（四十分钟）en utilisant le complément de durée.
\end{enumerate}
\end{exercice}

\begin{exercice}[Thème et version (type Bac)]
\textbf{Thème.} Traduis en chinois (caractères simplifiés) les phrases suivantes :
\begin{enumerate}
\item « J'apprends le chinois depuis deux ans. »
\item « Il n'a pas dormi huit heures. »
\item « Nous avons regardé le match de football pendant deux heures. »
\end{enumerate}

\textbf{Version.} Traduis en français les phrases suivantes :
\begin{enumerate}
\item \zh{他高高兴兴地回家了。} \emph{(Tā gāogāoxìngxìng de huí jiā le.)}
\item \zh{她看电视看了一个小时。} \emph{(Tā kàn diànshì kàn le yí ge xiǎoshí.)}
\item \zh{请你慢慢地写字。} \emph{(Qǐng nǐ mànman de xiězì.)}
\end{enumerate}
\end{exercice}

\begin{exercice}[Expression écrite guidée (type Bac)]
\textbf{Sujet :} \zh{我的一周} \emph{(Wǒ de yì zhōu)} — « Ma semaine ».

Rédige un court texte d'\textbf{au moins 60 caractères chinois} (environ 5 phrases) dans lequel tu racontes ta semaine type. Consignes obligatoires :
\begin{itemize}
\item utilise \textbf{au moins deux compléments de durée} (par exemple : 上课上了……分钟，睡觉睡……个小时) ;
\item utilise \textbf{au moins un adjectif redoublé} avec \zh{地} (par exemple : 慢慢地，高高兴兴地) ;
\item emploie un vocabulaire simple et des phrases courtes et correctes ;
\item soigne la ponctuation chinoise (，。).
\end{itemize}

\textbf{Aide :} tu peux t'inspirer du texte « Ma vie à Yaoundé » de l'exercice précédent, mais ton texte doit être personnel (tes cours, tes activités, ton sommeil, le football, le marché, etc.).
\end{exercice}

\newpage

\coursec{Corrigés des exercices}

\begin{corrige}[Exercice 1 — QCM : la bonne phrase]
\begin{enumerate}
\item \textbf{Réponse B} : \zh{她在中国住了三个月。} \emph{(Tā zài Zhōngguó zhù le sān ge yuè.)} « Elle a vécu trois mois en Chine. » Le complément de lieu \zh{在中国} se place \textbf{avant} le verbe ; le complément de durée \zh{三个月} se place \textbf{après} le verbe (avec \zh{了}). Les phrases A et C violent l'ordre des mots.
\item \textbf{Réponse A} : \zh{我们上课上了四十分钟。} \emph{(Wǒmen shàngkè shàng le sìshí fēnzhōng.)} « Nous avons cours pendant quarante minutes. » \zh{上课} est un verbe séparable (verbe + objet : 上 + 课) ; on répète donc le verbe : \cle{verbe + objet + verbe + 了 + durée}.
\item \textbf{Réponse B} : \zh{请你慢慢地说。} \emph{(Qǐng nǐ mànman de shuō.)} « Parle lentement, s'il te plaît. » Adjectif monosyllabique redoublé : \cle{AA + 地 + verbe}. La phrase A (sans redoublement) est incorrecte avec \zh{地} ; la phrase C a un ordre impossible.
\item \textbf{Réponse C} : \zh{他高高兴兴地回家了。} \emph{(Tā gāogāoxìngxìng de huí jiā le.)} « Il est rentré joyeusement à la maison. » Adjectif dissyllabique redoublé selon le schéma \cle{AABB + 地 + verbe}. En A, le redoublement est fautif (ABAB au lieu de AABB) ; en B, \zh{地} est mal placé.
\end{enumerate}
\end{corrige}

\begin{corrige}[Exercice 2 — Vrai ou faux ? Justifie ta réponse]
\begin{enumerate}
\item \textbf{F} — \zh{我没学了两年的中文。} est incorrecte. Règle : à la \textbf{négation} avec \zh{没(有)}, on \textbf{supprime 了}. Correction : \zh{我没学两年的中文。} ou mieux \zh{我没学过两年中文。} \emph{(Wǒ méi xuéguo liǎng nián Zhōngwén.)} « Je n'ai pas étudié le chinois pendant deux ans. »
\item \textbf{V} — \zh{她看电视看电视看了一个小时。} : correcte (lire : \zh{她看电视看了一个小时。}). \zh{看电视} est un verbe à objet ; le verbe est répété avant le complément de durée : \cle{V + O + V + 了 + durée}. « Elle a regardé la télévision pendant une heure. »
\item \textbf{F} — \zh{很大大大的眼睛。} est incorrecte. Règle : un adjectif \textbf{redoublé} ne peut pas être précédé de \zh{很} (le redoublement exprime déjà l'intensité ou la manière). Correction : \zh{大大的眼睛} \emph{(dàda de yǎnjing)} « de grands yeux » ou \zh{很大的眼睛} « de très grands yeux », mais jamais les deux ensemble.
\item \textbf{V} — \zh{他每天睡觉睡八个小时。} est correcte : \zh{睡觉} (dormir) est un verbe séparable (睡 + 觉) ; on répète le verbe \zh{睡} devant la durée. « Il dort huit heures chaque jour. »
\end{enumerate}
\end{corrige}

\begin{corrige}[Exercice 3 — Vocabulaire : durées et adjectifs]
\begin{center}
\begin{tabularx}{\linewidth}{|X|X|X|X|}
\hline
\rowcolor{popblueL} Caractères & Pinyin & Nature & Traduction \\
\hline
三个月 & sān ge yuè & complément de durée & trois mois \\
\hline
两年 & liǎng nián & complément de durée & deux ans \\
\hline
四十分钟 & sìshí fēnzhōng & complément de durée & quarante minutes \\
\hline
慢慢地 & mànman de & adverbe de manière (adj. redoublé + 地) & lentement \\
\hline
高高兴兴 & gāogāoxìngxìng & adjectif redoublé (AABB) & joyeux, joyeusement \\
\hline
一个小时 & yí ge xiǎoshí & complément de durée & une heure \\
\hline
\end{tabularx}
\end{center}
\textbf{Points de vigilance :} devant un classificateur, « deux » se dit \zh{两} et non \zh{二} (\zh{两年}, \zh{两个人}) ; \zh{一} devant un 4\up{e} ton se prononce \emph{yí} (\zh{一个小时}).
\end{corrige}

\begin{corrige}[Exercice 4 — Pinyin et caractères : à associer]
\begin{enumerate}
\item \emph{sān ge yuè} → \textbf{C} \zh{三个月}
\item \emph{mànman de} → \textbf{E} \zh{慢慢地}
\item \emph{bā ge xiǎoshí} → \textbf{F} \zh{八个小时}
\item \emph{gāogāoxìngxìng} → \textbf{A} \zh{高高兴兴}
\item \emph{liǎng nián} → \textbf{B} \zh{两年}
\item \emph{sìshí fēnzhōng} → \textbf{D} \zh{四十分钟}
\end{enumerate}
\end{corrige}

\begin{corrige}[Exercice 5 — Remise en ordre]
\begin{enumerate}
\item \zh{她在中国住了三个月。} \emph{(Tā zài Zhōngguó zhù le sān ge yuè.)} « Elle a vécu trois mois en Chine. » Schéma : Sujet + 在 + lieu + V + 了 + durée.
\item \zh{他看电视看了一个小时。} \emph{(Tā kàn diànshì kàn le yí ge xiǎoshí.)} « Il a regardé la télévision pendant une heure. » Verbe à objet : le verbe est répété.
\item \zh{我学中文学了两年。} ou \zh{我学了两年中文。} \emph{(Wǒ xué le liǎng nián Zhōngwén.)} « J'ai étudié le chinois pendant deux ans. » Quand l'objet est court et sans complément, la durée peut se placer entre le verbe et l'objet.
\item \zh{请你慢慢地说。} \emph{(Qǐng nǐ mànman de shuō.)} « Parle lentement, s'il te plaît. » Schéma : 请 + Sujet + AA + 地 + V.
\end{enumerate}
\end{corrige}

\begin{corrige}[Exercice 6 — Complète avec la durée]
\begin{enumerate}
\item \zh{我们上课上了四十分钟。} \emph{(Wǒmen shàngkè shàng le sìshí fēnzhōng.)} « Nous avons cours pendant quarante minutes. » (\zh{上课} : verbe séparable → répétition du verbe.)
\item \zh{他每天睡觉睡八个小时。} \emph{(Tā měi tiān shuìjiào shuì bā ge xiǎoshí.)} « Il dort huit heures par jour. » (Avec \zh{每天}, phrase habituelle : pas de \zh{了}.)
\item \zh{我在雅温得住了两年。} \emph{(Wǒ zài Yǎwēndé zhù le liǎng nián.)} « J'ai vécu deux ans à Yaoundé. »
\item \zh{她学习汉语学了三个月。} ou \zh{她学了三个月汉语。} \emph{(Tā xuéxí Hànyǔ xué le sān ge yuè.)} « Elle a étudié le chinois pendant trois mois. »
\end{enumerate}
\textbf{Point de vigilance :} ne jamais placer la durée avant le verbe (✱\zh{四十分钟上课}) ni après l'objet sans répéter le verbe (✱\zh{上课四十分钟}).
\end{corrige}

\begin{corrige}[Exercice 7 — Transforme les adjectifs]
\begin{enumerate}
\item \zh{请你慢慢地说。} \emph{(Qǐng nǐ mànman de shuō.)} « Parle lentement, s'il te plaît. » (A → AA : 慢 → 慢慢)
\item \zh{他高高兴兴地回家了。} \emph{(Tā gāogāoxìngxìng de huí jiā le.)} « Il est rentré joyeusement à la maison. » (AB → AABB : 高兴 → 高高兴兴)
\item \zh{老师认认真真地写字。} \emph{(Lǎoshī rènrenzhēnzhēn de xiězì.)} « Le professeur écrit avec sérieux. » (认真 → 认认真真)
\item \zh{我们快快乐乐地走进了教室。} \emph{(Wǒmen kuàikuàilèlè de zǒujìn le jiàoshì.)} « Nous sommes entrés joyeusement dans la salle. » (快乐 → 快快乐乐)
\end{enumerate}
\textbf{Règle rappelée :} adjectif d'une syllabe → \cle{AA + 地} ; adjectif de deux syllabes → \cle{AABB + 地} (jamais ABAB pour les adjectifs).
\end{corrige}

\begin{corrige}[Exercice 8 — Corrige les erreurs]
\begin{enumerate}
\item ✱\zh{我没学了两年的中文。} → \zh{我没学两年中文。} \emph{(Wǒ méi xué liǎng nián Zhōngwén.)} « Je n'ai pas étudié le chinois pendant deux ans. » \textbf{Règle :} avec la négation \zh{没(有)}, on supprime \zh{了} : la négation porte sur l'accomplissement, \zh{了} devient impossible.
\item ✱\zh{很大大大的眼睛。} → \zh{大大的眼睛} \emph{(dàda de yǎnjing)} « de grands yeux ». \textbf{Règle :} un adjectif redoublé est incompatible avec \zh{很} : le redoublement exprime déjà un degré.
\item ✱\zh{她在中国三个月住了。} → \zh{她在中国住了三个月。} \emph{(Tā zài Zhōngguó zhù le sān ge yuè.)} « Elle a vécu trois mois en Chine. » \textbf{Règle :} le complément de durée se place \textbf{après} le verbe (avec \zh{了}), jamais entre le complément de lieu et le verbe.
\item ✱\zh{请你慢地说。} → \zh{请你慢慢地说。} \emph{(Qǐng nǐ mànman de shuō.)} « Parle lentement, s'il te plaît. » \textbf{Règle :} devant \zh{地}, l'adjectif monosyllabique doit être redoublé (AA + 地 + V) ; \zh{慢地} seul est impossible.
\end{enumerate}
\end{corrige}

\begin{corrige}[Exercice 9 — Texte et questions de compréhension (type Bac)]
\textbf{Barème indicatif (10 pts) :} compréhension 4 × 1 pt ; questions de langue 3 × 2 pts.

\textbf{Questions de compréhension}
\begin{enumerate}
\item Marie apprend le chinois \textbf{depuis deux ans} (\zh{学习汉语学了两年了}).
\item Un cours dure \textbf{quarante minutes} (\zh{上课上四十分钟}).
\item Le professeur parle \textbf{lentement} (\zh{慢慢地说}) et les élèves apprennent \textbf{joyeusement} (\zh{高高兴兴地学习}).
\item Marie dort \textbf{huit heures} chaque nuit (\zh{睡觉睡八个小时}).
\end{enumerate}

\textbf{Questions de langue}
\begin{enumerate}
\item Deux compléments de durée : \zh{学了两年} (V + 了 + durée : 学 + 了 + 两年) et \zh{看一个小时} / \zh{睡八个小时} (V + durée, phrase habituelle). Schéma : \cle{verbe (+ répété) + 了 + durée}.
\item Adjectif redoublé : \zh{慢慢地} (AA + 地) ou \zh{高高兴兴地} (AABB + 地). Fonction : \textbf{complément circonstanciel de manière}, il décrit comment l'action est faite (« lentement », « joyeusement »).
\item Transformation : \zh{我们上课上了四十分钟。} \emph{(Wǒmen shàngkè shàng le sìshí fēnzhōng.)} « Nous avons cours pendant quarante minutes. »
\end{enumerate}
\textbf{Points de vigilance :} dans une réponse rédigée, reprendre la structure de la question ; ne pas confondre \zh{分钟} (minutes) et \zh{小时} (heures).
\end{corrige}

\begin{corrige}[Exercice 10 — Thème et version (type Bac)]
\textbf{Barème indicatif (10 pts) :} 2 pts par phrase (1 pt structure, 1 pt caractères/vocabulaire).

\textbf{Thème}
\begin{enumerate}
\item \zh{我学中文学了两年了。} \emph{(Wǒ xué Zhōngwén xué le liǎng nián le.)} — « depuis deux ans » = durée qui continue : \zh{了……了} en fin de phrase. Variante acceptée : \zh{我学了两年中文了。}
\item \zh{他没睡八个小时。} \emph{(Tā méi shuì bā ge xiǎoshí.)} — négation : \zh{没} sans \zh{了}.
\item \zh{我们看足球比赛看了两个小时。} \emph{(Wǒmen kàn zúqiú bǐsài kàn le liǎng ge xiǎoshí.)} — verbe à objet \zh{看比赛} → répétition du verbe.
\end{enumerate}

\textbf{Version}
\begin{enumerate}
\item \zh{他高高兴兴地回家了。} : « Il est rentré joyeusement à la maison. »
\item \zh{她看电视看了一个小时。} : « Elle a regardé la télévision pendant une heure. »
\item \zh{请你慢慢地写字。} : « Écris lentement, s'il te plaît. »
\end{enumerate}
\textbf{Points de vigilance :} « deux heures » = \zh{两个小时} (classificateur \zh{个} obligatoire) ; « deux ans » = \zh{两年} (sans classificateur, \zh{年} est lui-même l'unité).
\end{corrige}

\begin{corrige}[Exercice 11 — Expression écrite guidée (type Bac)]
\textbf{Barème indicatif (10 pts) :} respect des consignes (2 compléments de durée + 1 adjectif redoublé) 4 pts ; correction grammaticale 3 pts ; vocabulaire et ponctuation 2 pts ; longueur (≥ 60 caractères) 1 pt.

\textbf{Proposition de rédaction modèle :}

\begin{textebox}[我的一周 — modèle de rédaction]
我是杜阿拉的中学生。我的一周很忙。\\
\emph{Wǒ shì Dù'ālā de zhōngxuéshēng. Wǒ de yì zhōu hěn máng.}\\
Je suis lycéen à Douala. Ma semaine est bien remplie.\\[4pt]
每天我们上课上六个小时。我学习汉语学了两年了。\\
\emph{Měi tiān wǒmen shàngkè shàng liù ge xiǎoshí. Wǒ xuéxí Hànyǔ xué le liǎng nián le.}\\
Chaque jour, nous avons six heures de cours. J'apprends le chinois depuis deux ans.\\[4pt]
老师慢慢地说，我高高兴兴地学习。\\
\emph{Lǎoshī mànman de shuō, wǒ gāogāoxìngxìng de xuéxí.}\\
Le professeur parle lentement, et j'apprends joyeusement.\\[4pt]
晚上，我看电视看一个小时，然后睡觉睡八个小时。\\
\emph{Wǎnshang, wǒ kàn diànshì kàn yí ge xiǎoshí, ránhòu shuìjiào shuì bā ge xiǎoshí.}\\
Le soir, je regarde la télévision pendant une heure, puis je dors huit heures.
\end{textebox}

\textbf{Points de vigilance :}
\begin{itemize}
\item Vérifier la présence des \textbf{deux compléments de durée} : ici \zh{上课上六个小时}, \zh{学了两年}, \zh{看一个小时}, \zh{睡八个小时}.
\item Vérifier l'\textbf{adjectif redoublé + 地} : \zh{慢慢地说}, \zh{高高兴兴地学习}.
\item Ne pas écrire ✱\zh{很慢慢地} ni ✱\zh{我睡了八个小时没} ; à la négation, utiliser \zh{没} sans \zh{了}.
\item Ponctuation chinoise pleine chasse : \zh{，} et \zh{。}, sans espace entre les caractères.
\end{itemize}
\end{corrige}

\newpage

\coursec{Fiche bilan}

\courssub{Fiche bilan — Leçon 16 : Complément de durée \& Redoublement des adjectifs}

\begin{popfigure}
\begin{tikzpicture}[
    mindmap,
    concept color=popblue,
    level 1/.append style={level distance=3.8cm, sibling angle=360/7},
    level 2/.append style={level distance=2.2cm},
    every node/.style={font=\small, align=center},
    c1/.style={concept color=poporange},
    c2/.style={concept color=popgreen},
    c3/.style={concept color=poppurple},
    c4/.style={concept color=poppink},
    c5/.style={concept color=popgold},
    c6/.style={concept color=popblue!70!popdark},
    c7/.style={concept color=popmuted}
]
\node[concept, text=white, minimum size=2.2cm, align=center] (center){Leçon 16\\Durée \& Redoublement}
    child[c1] { node[concept, minimum size=1.8cm, align=center]{Complément\\de durée\\ \zh{补语}} 
        child { node[rectangle, rounded corners, fill=poporangeL, draw=poporange, inner sep=3pt] {S + V + Durée + \zh{了/过}} }
        child { node[rectangle, rounded corners, fill=poporangeL, draw=poporange, inner sep=3pt] {Négation : \zh{没(有) + V + Durée}} }
    }
    child[c2] { node[concept, minimum size=1.8cm, align=center]{COD \& Durée\\ \zh{动宾结构}} 
        child { node[rectangle, rounded corners, fill=popgreenL, draw=popgreen, inner sep=3pt] {COD court/pronom : V + Durée + COD} }
        child { node[rectangle, rounded corners, fill=popgreenL, draw=popgreen, inner sep=3pt] {COD long/spécifique : V + COD + V + Durée} }
    }
    child[c3] { node[concept, minimum size=1.8cm, align=center]{Redoublement\\Adj. monosyll.\\ \zh{AA $\to$ AABB / 一A一A}} 
        child { node[rectangle, rounded corners, fill=poppurpleL, draw=poppurple, inner sep=3pt] {\zh{高 $\to$ 高高 / 高高兴兴}} }
        child { node[rectangle, rounded corners, fill=poppurpleL, draw=poppurple, inner sep=3pt] {Fonction adverbiale + \zh{地}} }
    }
    child[c4] { node[concept, minimum size=1.8cm, align=center]{Redoublement\\Adj. disyll.\\ \zh{AB $\to$ AABB / ABAB}} 
        child { node[rectangle, rounded corners, fill=poppinkL, draw=poppink, inner sep=3pt] {\zh{高兴 $\to$ 高高兴兴}} }
        child { node[rectangle, rounded corners, fill=poppinkL, draw=poppink, inner sep=3pt] {\zh{认真 $\to$ 认认真真}} }
    }
    child[c5] { node[concept, minimum size=1.8cm, align=center]{Nuances\\expressives} 
        child { node[rectangle, rounded corners, fill=popgoldL, draw=popgold, inner sep=3pt] {Légèreté, décontraction} }
        child { node[rectangle, rounded corners, fill=popgoldL, draw=popgold, inner sep=3pt] {Description d'état} }
    }
    child[c6] { node[concept, minimum size=1.8cm, align=center]{Pièges\\francophones} 
        child { node[rectangle, rounded corners, fill=popblue!30, draw=popblue, inner sep=3pt] {Ordre : *V COD Durée} }
        child { node[rectangle, rounded corners, fill=popblue!30, draw=popblue, inner sep=3pt] {\zh{了} position / oubli \zh{没}} }
        child { node[rectangle, rounded corners, fill=popblue!30, draw=popblue, inner sep=3pt] {Confusion \zh{次/遍} vs Durée} }
    }
    child[c7] { node[concept, minimum size=1.8cm, align=center]{Mots-clés\\temps} 
        child { node[rectangle, rounded corners, fill=popcream, draw=popmuted, inner sep=3pt] {\zh{小时、分钟、天、年}} }
        child { node[rectangle, rounded corners, fill=popcream, draw=popmuted, inner sep=3pt] {\zh{次、遍、下、趟}} }
    };
\end{tikzpicture}
\legende{Carte mentale récapitulative : structures, cas de figure, redoublement et pièges à éviter.}
\end{popfigure}

\begin{bilan}

\end{bilan}

\courssub{Lexique clé — 重点词汇}
\begin{center}
\begin{tabularx}{\linewidth}{|X|X|X|X|}
\hline
\rowcolor{popblueL} \textbf{Caractères} & \textbf{Pinyin} & \textbf{Nature} & \textbf{Traduction} \\ \hline
小时 & xiǎoshí & N & heure (durée) \\ \hline
分钟 & fēnzhōng & N & minute \\ \hline
天 & tiān & N & jour / journée \\ \hline
年 & nián & N & an / année \\ \hline
次 & cì & M & fois (fréquence, action brève) \\ \hline
遍 & biàn & M & fois (fréquence, action complète) \\ \hline
下 & xià & M & fois (action brève, tentative) \\ \hline
趟 & tàng & M & fois (aller-retour, trajet) \\ \hline
红红 & hónghóng & Adj (red.) & rougeâtre, tout rouge \\ \hline
高高兴兴 & gāogāoxìngxìng & Adj (red.) & joyeusement, content \\ \hline
认认真真 & rènrènzhēnzhēn & Adv (red.) & sérieusement, soigneusement \\ \hline
一清二楚 & yīqīngèrchǔ & Id. (red.) & parfaitement clair \\ \hline
\end{tabularx}
\end{center}

\courssub{Règles de grammaire à connaître par cœur}
\begin{center}
\renewcommand{\arraystretch}{1.4}
\begin{tabularx}{\linewidth}{|>{\hsize=0.6\hsize}X|>{\hsize=1.2\hsize}X|>{\hsize=1.2\hsize}X|}
\hline
\rowcolor{popblueL} \textbf{Structure / Règle} & \textbf{Schéma / Exemple (Car. + Pinyin)} & \textbf{Traduction / Note} \\ \hline
\textbf{1. Compl. durée simple} & S + V + Durée (+ \zh{了/过}) & Action terminée (\zh{了}) ou vécue (\zh{过}). \\ \cline{2-3}
& \zh{我学了两年汉语。} & J'étudie le chinois depuis 2 ans (encore maintenant). \\ \hline
\textbf{2. Négation durée} & S + \zh{没(有)} + V + Durée & \textbf{Pas de \zh{了}}. \zh{过} possible pour « n'avoir jamais fait pendant... ». \\ \cline{2-3}
& \zh{我没睡多久。} & Je n'ai pas dormi longtemps. \\ \hline
\textbf{3. COD court / pronom} & S + V + Durée + COD & COD = \zh{它、这本书、饭} (court). \\ \cline{2-3}
& \zh{我看了两遍这本书。} & J'ai lu ce livre deux fois. \\ \hline
\textbf{4. COD long / spécifique} & S + V + COD + V + Durée & \textbf{Répétition du verbe} obligatoire. \\ \cline{2-3}
& \zh{我看这本书看了两年。} & Je lis ce livre depuis 2 ans. \\ \hline
\textbf{5. Redoublement Mono.} & AA $\to$ \zh{AABB} ou \zh{一A一A} & \zh{高 $\to$ 高高 / 高高兴兴} ; \zh{看 $\to$ 看看} (verbe). \\ \cline{2-3}
& \zh{天蓝蓝的。} & Le ciel est bleu (description vive). \\ \hline
\textbf{6. Redoublement Disyll.} & AB $\to$ \zh{AABB} (qualitatif) / \zh{ABAB} (manière) & \zh{干净 $\to$ 干干净净} ; \zh{认真 $\to$ 认认真真}. \\ \cline{2-3}
& \zh{他认认真真地学习。} & Il étudie sérieusement. \\ \hline
\textbf{7. Adj redoublé + \zh{地}} & AdjRedup + \zh{地} + V & Fonction adverbiale (manière). \\ \cline{2-3}
& \zh{高高兴兴地去。} & Y aller joyeusement. \\ \hline
\end{tabularx}
\end{center}

\courssub{Méthodes express — Réflexes de construction}
\begin{itemize}
    \item \textbf{Durée :} Repérer le \textbf{verbe d'action} $\to$ placer la durée \textbf{immédiatement après}.
    \item \textbf{COD présent ?} Court/pronom $\to$ après durée. Long/spécifique $\to$ \textbf{répéter le verbe} (V + COD + V + Durée).
    \item \textbf{Négation :} Remplacer \zh{了/过} par \zh{没(有)} devant le verbe, garder la durée à la fin.
    \item \textbf{Redoublement Adj :} Mono $\to$ \zh{AA} (souvent \zh{AABB} lexicalisé) ; Di $\to$ \zh{AABB} (état) ou \zh{ABAB} (manière). Ajouter \zh{地} pour modifier un verbe.
    \item \textbf{Traduction « Cela fait... que » :} \zh{S + V + Durée + 了} (ex: \zh{我来了两年了} = Cela fait 2 ans que je suis là).
\end{itemize}



\begin{aretenir}[Checklist avant le Bac — Leçon 16]
\begin{itemize}
    \item $\square$ Je sais placer le complément de durée \textbf{juste après le verbe} (S + V + Durée).
    \item $\square$ Je distingue \zh{了} (action achevée/état nouveau) et \zh{过} (expérience passée) avec une durée.
    \item $\square$ Je construis la négation avec \zh{没(有)} + V + Durée \textbf{sans \zh{了}}.
    \item $\square$ Je gère le COD : \textbf{court/pronom} après la durée ; \textbf{long/spécifique} $\to$ répétition du verbe (V + COD + V + Durée).
    \item $\square$ Je ne confonds pas \textbf{durée} (\zh{两小时}) et \textbf{fréquence} (\zh{两次/两遍}) : place différente (fréquence avant V ou après V selon structure).
    \item $\square$ Je redouble correctement les adjectifs monosyllabiques (\zh{AA} $\to$ \zh{AABB} / \zh{一A一A}) et disyllabiques (\zh{AB} $\to$ \zh{AABB} / \zh{ABAB}).
    \item $\square$ J'utilise \zh{地} après un adjectif redoublé pour en faire un adverbe de manière (\zh{高高兴兴地}).
    \item $\square$ Je perçois la nuance \textbf{affective / descriptive / décontractée} du redoublement (pas juste « plus intense »).
    \item $\square$ Je sais traduire « Cela fait [durée] que... » par \zh{S + V + Durée + 了} (action continue) ou \zh{没 + V + Durée + 了} (action interrompue).
    \item $\square$ J'évite l'erreur classique \emph{*Sujet + Verbe + COD + Durée} (sans répétition du verbe) pour un COD long.
    \item $\square$ Je reconnais les classificateurs de fréquence (\zh{次、遍、下、趟}) et leur place (souvent avant V ou après V + COD).
\end{itemize}
\end{aretenir}

\findecours

\end{document}
